\documentclass[11pt]{article}

% ---- theme variables (passed via pandoc -V) ----
\newcommand{\ThemeName}{Design Patterns}
\newcommand{\ThemeKicker}{Design Track}
\newcommand{\ThemeNumber}{}

\usepackage{interviewstyle}
\setthemecolor{0284c7}{38bdf8}

% pandoc-required plumbing
\providecommand{\tightlist}{\setlength{\itemsep}{1pt}\setlength{\parskip}{0pt}}
\setlength{\emergencystretch}{3em}
\providecommand{\pandocbounded}[1]{#1}

% syntax-highlighting macros emitted by pandoc (skylighting)
\usepackage{color}
\usepackage{fancyvrb}
\newcommand{\VerbBar}{|}
\newcommand{\VERB}{\Verb[commandchars=\\\{\}]}
\DefineVerbatimEnvironment{Highlighting}{Verbatim}{commandchars=\\\{\}}
% Add ',fontsize=\small' for more characters per line
\usepackage{framed}
\definecolor{shadecolor}{RGB}{35,38,41}
\newenvironment{Shaded}{\begin{snugshade}}{\end{snugshade}}
\newcommand{\AlertTok}[1]{\textcolor[rgb]{0.58,0.85,0.30}{\textbf{\colorbox[rgb]{0.30,0.12,0.14}{#1}}}}
\newcommand{\AnnotationTok}[1]{\textcolor[rgb]{0.25,0.50,0.35}{#1}}
\newcommand{\AttributeTok}[1]{\textcolor[rgb]{0.16,0.50,0.73}{#1}}
\newcommand{\BaseNTok}[1]{\textcolor[rgb]{0.96,0.45,0.00}{#1}}
\newcommand{\BuiltInTok}[1]{\textcolor[rgb]{0.50,0.55,0.55}{#1}}
\newcommand{\CharTok}[1]{\textcolor[rgb]{0.24,0.68,0.91}{#1}}
\newcommand{\CommentTok}[1]{\textcolor[rgb]{0.48,0.49,0.49}{#1}}
\newcommand{\CommentVarTok}[1]{\textcolor[rgb]{0.50,0.55,0.55}{#1}}
\newcommand{\ConstantTok}[1]{\textcolor[rgb]{0.15,0.68,0.68}{\textbf{#1}}}
\newcommand{\ControlFlowTok}[1]{\textcolor[rgb]{0.99,0.74,0.29}{\textbf{#1}}}
\newcommand{\DataTypeTok}[1]{\textcolor[rgb]{0.16,0.50,0.73}{#1}}
\newcommand{\DecValTok}[1]{\textcolor[rgb]{0.96,0.45,0.00}{#1}}
\newcommand{\DocumentationTok}[1]{\textcolor[rgb]{0.64,0.20,0.25}{#1}}
\newcommand{\ErrorTok}[1]{\textcolor[rgb]{0.85,0.27,0.33}{\underline{#1}}}
\newcommand{\ExtensionTok}[1]{\textcolor[rgb]{0.00,0.60,1.00}{\textbf{#1}}}
\newcommand{\FloatTok}[1]{\textcolor[rgb]{0.96,0.45,0.00}{#1}}
\newcommand{\FunctionTok}[1]{\textcolor[rgb]{0.56,0.27,0.68}{#1}}
\newcommand{\ImportTok}[1]{\textcolor[rgb]{0.15,0.68,0.38}{#1}}
\newcommand{\InformationTok}[1]{\textcolor[rgb]{0.77,0.36,0.00}{#1}}
\newcommand{\KeywordTok}[1]{\textcolor[rgb]{0.81,0.81,0.76}{\textbf{#1}}}
\newcommand{\NormalTok}[1]{\textcolor[rgb]{0.81,0.81,0.76}{#1}}
\newcommand{\OperatorTok}[1]{\textcolor[rgb]{0.81,0.81,0.76}{#1}}
\newcommand{\OtherTok}[1]{\textcolor[rgb]{0.15,0.68,0.38}{#1}}
\newcommand{\PreprocessorTok}[1]{\textcolor[rgb]{0.15,0.68,0.38}{#1}}
\newcommand{\RegionMarkerTok}[1]{\textcolor[rgb]{0.16,0.50,0.73}{\colorbox[rgb]{0.08,0.19,0.26}{#1}}}
\newcommand{\SpecialCharTok}[1]{\textcolor[rgb]{0.24,0.68,0.91}{#1}}
\newcommand{\SpecialStringTok}[1]{\textcolor[rgb]{0.85,0.27,0.33}{#1}}
\newcommand{\StringTok}[1]{\textcolor[rgb]{0.96,0.31,0.31}{#1}}
\newcommand{\VariableTok}[1]{\textcolor[rgb]{0.15,0.68,0.68}{#1}}
\newcommand{\VerbatimStringTok}[1]{\textcolor[rgb]{0.85,0.27,0.33}{#1}}
\newcommand{\WarningTok}[1]{\textcolor[rgb]{0.85,0.27,0.33}{#1}}

\begin{document}

\makecoverpage

\thispagestyle{empty}
{\hypersetup{hidelinks}\tableofcontents}
\clearpage

\begin{quote}
\textbf{How to use this file:} Read top-to-bottom for deep mastery. For last-minute revision, jump to §Common Interview Questions + §Revision Cheat Sheet. Patterns are the working vocabulary of \textbf{Low-Level Design} (\hyperlink{section.06}{\textcolor{acc}{§06}}) --- internalize each pattern\textquotesingle s \emph{intent} and the \emph{force} it relieves, not just the UML box-and-arrow. Every pattern below has: Intent, the concrete Problem it solves, a Structure diagram, real Java/Python code, when (and when NOT) to use it, a real framework that ships it, and trade-offs.
\end{quote}

\section{Overview --- What It Is}\label{overview--what-it-is}

A \textbf{design pattern} is a named, reusable \emph{template} for solving a recurring object-oriented design problem. It is \textbf{not} a library, a framework, or copy-pasteable code. It is a description of \emph{which classes to introduce, how they collaborate, and which responsibilities to assign to which object} --- captured at a level of abstraction you can re-apply across wildly different domains.

The canonical 23 patterns come from the 1994 book \emph{Design Patterns: Elements of Reusable Object-Oriented Software} by Gamma, Helm, Johnson, and Vlissides --- the \textbf{Gang of Four (GoF)}. They sorted patterns into three families by \emph{what aspect of the system they manage}:

\setcodelang{}

\begin{verbatim}
CREATIONAL              STRUCTURAL              BEHAVIORAL
"how objects are        "how objects are        "how objects
 created & who          composed into            communicate &
 owns creation"          larger structures"       distribute work"

Singleton               Adapter                 Strategy
Factory Method          Decorator               Observer
Abstract Factory        Facade                  Command
Builder                 Proxy                   State
Prototype               Composite               Template Method
                        Bridge                  Iterator
                        Flyweight               Chain of Responsibility
                                                Mediator
                                                Memento
                                                Visitor
                                                Interpreter
\end{verbatim}

Each pattern is described by four things:

\begin{longtable}[]{@{}
  >{\raggedright\arraybackslash}p{(\columnwidth - 4\tabcolsep) * \real{0.1524}}
  >{\raggedright\arraybackslash}p{(\columnwidth - 4\tabcolsep) * \real{0.3852}}
  >{\raggedright\arraybackslash}p{(\columnwidth - 4\tabcolsep) * \real{0.4624}}@{}}
\toprule\noalign{}
\rowcolor{acc!16}
\begin{minipage}[b]{\linewidth}\raggedright
Element
\end{minipage} & \begin{minipage}[b]{\linewidth}\raggedright
Meaning
\end{minipage} & \begin{minipage}[b]{\linewidth}\raggedright
Example (Strategy)
\end{minipage} \\
\midrule\noalign{}
\endhead
\bottomrule\noalign{}
\endlastfoot
\textbf{Intent} & The one-line purpose & "Make a family of algorithms interchangeable" \\
\rowcolor{zebra}
\textbf{Problem / Forces} & The tension it resolves & "Algorithm must vary at runtime without editing the caller" \\
\textbf{Structure} & The participating classes + relationships & \texttt{Context} holds a \texttt{Strategy} interface, swaps concrete impls \\
\rowcolor{zebra}
\textbf{Consequences} & The trade-offs you accept & More classes, runtime flexibility, OCP compliance \\
\end{longtable}

Beyond GoF, \textbf{modern/architectural patterns} --- Dependency Injection, MVC/MVP/MVVM, Repository, Null Object, Object Pool --- appear constantly in real codebases and interviews, so they get full treatment here too.

\textbf{Interview takeaway:} A pattern is \emph{a name for a shape that good OO code naturally takes}. Your job in an interview is to \textbf{recognize the shape} the requirement is pushing you toward, name it, and justify the trade-off --- not to recite UML.

\section{Why It Exists}\label{why-it-exists}

Before patterns, every experienced designer kept re-deriving the \emph{same} solutions to the \emph{same} structural problems --- and every team described those solutions with different words. Patterns exist to fix three things:

\begin{enumerate}
\def\labelenumi{\arabic{enumi}.}
\item
  \textbf{Re-invention.} The problem "I need to add behavior to objects at runtime without a combinatorial explosion of subclasses" has been solved thousands of times. The Decorator pattern is the distilled, vetted answer. You shouldn\textquotesingle t re-discover it under interview pressure.
\item
  \textbf{Communication.} Saying \emph{"wrap it in a Strategy"} transmits an entire micro-architecture --- an interface, a context that delegates, concrete implementations, runtime swapping --- in three words. Patterns are a \textbf{compression codec for design intent}. On a team, this is the difference between a 5-minute alignment and a 50-minute argument.
\item
  \textbf{Principle-in-practice.} Every pattern is SOLID + composition-over-inheritance applied to a \emph{specific shape} of problem. Strategy \emph{is} the Open/Closed Principle made concrete. Observer \emph{is} loose coupling between producer and consumer. Patterns are how abstract principles become muscle memory.
\end{enumerate}

Symptoms a codebase is \emph{missing} the right patterns:

\begin{itemize}
\tightlist
\item
  A \texttt{switch\ (type)} statement that grows every time a new type is added (\textbf{needs Strategy/Factory/State}).
\item
  Adding a feature requires editing 12 files (\textbf{shotgun surgery --- needs better encapsulation of what varies}).
\item
  A class with 4,000 lines doing everything (\textbf{God object --- needs SRP + delegation}).
\item
  Subclass explosion: \texttt{BufferedEncryptedCompressedStream} (\textbf{needs Decorator}).
\end{itemize}

\begin{quote}
Patterns are \emph{discovered, not invented}. They are the recurring crystal shapes good object-oriented code condenses into under the pressure of change. The skill is recognizing the shape early.
\end{quote}

\section{Why FAANG Cares}\label{why-faang-cares}

Design patterns sit at the intersection of three signals interviewers explicitly grade: \textbf{code quality}, \textbf{extensibility}, and \textbf{judgment}.

\begin{itemize}
\item
  \textbf{Amazon} --- Bar Raisers tie clean OOD to "Insist on Highest Standards" and "Dive Deep." In machine-coding / LLD rounds (parking lot, ride-share, in-memory cache) they watch whether you reach for Strategy/Factory/Observer \emph{appropriately} and whether you can extend your own design mid-interview without rewriting it (OCP in action).
\item
  \textbf{Google} --- Grades \emph{code quality} as a distinct hiring signal. The bar is \textbf{right-sized abstraction}: do you know when a pattern earns its indirection vs. when a plain function is better? Over-patterning is penalized as hard as under-designing.
\item
  \textbf{Meta} --- "Move fast" without leaving landmines. LLD rounds test whether your design absorbs new requirements gracefully. Patterns are the mechanism; the judgment to \emph{not} use them is the senior signal.
\item
  \textbf{Microsoft} --- Heavy on SOLID and extensibility. Expect "now add requirement X" to your design --- patterns are what make that a 2-line answer instead of a refactor.
\item
  \textbf{Apple} --- Cares about ownership semantics, interface design, and correct composition vs. inheritance choices.
\end{itemize}

\textbf{The specific signals being tested:}

\begin{enumerate}
\def\labelenumi{\arabic{enumi}.}
\tightlist
\item
  Can you \textbf{map a requirement to a pattern}? ("Pricing must be pluggable" \(\rightarrow\) Strategy.)
\item
  Can you \textbf{articulate the trade-off}? ("Decorator avoids subclass explosion but produces many small wrapper objects.")
\item
  Do you know \textbf{when NOT to use a pattern}? (Two states that never change \(\rightarrow\) if/else beats State.)
\item
  Can you \textbf{read framework code}? (Spring beans = Factory + Singleton + Proxy; React = Observer + Composite.)
\item
  Do you recognize \textbf{anti-patterns} and know the refactor that fixes them?
\end{enumerate}

\begin{quote}
\textbf{Interview signal:} Naming a pattern \emph{and} justifying its trade-off marks a senior engineer. Naming a pattern you don\textquotesingle t need marks a junior one.
\end{quote}

\section{Core Concepts}\label{core-concepts}

Patterns are not random tricks --- they are all derived from a small set of \textbf{design principles}. Master these five and most patterns become obvious consequences.

\subsection{1. Program to an Interface, Not an Implementation}\label{1-program-to-an-interface-not-an-implementation}

The foundational GoF mantra. Depend on \textbf{abstractions} (interfaces / abstract base classes), never on concrete classes. This is the \textbf{Dependency Inversion Principle (DIP)} in one sentence.

\setcodelang{Java}

\begin{Shaded}
\begin{Highlighting}[]
\CommentTok{// BAD — caller welded to a concrete class; swapping it edits this code}
\BuiltInTok{ArrayList}\OperatorTok{\textless{}}\NormalTok{Order}\OperatorTok{\textgreater{}}\NormalTok{ orders }\OperatorTok{=} \KeywordTok{new} \BuiltInTok{ArrayList}\OperatorTok{\textless{}\textgreater{}();}

\CommentTok{// GOOD — caller depends only on the List contract; impl is swappable}
\BuiltInTok{List}\OperatorTok{\textless{}}\NormalTok{Order}\OperatorTok{\textgreater{}}\NormalTok{ orders }\OperatorTok{=} \KeywordTok{new} \BuiltInTok{ArrayList}\OperatorTok{\textless{}\textgreater{}();}   \CommentTok{// could be LinkedList, CopyOnWriteArrayList…}
\end{Highlighting}
\end{Shaded}

The payoff: you can substitute mocks in tests, swap a \texttt{MySQLRepository} for an \texttt{InMemoryRepository}, and the calling code never changes. \textbf{Every creational pattern exists to help you obey this rule} --- they let the caller request an abstraction and receive a concrete object without naming the concrete class.

\subsection{2. Favor Composition Over Inheritance}\label{2-favor-composition-over-inheritance}

Inheritance ("white-box reuse") exposes the parent\textquotesingle s internals to the child, creating the \textbf{fragile base class problem}: a change to the parent silently breaks children. Composition ("black-box reuse") delegates to a held object through its public interface --- looser, runtime-swappable.

\setcodelang{Python}

\begin{Shaded}
\begin{Highlighting}[]
\CommentTok{\# Inheritance — Logger is permanently a FileHandler; can\textquotesingle{}t switch to network logging}
\KeywordTok{class}\NormalTok{ Logger(FileHandler): ...}

\CommentTok{\# Composition — swap the handler at runtime; Logger is decoupled from storage}
\KeywordTok{class}\NormalTok{ Logger:}
    \KeywordTok{def} \FunctionTok{\_\_init\_\_}\NormalTok{(}\VariableTok{self}\NormalTok{, handler: LogHandler):  }\CommentTok{\# FileHandler | ConsoleHandler | KafkaHandler}
        \VariableTok{self}\NormalTok{.\_handler }\OperatorTok{=}\NormalTok{ handler}
    \KeywordTok{def}\NormalTok{ log(}\VariableTok{self}\NormalTok{, msg): }\VariableTok{self}\NormalTok{.\_handler.emit(msg)}
\end{Highlighting}
\end{Shaded}

Most structural and behavioral patterns (Strategy, Decorator, Bridge, State, Adapter, Proxy) are \textbf{composition} mechanisms. When you find yourself subclassing purely to reuse code, stop --- compose instead.

\subsection{3. Encapsulate What Varies}\label{3-encapsulate-what-varies}

Identify the part of the system that \textbf{changes} (the algorithm, the product type, the notification target) and isolate it behind a stable interface. The stable rest of the system then \emph{never changes} when the varying part does --- this \textbf{is} the Open/Closed Principle (open for extension, closed for modification).

\setcodelang{}

\begin{verbatim}
The whole catalog of patterns, in one sentence:
  "Find what varies, wrap it behind an interface, and depend on the interface."
  - Algorithm varies      → Strategy
  - Object created varies  → Factory
  - State-dependent behavior varies → State
  - Wrapped responsibilities vary    → Decorator
\end{verbatim}

\subsection{4. SOLID Recap (the principles patterns enforce)}\label{4-solid-recap-the-principles-patterns-enforce}

\begin{longtable}[]{@{}
  >{\raggedright\arraybackslash}p{(\columnwidth - 6\tabcolsep) * \real{0.0735}}
  >{\raggedright\arraybackslash}p{(\columnwidth - 6\tabcolsep) * \real{0.2236}}
  >{\raggedright\arraybackslash}p{(\columnwidth - 6\tabcolsep) * \real{0.3727}}
  >{\raggedright\arraybackslash}p{(\columnwidth - 6\tabcolsep) * \real{0.3301}}@{}}
\toprule\noalign{}
\rowcolor{acc!16}
\begin{minipage}[b]{\linewidth}\raggedright
Letter
\end{minipage} & \begin{minipage}[b]{\linewidth}\raggedright
Principle
\end{minipage} & \begin{minipage}[b]{\linewidth}\raggedright
One-liner
\end{minipage} & \begin{minipage}[b]{\linewidth}\raggedright
Pattern that embodies it
\end{minipage} \\
\midrule\noalign{}
\endhead
\bottomrule\noalign{}
\endlastfoot
\textbf{S} & Single Responsibility & One class, one reason to change & Command (one action per object) \\
\rowcolor{zebra}
\textbf{O} & Open/Closed & Extend by adding code, not editing & Strategy, Factory, Decorator \\
\textbf{L} & Liskov Substitution & Subtypes honor the base contract & Template Method, Null Object \\
\rowcolor{zebra}
\textbf{I} & Interface Segregation & Many small interfaces \textgreater{} one fat one & Adapter, Visitor \\
\textbf{D} & Dependency Inversion & Depend on abstractions & Abstract Factory, DI \\
\end{longtable}

\begin{quote}
\textbf{Mnemonic:} A pattern that \emph{violates} SOLID is being misused. If your "Singleton" has 20 unrelated methods, it\textquotesingle s also a God object (SRP violation) --- the pattern label doesn\textquotesingle t excuse it.
\end{quote}

\subsection{5. The Hollywood Principle --- "Don\textquotesingle t call us, we\textquotesingle ll call you"}\label{5-the-hollywood-principle--dont-call-us-well-call-you}

Inversion of Control (IoC): high-level code defines the \emph{skeleton} and calls down into low-level hooks; low-level components do not call up into the framework. This is the heart of \textbf{Template Method} (base class calls your overridden steps), \textbf{Observer} (subject calls your \texttt{update}), and \textbf{Dependency Injection} (the container constructs and hands you your dependencies).

\setcodelang{}

\begin{verbatim}
Traditional flow:        your code → library
Hollywood / IoC:         framework → your hooks   (framework owns the control loop)
\end{verbatim}

This single inversion is \emph{why} frameworks (Spring, JUnit, React) feel different from libraries: you fill in blanks, the framework runs the show.

\textbf{Interview takeaway:} If asked "what\textquotesingle s the principle behind this pattern?", almost every correct answer is one of: \emph{encapsulate what varies (OCP), program to interfaces (DIP), favor composition, or invert control (Hollywood)}.

\section{Creational Patterns}\label{creational-patterns}

Creational patterns abstract the \textbf{instantiation process} --- they let a system be independent of \emph{how its objects are created, composed, and represented}. The common thread: the caller asks for an abstraction and receives a concrete object \textbf{without writing \texttt{new\ ConcreteClass()}} itself.

\subsection{Singleton}\label{singleton}

\textbf{Intent:} Ensure a class has exactly one instance and provide a single global access point to it.

\textbf{Problem it solves:} Some resources must be unique and shared: a connection pool, an in-memory cache, a logger, an app-wide config registry. Creating two would be wrong (two pools double the DB connections) or wasteful. You need \emph{one} instance and a way for any code to reach it.

\textbf{Structure:}

\setcodelang{}

\begin{verbatim}
+----------------------+
|     Singleton        |
+----------------------+
| - instance: Singleton|  (static, private)
| - Singleton()        |  (private ctor — no outside `new`)
+----------------------+
| + getInstance(): S   |  (static accessor; lazily/eagerly creates)
+----------------------+
        ▲ returns the same object every call
\end{verbatim}

\textbf{The four canonical thread-safe Java variants} (this is the \#1 Singleton interview question):

\setcodelang{Java}

\begin{Shaded}
\begin{Highlighting}[]
\CommentTok{// 1) EAGER — instance built at class{-}load. Simple, thread{-}safe by JVM guarantee.}
\CommentTok{//    Downside: built even if never used.}
\KeywordTok{public} \KeywordTok{class}\NormalTok{ EagerConfig }\OperatorTok{\{}
    \KeywordTok{private} \DataTypeTok{static} \DataTypeTok{final}\NormalTok{ EagerConfig INSTANCE }\OperatorTok{=} \KeywordTok{new} \FunctionTok{EagerConfig}\OperatorTok{();}
    \KeywordTok{private} \FunctionTok{EagerConfig}\OperatorTok{()} \OperatorTok{\{\}}
    \KeywordTok{public} \DataTypeTok{static}\NormalTok{ EagerConfig }\FunctionTok{get}\OperatorTok{()} \OperatorTok{\{} \ControlFlowTok{return}\NormalTok{ INSTANCE}\OperatorTok{;} \OperatorTok{\}}
\OperatorTok{\}}

\CommentTok{// 2) DOUBLE{-}CHECKED LOCKING — lazy + thread{-}safe. \textasciigrave{}volatile\textasciigrave{} is MANDATORY}
\CommentTok{//    (without it, partially{-}constructed object can leak due to reordering).}
\KeywordTok{public} \KeywordTok{class}\NormalTok{ DclConfig }\OperatorTok{\{}
    \KeywordTok{private} \DataTypeTok{static} \KeywordTok{volatile}\NormalTok{ DclConfig instance}\OperatorTok{;}   \CommentTok{// volatile = no reordering}
    \KeywordTok{private} \FunctionTok{DclConfig}\OperatorTok{()} \OperatorTok{\{\}}
    \KeywordTok{public} \DataTypeTok{static}\NormalTok{ DclConfig }\FunctionTok{get}\OperatorTok{()} \OperatorTok{\{}
        \ControlFlowTok{if} \OperatorTok{(}\NormalTok{instance }\OperatorTok{==} \KeywordTok{null}\OperatorTok{)} \OperatorTok{\{}                   \CommentTok{// 1st check (no lock, fast path)}
            \KeywordTok{synchronized} \OperatorTok{(}\NormalTok{DclConfig}\OperatorTok{.}\FunctionTok{class}\OperatorTok{)} \OperatorTok{\{}
                \ControlFlowTok{if} \OperatorTok{(}\NormalTok{instance }\OperatorTok{==} \KeywordTok{null}\OperatorTok{)} \OperatorTok{\{}           \CommentTok{// 2nd check (locked)}
\NormalTok{                    instance }\OperatorTok{=} \KeywordTok{new} \FunctionTok{DclConfig}\OperatorTok{();}
                \OperatorTok{\}}
            \OperatorTok{\}}
        \OperatorTok{\}}
        \ControlFlowTok{return}\NormalTok{ instance}\OperatorTok{;}
    \OperatorTok{\}}
\OperatorTok{\}}

\CommentTok{// 3) BILL PUGH / INITIALIZATION{-}ON{-}DEMAND HOLDER — best of both.}
\CommentTok{//    Lazy (holder class loads only on first get()), thread{-}safe (JVM class{-}init lock),}
\CommentTok{//    no synchronization cost on the hot path. The recommended idiom.}
\KeywordTok{public} \KeywordTok{class}\NormalTok{ HolderConfig }\OperatorTok{\{}
    \KeywordTok{private} \FunctionTok{HolderConfig}\OperatorTok{()} \OperatorTok{\{\}}
    \KeywordTok{private} \DataTypeTok{static} \KeywordTok{class} \BuiltInTok{Holder} \OperatorTok{\{} \DataTypeTok{static} \DataTypeTok{final}\NormalTok{ HolderConfig I }\OperatorTok{=} \KeywordTok{new} \FunctionTok{HolderConfig}\OperatorTok{();} \OperatorTok{\}}
    \KeywordTok{public} \DataTypeTok{static}\NormalTok{ HolderConfig }\FunctionTok{get}\OperatorTok{()} \OperatorTok{\{} \ControlFlowTok{return} \BuiltInTok{Holder}\OperatorTok{.}\FunctionTok{I}\OperatorTok{;} \OperatorTok{\}}
\OperatorTok{\}}

\CommentTok{// 4) ENUM — Joshua Bloch\textquotesingle{}s preferred form. Thread{-}safe, serialization{-}safe,}
\CommentTok{//    reflection{-}proof (you cannot reflectively instantiate a second enum constant).}
\KeywordTok{public} \KeywordTok{enum}\NormalTok{ EnumConfig }\OperatorTok{\{}
\NormalTok{    INSTANCE}\OperatorTok{;}
    \KeywordTok{public} \DataTypeTok{void} \FunctionTok{load}\OperatorTok{()} \OperatorTok{\{} \CommentTok{/* ... */} \OperatorTok{\}}
\OperatorTok{\}}
\end{Highlighting}
\end{Shaded}

\setcodelang{Python}

\begin{Shaded}
\begin{Highlighting}[]
\CommentTok{\# Python idiom — \_\_new\_\_ guard (GIL makes the simple version usually fine, but}
\CommentTok{\# add a lock for correctness under true parallelism / asyncio)}
\ImportTok{import}\NormalTok{ threading}
\KeywordTok{class}\NormalTok{ Config:}
\NormalTok{    \_instance }\OperatorTok{=} \VariableTok{None}
\NormalTok{    \_lock }\OperatorTok{=}\NormalTok{ threading.Lock()}
    \KeywordTok{def} \FunctionTok{\_\_new\_\_}\NormalTok{(cls):}
        \ControlFlowTok{if}\NormalTok{ cls.\_instance }\KeywordTok{is} \VariableTok{None}\NormalTok{:}
            \ControlFlowTok{with}\NormalTok{ cls.\_lock:}
                \ControlFlowTok{if}\NormalTok{ cls.\_instance }\KeywordTok{is} \VariableTok{None}\NormalTok{:}
\NormalTok{                    cls.\_instance }\OperatorTok{=} \BuiltInTok{super}\NormalTok{().}\FunctionTok{\_\_new\_\_}\NormalTok{(cls)}
        \ControlFlowTok{return}\NormalTok{ cls.\_instance}
\end{Highlighting}
\end{Shaded}

\textbf{When to use:} Genuinely-single shared resources --- connection/thread pools, caches, config, logging facades. \textbf{When NOT to use:} As a convenient global variable. Singletons are \emph{hidden} dependencies (a method that calls \texttt{Config.get()} doesn\textquotesingle t declare it needs config), they poison unit tests (can\textquotesingle t inject a mock; state leaks between tests), and they couple everything to a concrete class. \textbf{Prefer dependency injection of a single instance} managed by a container --- you keep "one instance" without the global-access pathologies.

\textbf{Real framework example:} \texttt{java.lang.Runtime.getRuntime()}, \texttt{java.awt.Desktop.getDesktop()}, Spring beans (singleton-scoped by default, but DI-managed so they\textquotesingle re \emph{testable} singletons).

\textbf{Trade-offs:} Controlled single instance, lazy init possible, global reach. Hidden global state, hard to test, thread-safety complexity, violates SRP if it accretes logic, can mask poor design ("everything is a singleton").

\subsection{Factory Method}\label{factory-method}

\textbf{Intent:} Define an interface for creating an object, but let \textbf{subclasses} decide which concrete class to instantiate. Deferring instantiation to subclasses.

\textbf{Problem it solves:} A \texttt{LogisticsApp} defines the workflow \texttt{planDelivery()} which needs a \texttt{Transport}. But the \emph{type} of transport depends on the subclass: \texttt{RoadLogistics} creates \texttt{Truck}, \texttt{SeaLogistics} creates \texttt{Ship}. The base class can\textquotesingle t hard-code \texttt{new\ Truck()} --- it would break sea logistics. So it calls an abstract \texttt{createTransport()} hook that subclasses override.

\textbf{Structure:}

\setcodelang{}

\begin{verbatim}
   «abstract» Creator                  «interface» Product
   + someOperation() {                       ▲ realizes
       Product p = createProduct();    ┌──────┴──────┐
       p.use();                     ConcreteA      ConcreteB
   }                                                
   + createProduct(): Product  ◄─── abstract "factory method"
        ▲ overridden by
   ┌────┴─────┐
ConcreteCreatorA    ConcreteCreatorB
 createProduct()=A   createProduct()=B
\end{verbatim}

\setcodelang{Java}

\begin{Shaded}
\begin{Highlighting}[]
\KeywordTok{interface}\NormalTok{ Transport }\OperatorTok{\{} \DataTypeTok{void} \FunctionTok{deliver}\OperatorTok{();} \OperatorTok{\}}
\KeywordTok{class}\NormalTok{ Truck }\KeywordTok{implements}\NormalTok{ Transport }\OperatorTok{\{} \KeywordTok{public} \DataTypeTok{void} \FunctionTok{deliver}\OperatorTok{()\{} \BuiltInTok{System}\OperatorTok{.}\FunctionTok{out}\OperatorTok{.}\FunctionTok{println}\OperatorTok{(}\StringTok{"by road"}\OperatorTok{);} \OperatorTok{\}} \OperatorTok{\}}
\KeywordTok{class}\NormalTok{ Ship  }\KeywordTok{implements}\NormalTok{ Transport }\OperatorTok{\{} \KeywordTok{public} \DataTypeTok{void} \FunctionTok{deliver}\OperatorTok{()\{} \BuiltInTok{System}\OperatorTok{.}\FunctionTok{out}\OperatorTok{.}\FunctionTok{println}\OperatorTok{(}\StringTok{"by sea"}\OperatorTok{);} \OperatorTok{\}} \OperatorTok{\}}

\KeywordTok{abstract} \KeywordTok{class}\NormalTok{ Logistics }\OperatorTok{\{}
    \KeywordTok{public} \DataTypeTok{void} \FunctionTok{planDelivery}\OperatorTok{()} \OperatorTok{\{}           \CommentTok{// the fixed workflow…}
\NormalTok{        Transport t }\OperatorTok{=} \FunctionTok{createTransport}\OperatorTok{();}   \CommentTok{// …uses the factory hook}
\NormalTok{        t}\OperatorTok{.}\FunctionTok{deliver}\OperatorTok{();}
    \OperatorTok{\}}
    \KeywordTok{protected} \KeywordTok{abstract}\NormalTok{ Transport }\FunctionTok{createTransport}\OperatorTok{();}   \CommentTok{// subclasses decide the type}
\OperatorTok{\}}
\KeywordTok{class}\NormalTok{ RoadLogistics }\KeywordTok{extends}\NormalTok{ Logistics }\OperatorTok{\{}
    \KeywordTok{protected}\NormalTok{ Transport }\FunctionTok{createTransport}\OperatorTok{()} \OperatorTok{\{} \ControlFlowTok{return} \KeywordTok{new} \FunctionTok{Truck}\OperatorTok{();} \OperatorTok{\}}
\OperatorTok{\}}
\KeywordTok{class}\NormalTok{ SeaLogistics }\KeywordTok{extends}\NormalTok{ Logistics }\OperatorTok{\{}
    \KeywordTok{protected}\NormalTok{ Transport }\FunctionTok{createTransport}\OperatorTok{()} \OperatorTok{\{} \ControlFlowTok{return} \KeywordTok{new} \FunctionTok{Ship}\OperatorTok{();} \OperatorTok{\}}
\OperatorTok{\}}
\end{Highlighting}
\end{Shaded}

\textbf{When to use:} When a class can\textquotesingle t anticipate the concrete type it must create, or you want subclasses to specify it. When you want to decouple object construction from its use (DIP/OCP). \textbf{When NOT to use:} When there\textquotesingle s only ever one concrete product --- a direct constructor is clearer (YAGNI).

\textbf{Real framework example:} \texttt{java.util.Calendar.getInstance()}, \texttt{java.text.NumberFormat.getInstance()}, the JDBC \texttt{Connection.createStatement()}, Django\textquotesingle s \texttt{Model.objects.create()}.

\textbf{Trade-offs:} Decouples client from concrete classes, single place to change creation, OCP-friendly. Requires a subclass per product variation, can proliferate classes for trivial creation.

\begin{quote}
\textbf{Factory Method vs. simple/static factory:} The GoF Factory Method uses \emph{inheritance} (subclass overrides the creator). A "simple factory" --- one \texttt{create(type)} method with a \texttt{switch} --- is not a GoF pattern but is common and fine for small cases. Mention the distinction in interviews.
\end{quote}

\subsection{Abstract Factory}\label{abstract-factory}

\textbf{Intent:} Provide an interface for creating \textbf{families of related objects} without specifying their concrete classes --- and guarantee the products are compatible with each other.

\textbf{Problem it solves:} A cross-platform UI must render a \texttt{Button} \emph{and} a \texttt{Checkbox}. On Windows both must be Windows-style; on macOS both must be Mac-style. You must never mix a \texttt{WinButton} with a \texttt{MacCheckbox}. An Abstract Factory bundles the creation of a \emph{whole family} so a single factory choice guarantees a consistent set.

\textbf{Structure:}

\setcodelang{}

\begin{verbatim}
        «interface» GUIFactory
        + createButton(): Button
        + createCheckbox(): Checkbox
              ▲ realizes
     ┌────────┴─────────┐
 WinFactory          MacFactory
  →WinButton           →MacButton
  →WinCheckbox         →MacCheckbox

Client depends ONLY on GUIFactory + Button + Checkbox interfaces.
Pick the factory once → every product matches.
\end{verbatim}

\setcodelang{Python}

\begin{Shaded}
\begin{Highlighting}[]
\ImportTok{from}\NormalTok{ abc }\ImportTok{import}\NormalTok{ ABC, abstractmethod}

\KeywordTok{class}\NormalTok{ Button(ABC):}
    \AttributeTok{@abstractmethod}
    \KeywordTok{def}\NormalTok{ render(}\VariableTok{self}\NormalTok{): ...}
\KeywordTok{class}\NormalTok{ Checkbox(ABC):}
    \AttributeTok{@abstractmethod}
    \KeywordTok{def}\NormalTok{ render(}\VariableTok{self}\NormalTok{): ...}

\KeywordTok{class}\NormalTok{ WinButton(Button):}
    \KeywordTok{def}\NormalTok{ render(}\VariableTok{self}\NormalTok{): }\BuiltInTok{print}\NormalTok{(}\StringTok{"[Win Button]"}\NormalTok{)}
\KeywordTok{class}\NormalTok{ WinCheckbox(Checkbox):}
    \KeywordTok{def}\NormalTok{ render(}\VariableTok{self}\NormalTok{): }\BuiltInTok{print}\NormalTok{(}\StringTok{"[Win Checkbox]"}\NormalTok{)}
\KeywordTok{class}\NormalTok{ MacButton(Button):}
    \KeywordTok{def}\NormalTok{ render(}\VariableTok{self}\NormalTok{): }\BuiltInTok{print}\NormalTok{(}\StringTok{"(Mac Button)"}\NormalTok{)}
\KeywordTok{class}\NormalTok{ MacCheckbox(Checkbox):}
    \KeywordTok{def}\NormalTok{ render(}\VariableTok{self}\NormalTok{): }\BuiltInTok{print}\NormalTok{(}\StringTok{"(Mac Checkbox)"}\NormalTok{)}

\KeywordTok{class}\NormalTok{ GUIFactory(ABC):}
    \AttributeTok{@abstractmethod}
    \KeywordTok{def}\NormalTok{ create\_button(}\VariableTok{self}\NormalTok{) }\OperatorTok{{-}\textgreater{}}\NormalTok{ Button: ...}
    \AttributeTok{@abstractmethod}
    \KeywordTok{def}\NormalTok{ create\_checkbox(}\VariableTok{self}\NormalTok{) }\OperatorTok{{-}\textgreater{}}\NormalTok{ Checkbox: ...}

\KeywordTok{class}\NormalTok{ WinFactory(GUIFactory):}
    \KeywordTok{def}\NormalTok{ create\_button(}\VariableTok{self}\NormalTok{): }\ControlFlowTok{return}\NormalTok{ WinButton()}
    \KeywordTok{def}\NormalTok{ create\_checkbox(}\VariableTok{self}\NormalTok{): }\ControlFlowTok{return}\NormalTok{ WinCheckbox()}
\KeywordTok{class}\NormalTok{ MacFactory(GUIFactory):}
    \KeywordTok{def}\NormalTok{ create\_button(}\VariableTok{self}\NormalTok{): }\ControlFlowTok{return}\NormalTok{ MacButton()}
    \KeywordTok{def}\NormalTok{ create\_checkbox(}\VariableTok{self}\NormalTok{): }\ControlFlowTok{return}\NormalTok{ MacCheckbox()}

\KeywordTok{def}\NormalTok{ build\_ui(factory: GUIFactory):       }\CommentTok{\# client is OS{-}agnostic}
\NormalTok{    factory.create\_button().render()}
\NormalTok{    factory.create\_checkbox().render()}
\end{Highlighting}
\end{Shaded}

\textbf{When to use:} Multiple product families that must stay internally consistent (UI toolkits per OS/theme, DB driver suites per vendor, cloud-provider SDKs). \textbf{When NOT to use:} Only one family exists today (YAGNI), or products aren\textquotesingle t actually related --- then a Factory Method per product is simpler.

\textbf{Real framework example:} \texttt{javax.xml.parsers.DocumentBuilderFactory}, \texttt{javax.xml.transform.TransformerFactory}, AWT\textquotesingle s \texttt{Toolkit}, cross-platform UI kits (Qt, Swing L\&F).

\textbf{Trade-offs:} Guarantees compatible product families, isolates concrete classes, swap a whole family by swapping one factory. Adding a \emph{new product type} (e.g., a \texttt{Slider}) forces editing every factory interface + impl --- the pattern is rigid against new product \emph{kinds} while flexible against new \emph{families}.

\begin{quote}
\textbf{Factory Method vs. Abstract Factory:} Factory Method creates \textbf{one} product via inheritance; Abstract Factory creates a \textbf{family} of related products via composition (a factory object with several create-methods). Abstract Factory is frequently \emph{implemented using} Factory Methods.
\end{quote}

\subsection{Builder}\label{builder}

\textbf{Intent:} Separate the construction of a complex object from its representation, so the same construction process can produce different representations --- and so an object with many parameters can be built step-by-step.

\textbf{Problem it solves:} \textbf{Telescoping constructors.} An object with 8 fields (3 required, 5 optional) forces either a dozen overloaded constructors (\texttt{Pizza(size)}, \texttt{Pizza(size,\ crust)}, \texttt{Pizza(size,\ crust,\ cheese)}\ldots) or a single constructor with 8 positional args you can\textquotesingle t read (\texttt{new\ Pizza("L",\ "thin",\ true,\ false,\ true,\ null,\ 2,\ "extra")} --- which boolean is which?). Builder gives a fluent, self-documenting, validated construction path and supports \textbf{immutable} results.

\textbf{Structure:}

\setcodelang{}

\begin{verbatim}
Client → Builder.size("L").crust("thin").cheese().build()  → Pizza (immutable)
            └ each step returns `this` (fluent chaining) ┘
            └ build() validates invariants, then constructs ┘
\end{verbatim}

\setcodelang{Java}

\begin{Shaded}
\begin{Highlighting}[]
\KeywordTok{public} \DataTypeTok{final} \KeywordTok{class}\NormalTok{ Pizza }\OperatorTok{\{}
    \KeywordTok{private} \DataTypeTok{final} \BuiltInTok{String}\NormalTok{ size}\OperatorTok{;}        \CommentTok{// immutable once built}
    \KeywordTok{private} \DataTypeTok{final} \DataTypeTok{boolean}\NormalTok{ cheese}\OperatorTok{,}\NormalTok{ pepperoni}\OperatorTok{;}
    \KeywordTok{private} \FunctionTok{Pizza}\OperatorTok{(}\NormalTok{Builder b}\OperatorTok{)} \OperatorTok{\{}        \CommentTok{// private ctor — only the Builder calls it}
        \KeywordTok{this}\OperatorTok{.}\FunctionTok{size} \OperatorTok{=}\NormalTok{ b}\OperatorTok{.}\FunctionTok{size}\OperatorTok{;} \KeywordTok{this}\OperatorTok{.}\FunctionTok{cheese} \OperatorTok{=}\NormalTok{ b}\OperatorTok{.}\FunctionTok{cheese}\OperatorTok{;} \KeywordTok{this}\OperatorTok{.}\FunctionTok{pepperoni} \OperatorTok{=}\NormalTok{ b}\OperatorTok{.}\FunctionTok{pepperoni}\OperatorTok{;}
    \OperatorTok{\}}
    \KeywordTok{public} \DataTypeTok{static} \KeywordTok{class}\NormalTok{ Builder }\OperatorTok{\{}
        \KeywordTok{private} \DataTypeTok{final} \BuiltInTok{String}\NormalTok{ size}\OperatorTok{;}    \CommentTok{// required → in Builder\textquotesingle{}s ctor}
        \KeywordTok{private} \DataTypeTok{boolean}\NormalTok{ cheese}\OperatorTok{,}\NormalTok{ pepperoni}\OperatorTok{;}          \CommentTok{// optional → defaults}
        \KeywordTok{public} \FunctionTok{Builder}\OperatorTok{(}\BuiltInTok{String}\NormalTok{ size}\OperatorTok{)} \OperatorTok{\{} \KeywordTok{this}\OperatorTok{.}\FunctionTok{size} \OperatorTok{=}\NormalTok{ size}\OperatorTok{;} \OperatorTok{\}}
        \KeywordTok{public}\NormalTok{ Builder }\FunctionTok{cheese}\OperatorTok{()}    \OperatorTok{\{} \KeywordTok{this}\OperatorTok{.}\FunctionTok{cheese} \OperatorTok{=} \KeywordTok{true}\OperatorTok{;} \ControlFlowTok{return} \KeywordTok{this}\OperatorTok{;} \OperatorTok{\}}     \CommentTok{// fluent}
        \KeywordTok{public}\NormalTok{ Builder }\FunctionTok{pepperoni}\OperatorTok{()} \OperatorTok{\{} \KeywordTok{this}\OperatorTok{.}\FunctionTok{pepperoni} \OperatorTok{=} \KeywordTok{true}\OperatorTok{;} \ControlFlowTok{return} \KeywordTok{this}\OperatorTok{;} \OperatorTok{\}}
        \KeywordTok{public}\NormalTok{ Pizza }\FunctionTok{build}\OperatorTok{()} \OperatorTok{\{}
            \ControlFlowTok{if} \OperatorTok{(}\NormalTok{size }\OperatorTok{==} \KeywordTok{null}\OperatorTok{)} \ControlFlowTok{throw} \KeywordTok{new} \BuiltInTok{IllegalStateException}\OperatorTok{(}\StringTok{"size required"}\OperatorTok{);}
            \ControlFlowTok{return} \KeywordTok{new} \FunctionTok{Pizza}\OperatorTok{(}\KeywordTok{this}\OperatorTok{);}                  \CommentTok{// validate{-}then{-}construct}
        \OperatorTok{\}}
    \OperatorTok{\}}
\OperatorTok{\}}
\CommentTok{// Usage — reads like prose, order{-}independent for optionals:}
\NormalTok{Pizza p }\OperatorTok{=} \KeywordTok{new}\NormalTok{ Pizza}\OperatorTok{.}\FunctionTok{Builder}\OperatorTok{(}\StringTok{"L"}\OperatorTok{).}\FunctionTok{cheese}\OperatorTok{().}\FunctionTok{pepperoni}\OperatorTok{().}\FunctionTok{build}\OperatorTok{();}
\end{Highlighting}
\end{Shaded}

\textbf{When to use:} Many optional/configurable parameters; you want immutability; construction needs validation or multiple steps; you want to build \emph{different representations} with the same steps (a \texttt{Director} orchestrating builders). \textbf{When NOT to use:} Objects with 2--3 fields --- a plain constructor or a small record is simpler.

\textbf{Real framework example:} \texttt{StringBuilder}, \texttt{java.net.http.HttpRequest.Builder}, \texttt{Stream.Builder}, Lombok \texttt{@Builder}, Guava\textquotesingle s \texttt{ImmutableList.builder()}, Kotlin\textquotesingle s named/default args (a language feature that \emph{replaces} Builder), protobuf message builders.

\textbf{Trade-offs:} Readable construction, immutability, validation in one place, no telescoping ctors. More boilerplate (a whole Builder class), slight object-allocation overhead, overkill for simple objects.

\subsection{Prototype}\label{prototype}

\textbf{Intent:} Specify the kinds of objects to create using a prototypical instance, and create new objects by \textbf{copying (cloning)} this prototype rather than constructing from scratch.

\textbf{Problem it solves:} Object creation is expensive (a DB round-trip, heavy computation, deserialization) or the exact class is decided at runtime. Instead of re-running the expensive setup, you keep a pre-built prototype and \texttt{clone()} it, then tweak the copy. Also useful when a system must be independent of how its products are created.

\textbf{Structure:}

\setcodelang{}

\begin{verbatim}
«interface» Prototype { clone(): Prototype }
        ▲
   ConcreteShape { clone(): returns a copy of itself }

registry: { "warrior": warriorPrototype, "mage": magePrototype }
spawn("warrior") → warriorPrototype.clone()   (cheap copy, no DB hit)
\end{verbatim}

\setcodelang{Python}

\begin{Shaded}
\begin{Highlighting}[]
\ImportTok{import}\NormalTok{ copy}

\KeywordTok{class}\NormalTok{ GameEnemy:}
    \KeywordTok{def} \FunctionTok{\_\_init\_\_}\NormalTok{(}\VariableTok{self}\NormalTok{, sprite, stats, ai\_tree):}
        \VariableTok{self}\NormalTok{.sprite }\OperatorTok{=}\NormalTok{ sprite          }\CommentTok{\# imagine: loaded from disk (expensive)}
        \VariableTok{self}\NormalTok{.stats }\OperatorTok{=}\NormalTok{ stats}
        \VariableTok{self}\NormalTok{.ai\_tree }\OperatorTok{=}\NormalTok{ ai\_tree        }\CommentTok{\# imagine: compiled behavior tree (expensive)}

    \KeywordTok{def}\NormalTok{ clone(}\VariableTok{self}\NormalTok{) }\OperatorTok{{-}\textgreater{}} \StringTok{"GameEnemy"}\NormalTok{:}
        \ControlFlowTok{return}\NormalTok{ copy.deepcopy(}\VariableTok{self}\NormalTok{)    }\CommentTok{\# deep vs shallow matters! (see trade{-}offs)}

\CommentTok{\# Build the costly prototype ONCE, then spawn a wave cheaply:}
\NormalTok{orc\_prototype }\OperatorTok{=}\NormalTok{ GameEnemy(load\_sprite(}\StringTok{"orc.png"}\NormalTok{), \{}\StringTok{"hp"}\NormalTok{: }\DecValTok{50}\NormalTok{\}, compile\_ai(}\StringTok{"orc.bt"}\NormalTok{))}
\NormalTok{wave }\OperatorTok{=}\NormalTok{ [orc\_prototype.clone() }\ControlFlowTok{for}\NormalTok{ \_ }\KeywordTok{in} \BuiltInTok{range}\NormalTok{(}\DecValTok{100}\NormalTok{)]   }\CommentTok{\# no disk/compile per enemy}
\NormalTok{wave[}\DecValTok{0}\NormalTok{].stats[}\StringTok{"hp"}\NormalTok{] }\OperatorTok{=} \DecValTok{80}                              \CommentTok{\# tweak the copy independently}
\end{Highlighting}
\end{Shaded}

\textbf{When to use:} Costly construction you can amortize via copy; many similar objects differing in a few fields; the concrete type is determined at runtime (clone a registered prototype by key). \textbf{When NOT to use:} Cheap-to-construct objects, or objects with complex reference graphs / external resources (file handles, sockets) where copying semantics are murky.

\textbf{Real framework example:} Java\textquotesingle s \texttt{Cloneable} / \texttt{Object.clone()}, JavaScript\textquotesingle s \texttt{Object.create(proto)} and prototypal inheritance, undo systems that snapshot via clone, ML pipeline \texttt{clone()} (scikit-learn \texttt{clone(estimator)} copies hyperparameters).

\textbf{Trade-offs:} Avoids expensive re-init, decouples from concrete class, runtime-configurable. \textbf{Shallow vs. deep copy is the classic bug} --- a shallow clone shares mutable sub-objects (two "copies" mutating the same list). Cloning objects with cycles or external resources is hard. \texttt{Cloneable} in Java is widely considered a flawed API.

\subsection{Creational Patterns --- Comparison}\label{creational-patterns--comparison}

\begin{longtable}[]{@{}
  >{\raggedright\arraybackslash}p{(\columnwidth - 4\tabcolsep) * \real{0.1592}}
  >{\raggedright\arraybackslash}p{(\columnwidth - 4\tabcolsep) * \real{0.4114}}
  >{\raggedright\arraybackslash}p{(\columnwidth - 4\tabcolsep) * \real{0.4294}}@{}}
\toprule\noalign{}
\rowcolor{acc!16}
\begin{minipage}[b]{\linewidth}\raggedright
Pattern
\end{minipage} & \begin{minipage}[b]{\linewidth}\raggedright
What it abstracts
\end{minipage} & \begin{minipage}[b]{\linewidth}\raggedright
Reach for it when
\end{minipage} \\
\midrule\noalign{}
\endhead
\bottomrule\noalign{}
\endlastfoot
\textbf{Singleton} & \emph{how many} instances (exactly one) & Logger, config, pool --- sparingly; prefer DI \\
\rowcolor{zebra}
\textbf{Factory Method} & \emph{which subclass} creates the product & One product, concrete type chosen by subclass \\
\textbf{Abstract Factory} & a \emph{family} of compatible products & Cross-platform/pluggable suites that must match \\
\rowcolor{zebra}
\textbf{Builder} & \emph{step-by-step} assembly of one complex object & Many optional fields, immutability, validation \\
\textbf{Prototype} & creation via \emph{copy} of an exemplar & Expensive construction, runtime-decided type \\
\end{longtable}

\section{Structural Patterns}\label{structural-patterns}

Structural patterns are about \textbf{composing classes and objects into larger structures} while keeping those structures flexible and efficient. The recurring theme: a wrapper or tree that lets objects collaborate without rigid coupling.

\subsection{Adapter}\label{adapter}

\textbf{Intent:} Convert the interface of a class into another interface clients expect. Lets classes with incompatible interfaces work together.

\textbf{Problem it solves:} Your code is written against a \texttt{PaymentGateway} interface with \texttt{pay(amount)}. You must integrate a third-party \texttt{StripeSdk} that exposes \texttt{makeCharge(cents,\ currency)}. You can\textquotesingle t edit the SDK and you don\textquotesingle t want to rewrite your call sites. An Adapter wraps the SDK and presents the \texttt{PaymentGateway} interface, translating each call.

\textbf{Structure:}

\setcodelang{}

\begin{verbatim}
Client → «interface» Target.request()
                 ▲ realizes
            Adapter ──holds──> Adaptee (incompatible 3rd-party class)
            request() { adaptee.specificRequest(); }   // translates the call
\end{verbatim}

\setcodelang{Java}

\begin{Shaded}
\begin{Highlighting}[]
\KeywordTok{interface}\NormalTok{ PaymentGateway }\OperatorTok{\{} \DataTypeTok{boolean} \FunctionTok{pay}\OperatorTok{(}\DataTypeTok{double}\NormalTok{ dollars}\OperatorTok{);} \OperatorTok{\}}   \CommentTok{// what our code expects}

\KeywordTok{class}\NormalTok{ StripeSdk }\OperatorTok{\{}                                           \CommentTok{// 3rd{-}party, can\textquotesingle{}t change}
    \DataTypeTok{boolean} \FunctionTok{makeCharge}\OperatorTok{(}\DataTypeTok{long}\NormalTok{ cents}\OperatorTok{,} \BuiltInTok{String}\NormalTok{ currency}\OperatorTok{)} \OperatorTok{\{} \CommentTok{/* ... */} \ControlFlowTok{return} \KeywordTok{true}\OperatorTok{;} \OperatorTok{\}}
\OperatorTok{\}}

\KeywordTok{class}\NormalTok{ StripeAdapter }\KeywordTok{implements}\NormalTok{ PaymentGateway }\OperatorTok{\{}            \CommentTok{// the bridge}
    \KeywordTok{private} \DataTypeTok{final}\NormalTok{ StripeSdk stripe}\OperatorTok{;}
    \FunctionTok{StripeAdapter}\OperatorTok{(}\NormalTok{StripeSdk s}\OperatorTok{)} \OperatorTok{\{} \KeywordTok{this}\OperatorTok{.}\FunctionTok{stripe} \OperatorTok{=}\NormalTok{ s}\OperatorTok{;} \OperatorTok{\}}
    \KeywordTok{public} \DataTypeTok{boolean} \FunctionTok{pay}\OperatorTok{(}\DataTypeTok{double}\NormalTok{ dollars}\OperatorTok{)} \OperatorTok{\{}
        \ControlFlowTok{return}\NormalTok{ stripe}\OperatorTok{.}\FunctionTok{makeCharge}\OperatorTok{(}\BuiltInTok{Math}\OperatorTok{.}\FunctionTok{round}\OperatorTok{(}\NormalTok{dollars }\OperatorTok{*} \DecValTok{100}\OperatorTok{),} \StringTok{"USD"}\OperatorTok{);}  \CommentTok{// translate}
    \OperatorTok{\}}
\OperatorTok{\}}
\CommentTok{// Client code stays clean: PaymentGateway g = new StripeAdapter(new StripeSdk());}
\end{Highlighting}
\end{Shaded}

\textbf{When to use:} Integrating third-party libs, legacy code, or two subsystems with mismatched interfaces; making old code satisfy a new interface without rewriting it. \textbf{When NOT to use:} When you control both sides --- just align the interfaces. When the impedance mismatch is so large the adapter would reimplement the adaptee (that\textquotesingle s a rewrite, not an adapter).

\textbf{Real framework example:} \texttt{java.io.InputStreamReader} (adapts byte \texttt{InputStream} \(\rightarrow\) char \texttt{Reader}), \texttt{java.util.Arrays.asList()} (array \(\rightarrow\) List), \texttt{java.util.Collections.list(Enumeration)} (old \texttt{Enumeration} \(\rightarrow\) modern \texttt{List}), Slf4j bridges (\texttt{slf4j-log4j12}).

\textbf{Trade-offs:} Reuse incompatible code, isolate the dependency in one place, SRP (translation logic localized). Extra indirection; deep adapter chains get confusing; can hide a mismatch that should be fixed at the source.

\begin{quote}
\textbf{Object Adapter (composition, shown above) vs. Class Adapter (multiple inheritance).} Object adapter is preferred and the only option in single-inheritance languages like Java for combining with concrete adaptees.
\end{quote}

\subsection{Decorator}\label{decorator}

\textbf{Intent:} Attach additional responsibilities to an object \textbf{dynamically}. A flexible alternative to subclassing for extending functionality.

\textbf{Problem it solves:} You have a \texttt{DataSource} and want optional, \emph{combinable} features: compression, encryption, buffering. With inheritance you\textquotesingle d need a class per combination --- \texttt{CompressedSource}, \texttt{EncryptedSource}, \texttt{CompressedEncryptedSource}, \texttt{BufferedCompressedEncryptedSource}\ldots{} an exponential explosion. Decorator wraps the object in layers, each adding one responsibility, composed at runtime in any order.

\textbf{Structure:}

\setcodelang{}

\begin{verbatim}
        «interface» Component { operation() }
              ▲                      ▲
        ConcreteComponent      Decorator ──wraps──> Component
        (the core object)         operation(){ extra(); wrapped.operation(); }
                                       ▲
                          ┌────────────┼────────────┐
                   BufferDecorator  GzipDecorator  CryptoDecorator
\end{verbatim}

\setcodelang{Java}

\begin{Shaded}
\begin{Highlighting}[]
\KeywordTok{interface} \BuiltInTok{DataSource} \OperatorTok{\{} \DataTypeTok{void} \FunctionTok{write}\OperatorTok{(}\BuiltInTok{String}\NormalTok{ data}\OperatorTok{);} \BuiltInTok{String} \FunctionTok{read}\OperatorTok{();} \OperatorTok{\}}

\KeywordTok{class} \BuiltInTok{FileDataSource} \KeywordTok{implements} \BuiltInTok{DataSource} \OperatorTok{\{}                 \CommentTok{// the core}
    \KeywordTok{public} \DataTypeTok{void} \FunctionTok{write}\OperatorTok{(}\BuiltInTok{String}\NormalTok{ d}\OperatorTok{)\{} \CommentTok{/* to file */} \OperatorTok{\}}
    \KeywordTok{public} \BuiltInTok{String} \FunctionTok{read}\OperatorTok{()\{} \ControlFlowTok{return} \StringTok{"raw"}\OperatorTok{;} \OperatorTok{\}}
\OperatorTok{\}}

\KeywordTok{abstract} \KeywordTok{class}\NormalTok{ DataSourceDecorator }\KeywordTok{implements} \BuiltInTok{DataSource} \OperatorTok{\{}   \CommentTok{// base wrapper}
    \KeywordTok{protected} \DataTypeTok{final} \BuiltInTok{DataSource}\NormalTok{ wrappee}\OperatorTok{;}
    \FunctionTok{DataSourceDecorator}\OperatorTok{(}\BuiltInTok{DataSource}\NormalTok{ s}\OperatorTok{)\{} \KeywordTok{this}\OperatorTok{.}\FunctionTok{wrappee} \OperatorTok{=}\NormalTok{ s}\OperatorTok{;} \OperatorTok{\}}
    \KeywordTok{public} \DataTypeTok{void} \FunctionTok{write}\OperatorTok{(}\BuiltInTok{String}\NormalTok{ d}\OperatorTok{)\{}\NormalTok{ wrappee}\OperatorTok{.}\FunctionTok{write}\OperatorTok{(}\NormalTok{d}\OperatorTok{);} \OperatorTok{\}}
    \KeywordTok{public} \BuiltInTok{String} \FunctionTok{read}\OperatorTok{()\{} \ControlFlowTok{return}\NormalTok{ wrappee}\OperatorTok{.}\FunctionTok{read}\OperatorTok{();} \OperatorTok{\}}
\OperatorTok{\}}

\KeywordTok{class}\NormalTok{ EncryptionDecorator }\KeywordTok{extends}\NormalTok{ DataSourceDecorator }\OperatorTok{\{}
    \FunctionTok{EncryptionDecorator}\OperatorTok{(}\BuiltInTok{DataSource}\NormalTok{ s}\OperatorTok{)\{} \KeywordTok{super}\OperatorTok{(}\NormalTok{s}\OperatorTok{);} \OperatorTok{\}}
    \KeywordTok{public} \DataTypeTok{void} \FunctionTok{write}\OperatorTok{(}\BuiltInTok{String}\NormalTok{ d}\OperatorTok{)\{} \KeywordTok{super}\OperatorTok{.}\FunctionTok{write}\OperatorTok{(}\FunctionTok{encrypt}\OperatorTok{(}\NormalTok{d}\OperatorTok{));} \OperatorTok{\}}   \CommentTok{// add behavior, then delegate}
    \KeywordTok{public} \BuiltInTok{String} \FunctionTok{read}\OperatorTok{()\{} \ControlFlowTok{return} \FunctionTok{decrypt}\OperatorTok{(}\KeywordTok{super}\OperatorTok{.}\FunctionTok{read}\OperatorTok{());} \OperatorTok{\}}
    \KeywordTok{private} \BuiltInTok{String} \FunctionTok{encrypt}\OperatorTok{(}\BuiltInTok{String}\NormalTok{ s}\OperatorTok{)\{} \ControlFlowTok{return} \StringTok{"enc("} \OperatorTok{+}\NormalTok{ s }\OperatorTok{+} \StringTok{")"}\OperatorTok{;} \OperatorTok{\}}
    \KeywordTok{private} \BuiltInTok{String} \FunctionTok{decrypt}\OperatorTok{(}\BuiltInTok{String}\NormalTok{ s}\OperatorTok{)\{} \ControlFlowTok{return}\NormalTok{ s}\OperatorTok{;} \OperatorTok{\}}
\OperatorTok{\}}
\CommentTok{// Stack at runtime, any order:}
\BuiltInTok{DataSource}\NormalTok{ src }\OperatorTok{=} \KeywordTok{new} \FunctionTok{EncryptionDecorator}\OperatorTok{(}\KeywordTok{new} \FunctionTok{CompressionDecorator}\OperatorTok{(}\KeywordTok{new} \BuiltInTok{FileDataSource}\OperatorTok{()));}
\end{Highlighting}
\end{Shaded}

\textbf{The canonical Java I/O Streams example} (the textbook real-world Decorator):

\setcodelang{Java}

\begin{Shaded}
\begin{Highlighting}[]
\CommentTok{// Each constructor wraps the previous, adding one capability:}
\BuiltInTok{BufferedReader}\NormalTok{ r }\OperatorTok{=} \KeywordTok{new} \BuiltInTok{BufferedReader}\OperatorTok{(}      \CommentTok{// adds buffering}
    \KeywordTok{new} \BuiltInTok{InputStreamReader}\OperatorTok{(}                  \CommentTok{// adapts bytes→chars (Adapter, actually)}
        \KeywordTok{new} \BuiltInTok{FileInputStream}\OperatorTok{(}\StringTok{"data.txt"}\OperatorTok{)));}  \CommentTok{// the core byte source}
\CommentTok{// Compose freely: GZIPInputStream, CipherInputStream, DataInputStream … all decorators.}
\end{Highlighting}
\end{Shaded}

\textbf{When to use:} Optional, stackable, order-flexible responsibilities; you want to add behavior without touching the original class or exploding subclasses (OCP). \textbf{When NOT to use:} When you need to inspect the \emph{concrete} wrapped type at runtime (decorators hide it); when there\textquotesingle s only one fixed feature (just subclass); when ordering bugs would be subtle.

\textbf{Real framework example:} \textbf{Java I/O streams} (the definitive example), \texttt{java.util.Collections.synchronizedList()} / \texttt{unmodifiableList()}, Python\textquotesingle s \texttt{@functools.lru\_cache} and other function decorators, HTTP middleware chains (Express, Django middleware), \texttt{javax.servlet.http.HttpServletRequestWrapper}.

\textbf{Trade-offs:} Runtime composition, avoids subclass explosion, single-responsibility layers, OCP. Many small objects (debugging a 5-layer stack is painful), order can matter and be non-obvious, hard to remove a specific decorator from the middle of a stack, identity checks (\texttt{instanceof}) break through wrappers.

\begin{quote}
\textbf{Decorator vs. Inheritance:} inheritance fixes the augmentation at compile time and can\textquotesingle t combine; Decorator composes at runtime and stacks arbitrarily. \textbf{Decorator vs. Proxy:} identical structure (both wrap, same interface) --- \emph{intent} differs: Decorator \emph{adds responsibilities}; Proxy \emph{controls access}.
\end{quote}

\subsection{Facade}\label{facade}

\textbf{Intent:} Provide a unified, simplified interface to a set of interfaces in a subsystem. Makes the subsystem easier to use.

\textbf{Problem it solves:} Converting a video involves codecs, audio mixers, bitrate readers, format detectors --- a dozen subsystem classes with intricate setup order. Most clients just want \texttt{convert(file,\ "mp4")}. A Facade exposes one simple method and orchestrates the messy internals, decoupling clients from the subsystem.

\textbf{Structure:}

\setcodelang{}

\begin{verbatim}
Client → Facade.watchMovie()
              │ orchestrates:
              ├─ amp.on(); amp.setVolume(5)
              ├─ projector.on(); projector.wideScreen()
              └─ player.on(); player.play(movie)
   (subsystem classes remain usable directly for power users)
\end{verbatim}

\setcodelang{Python}

\begin{Shaded}
\begin{Highlighting}[]
\KeywordTok{class}\NormalTok{ Amplifier:}
    \KeywordTok{def}\NormalTok{ on(}\VariableTok{self}\NormalTok{): }\BuiltInTok{print}\NormalTok{(}\StringTok{"Amp on"}\NormalTok{)}
    \KeywordTok{def}\NormalTok{ set\_volume(}\VariableTok{self}\NormalTok{, v): }\BuiltInTok{print}\NormalTok{(}\SpecialStringTok{f"Volume }\SpecialCharTok{\{}\NormalTok{v}\SpecialCharTok{\}}\SpecialStringTok{"}\NormalTok{)}
\KeywordTok{class}\NormalTok{ Projector:}
    \KeywordTok{def}\NormalTok{ on(}\VariableTok{self}\NormalTok{): }\BuiltInTok{print}\NormalTok{(}\StringTok{"Projector on"}\NormalTok{)}
    \KeywordTok{def}\NormalTok{ wide(}\VariableTok{self}\NormalTok{): }\BuiltInTok{print}\NormalTok{(}\StringTok{"Wide screen"}\NormalTok{)}
\KeywordTok{class}\NormalTok{ Player:}
    \KeywordTok{def}\NormalTok{ on(}\VariableTok{self}\NormalTok{): }\BuiltInTok{print}\NormalTok{(}\StringTok{"Player on"}\NormalTok{)}
    \KeywordTok{def}\NormalTok{ play(}\VariableTok{self}\NormalTok{, m): }\BuiltInTok{print}\NormalTok{(}\SpecialStringTok{f"Playing }\SpecialCharTok{\{}\NormalTok{m}\SpecialCharTok{\}}\SpecialStringTok{"}\NormalTok{)}

\KeywordTok{class}\NormalTok{ HomeTheaterFacade:                 }\CommentTok{\# one simple front}
    \KeywordTok{def} \FunctionTok{\_\_init\_\_}\NormalTok{(}\VariableTok{self}\NormalTok{, amp, proj, player):}
        \VariableTok{self}\NormalTok{.\_amp, }\VariableTok{self}\NormalTok{.\_proj, }\VariableTok{self}\NormalTok{.\_player }\OperatorTok{=}\NormalTok{ amp, proj, player}
    \KeywordTok{def}\NormalTok{ watch\_movie(}\VariableTok{self}\NormalTok{, movie):        }\CommentTok{\# hides ordering \& coordination}
        \VariableTok{self}\NormalTok{.\_amp.on()}\OperatorTok{;} \VariableTok{self}\NormalTok{.\_amp.set\_volume(}\DecValTok{5}\NormalTok{)}
        \VariableTok{self}\NormalTok{.\_proj.on()}\OperatorTok{;} \VariableTok{self}\NormalTok{.\_proj.wide()}
        \VariableTok{self}\NormalTok{.\_player.on()}\OperatorTok{;} \VariableTok{self}\NormalTok{.\_player.play(movie)}

\NormalTok{HomeTheaterFacade(Amplifier(), Projector(), Player()).watch\_movie(}\StringTok{"Inception"}\NormalTok{)}
\end{Highlighting}
\end{Shaded}

\textbf{When to use:} A complex subsystem with intricate APIs; you want a simple entry point and to reduce client\(\leftrightarrow\)subsystem coupling; layering an architecture (each layer exposes a facade to the next). \textbf{When NOT to use:} When the subsystem is already simple (the facade adds a useless layer), or when clients genuinely need fine-grained control (then expose both).

\textbf{Real framework example:} SLF4J (\texttt{LoggerFactory.getLogger()} fronts log4j/logback), \texttt{javax.faces.context.FacesContext}, Spring\textquotesingle s \texttt{JdbcTemplate} (facade over raw JDBC), AWS SDK service clients, a REST gateway fronting microservices.

\textbf{Trade-offs:} Simpler client code, decouples clients from subsystem internals, defines a clear layer boundary. Can become a God object if it accretes logic; risks hiding power users need; another layer to maintain.

\begin{quote}
\textbf{Facade vs. Adapter:} Adapter makes one interface \emph{match} an expected one (interface conversion); Facade \emph{simplifies} many interfaces into a new, easier one (interface reduction). Facade is about ease; Adapter is about compatibility.
\end{quote}

\subsection{Proxy}\label{proxy}

\textbf{Intent:} Provide a surrogate or placeholder for another object to \textbf{control access} to it --- same interface as the real subject, with added control logic.

\textbf{Problem it solves:} You need to do \emph{something around} access to an object without changing the object or its clients: defer its expensive creation, check permissions, cache results, log, or reach it across a network. A Proxy implements the same interface and inserts that logic.

\textbf{Flavors:}

\begin{longtable}[]{@{}
  >{\raggedright\arraybackslash}p{(\columnwidth - 4\tabcolsep) * \real{0.1521}}
  >{\raggedright\arraybackslash}p{(\columnwidth - 4\tabcolsep) * \real{0.4376}}
  >{\raggedright\arraybackslash}p{(\columnwidth - 4\tabcolsep) * \real{0.4102}}@{}}
\toprule\noalign{}
\rowcolor{acc!16}
\begin{minipage}[b]{\linewidth}\raggedright
Flavor
\end{minipage} & \begin{minipage}[b]{\linewidth}\raggedright
Purpose
\end{minipage} & \begin{minipage}[b]{\linewidth}\raggedright
Example
\end{minipage} \\
\midrule\noalign{}
\endhead
\bottomrule\noalign{}
\endlastfoot
\textbf{Virtual} & Lazy-create an expensive object on first use & High-res image not loaded until displayed \\
\rowcolor{zebra}
\textbf{Protection} & Access control / authorization & Reject \texttt{delete()} unless caller is admin \\
\textbf{Remote} & Local stand-in for an object in another process & RMI/gRPC client stub \\
\rowcolor{zebra}
\textbf{Caching} & Memoize expensive results & Cache HTTP/DB responses \\
\textbf{Smart reference} & Extra bookkeeping (ref-count, logging) & Hibernate lazy proxies, AOP \\
\end{longtable}

\textbf{Structure:}

\setcodelang{}

\begin{verbatim}
Client → «interface» Subject { request() }
              ▲                    ▲
         RealSubject            Proxy ──holds──> RealSubject
         (heavy/remote)         request(){ check/lazy/cache; real.request(); }
\end{verbatim}

\setcodelang{Python}

\begin{Shaded}
\begin{Highlighting}[]
\ImportTok{from}\NormalTok{ abc }\ImportTok{import}\NormalTok{ ABC, abstractmethod}

\KeywordTok{class}\NormalTok{ Image(ABC):}
    \AttributeTok{@abstractmethod}
    \KeywordTok{def}\NormalTok{ display(}\VariableTok{self}\NormalTok{): ...}

\KeywordTok{class}\NormalTok{ RealImage(Image):}
    \KeywordTok{def} \FunctionTok{\_\_init\_\_}\NormalTok{(}\VariableTok{self}\NormalTok{, filename):}
        \VariableTok{self}\NormalTok{.filename }\OperatorTok{=}\NormalTok{ filename}
        \VariableTok{self}\NormalTok{.\_load\_from\_disk()          }\CommentTok{\# EXPENSIVE — happens at construction}
    \KeywordTok{def}\NormalTok{ \_load\_from\_disk(}\VariableTok{self}\NormalTok{): }\BuiltInTok{print}\NormalTok{(}\SpecialStringTok{f"Loading }\SpecialCharTok{\{}\VariableTok{self}\SpecialCharTok{.}\NormalTok{filename}\SpecialCharTok{\}}\SpecialStringTok{ (heavy I/O)"}\NormalTok{)}
    \KeywordTok{def}\NormalTok{ display(}\VariableTok{self}\NormalTok{): }\BuiltInTok{print}\NormalTok{(}\SpecialStringTok{f"Displaying }\SpecialCharTok{\{}\VariableTok{self}\SpecialCharTok{.}\NormalTok{filename}\SpecialCharTok{\}}\SpecialStringTok{"}\NormalTok{)}

\KeywordTok{class}\NormalTok{ ImageProxy(Image):                 }\CommentTok{\# VIRTUAL proxy — defers the heavy load}
    \KeywordTok{def} \FunctionTok{\_\_init\_\_}\NormalTok{(}\VariableTok{self}\NormalTok{, filename):}
        \VariableTok{self}\NormalTok{.filename }\OperatorTok{=}\NormalTok{ filename}
        \VariableTok{self}\NormalTok{.\_real }\OperatorTok{=} \VariableTok{None}
    \KeywordTok{def}\NormalTok{ display(}\VariableTok{self}\NormalTok{):}
        \ControlFlowTok{if} \VariableTok{self}\NormalTok{.\_real }\KeywordTok{is} \VariableTok{None}\NormalTok{:           }\CommentTok{\# create only on first actual use}
            \VariableTok{self}\NormalTok{.\_real }\OperatorTok{=}\NormalTok{ RealImage(}\VariableTok{self}\NormalTok{.filename)}
        \VariableTok{self}\NormalTok{.\_real.display()}

\NormalTok{img }\OperatorTok{=}\NormalTok{ ImageProxy(}\StringTok{"huge.png"}\NormalTok{)   }\CommentTok{\# no disk I/O yet}
\NormalTok{img.display()                  }\CommentTok{\# NOW it loads, then displays}
\NormalTok{img.display()                  }\CommentTok{\# already loaded — just displays}
\end{Highlighting}
\end{Shaded}

\textbf{When to use:} Lazy initialization of heavy objects, access control, remote access, caching, instrumentation --- any "do work around access" need. \textbf{When NOT to use:} When the added latency/indirection isn\textquotesingle t justified; when a Decorator (adds behavior) or a cache library would be clearer for the specific need.

\textbf{Real framework example:} Hibernate lazy-loading proxies, Spring AOP / \texttt{@Transactional} (dynamic proxy wraps the bean to open/commit transactions), \texttt{java.lang.reflect.Proxy} (dynamic proxies), gRPC/RMI stubs, mock objects in tests (a proxy that records calls).

\textbf{Trade-offs:} Control access transparently, lazy/remote/caching without changing clients, separation of concerns. Extra indirection/latency, response-time can become unpredictable (lazy load surprises), more classes; remote proxies hide network failures behind a local-looking call.

\subsection{Composite}\label{composite}

\textbf{Intent:} Compose objects into \textbf{tree structures} to represent part-whole hierarchies. Lets clients treat individual objects and compositions of objects \textbf{uniformly}.

\textbf{Problem it solves:} A file system has files (leaves) and directories (which contain files \emph{and} directories). You want \texttt{size()} or \texttt{print()} to work the same whether you call it on a single file or a whole directory tree --- without the client writing \texttt{if\ (isDirectory)\ …\ else\ …} everywhere. Composite gives leaves and containers a common interface and lets containers recurse into children.

\textbf{Structure:}

\setcodelang{}

\begin{verbatim}
        «interface» Component { operation() }
              ▲                       ▲
            Leaf                  Composite ──contains 0..*──> Component
        operation()              operation(){ for c in children: c.operation() }
                                 add(c) / remove(c)
\end{verbatim}

\setcodelang{Python}

\begin{Shaded}
\begin{Highlighting}[]
\ImportTok{from}\NormalTok{ abc }\ImportTok{import}\NormalTok{ ABC, abstractmethod}

\KeywordTok{class}\NormalTok{ FsNode(ABC):}
    \AttributeTok{@abstractmethod}
    \KeywordTok{def}\NormalTok{ size(}\VariableTok{self}\NormalTok{) }\OperatorTok{{-}\textgreater{}} \BuiltInTok{int}\NormalTok{: ...}

\KeywordTok{class}\NormalTok{ File(FsNode):                      }\CommentTok{\# LEAF}
    \KeywordTok{def} \FunctionTok{\_\_init\_\_}\NormalTok{(}\VariableTok{self}\NormalTok{, name, bytes\_): }\VariableTok{self}\NormalTok{.name, }\VariableTok{self}\NormalTok{.\_bytes }\OperatorTok{=}\NormalTok{ name, bytes\_}
    \KeywordTok{def}\NormalTok{ size(}\VariableTok{self}\NormalTok{): }\ControlFlowTok{return} \VariableTok{self}\NormalTok{.\_bytes}

\KeywordTok{class}\NormalTok{ Directory(FsNode):                  }\CommentTok{\# COMPOSITE}
    \KeywordTok{def} \FunctionTok{\_\_init\_\_}\NormalTok{(}\VariableTok{self}\NormalTok{, name): }\VariableTok{self}\NormalTok{.name, }\VariableTok{self}\NormalTok{.\_children }\OperatorTok{=}\NormalTok{ name, []}
    \KeywordTok{def}\NormalTok{ add(}\VariableTok{self}\NormalTok{, node: FsNode): }\VariableTok{self}\NormalTok{.\_children.append(node)}
    \KeywordTok{def}\NormalTok{ size(}\VariableTok{self}\NormalTok{):                       }\CommentTok{\# recurse uniformly — no type checks}
        \ControlFlowTok{return} \BuiltInTok{sum}\NormalTok{(child.size() }\ControlFlowTok{for}\NormalTok{ child }\KeywordTok{in} \VariableTok{self}\NormalTok{.\_children)}

\NormalTok{root }\OperatorTok{=}\NormalTok{ Directory(}\StringTok{"root"}\NormalTok{)}
\NormalTok{root.add(File(}\StringTok{"a.txt"}\NormalTok{, }\DecValTok{100}\NormalTok{))}
\NormalTok{docs }\OperatorTok{=}\NormalTok{ Directory(}\StringTok{"docs"}\NormalTok{)}\OperatorTok{;}\NormalTok{ docs.add(File(}\StringTok{"b.pdf"}\NormalTok{, }\DecValTok{2000}\NormalTok{))}
\NormalTok{root.add(docs)}
\BuiltInTok{print}\NormalTok{(root.size())   }\CommentTok{\# 2100 — works the same on file or directory}
\end{Highlighting}
\end{Shaded}

\textbf{When to use:} Genuine part-whole hierarchies (file systems, UI widget trees, org charts, arithmetic expression trees, graphics scene graphs). \textbf{When NOT to use:} Flat collections (no tree) or when leaves and composites are too different to share a meaningful interface (forcing \texttt{add()} onto a Leaf is awkward --- a known design tension).

\textbf{Real framework example:} \texttt{java.awt.Container} / Swing \texttt{JComponent} trees, the DOM (nodes \& elements), React component trees, \texttt{java.util.Map.Entry} collections, Composite UI in Android \texttt{ViewGroup}/\texttt{View}.

\textbf{Trade-offs:} Uniform treatment of leaf/composite, easy to add new node types, natural recursion. The shared interface can over-generalize (a \texttt{File} shouldn\textquotesingle t really have \texttt{add()}); can make the design too permissive (type safety lost); deep trees risk stack overflow on naive recursion.

\subsection{Bridge}\label{bridge}

\textbf{Intent:} Decouple an \textbf{abstraction} from its \textbf{implementation} so the two can vary independently --- avoiding a Cartesian-product class explosion.

\textbf{Problem it solves:} You have \texttt{Shape} (Circle, Square) and rendering APIs (Vector, Raster). With inheritance you\textquotesingle d need \texttt{VectorCircle,\ RasterCircle,\ VectorSquare,\ RasterSquare} --- and every new shape \emph{or} new renderer multiplies the matrix. Bridge splits these into two hierarchies: \texttt{Shape} (abstraction) \emph{has-a} \texttt{Renderer} (implementor). Now \texttt{n} shapes + \texttt{m} renderers = \texttt{n\ +\ m} classes, not \texttt{n\ ×\ m}.

\textbf{Structure:}

\setcodelang{}

\begin{verbatim}
   Abstraction ──has-a──> «interface» Implementor
   (Shape)                  (Renderer: drawCircle/drawLine)
      ▲                          ▲
  RefinedAbstraction        ConcreteImplementor
  (Circle, Square)          (VectorRenderer, RasterRenderer)

  Two hierarchies vary INDEPENDENTLY across the "bridge".
\end{verbatim}

\setcodelang{Java}

\begin{Shaded}
\begin{Highlighting}[]
\KeywordTok{interface} \BuiltInTok{Renderer} \OperatorTok{\{} \DataTypeTok{void} \FunctionTok{renderCircle}\OperatorTok{(}\DataTypeTok{float}\NormalTok{ radius}\OperatorTok{);} \OperatorTok{\}}                 \CommentTok{// implementor}
\KeywordTok{class}\NormalTok{ VectorRenderer }\KeywordTok{implements} \BuiltInTok{Renderer} \OperatorTok{\{}
    \KeywordTok{public} \DataTypeTok{void} \FunctionTok{renderCircle}\OperatorTok{(}\DataTypeTok{float}\NormalTok{ r}\OperatorTok{)\{} \BuiltInTok{System}\OperatorTok{.}\FunctionTok{out}\OperatorTok{.}\FunctionTok{println}\OperatorTok{(}\StringTok{"vector circle r="} \OperatorTok{+}\NormalTok{ r}\OperatorTok{);} \OperatorTok{\}}
\OperatorTok{\}}
\KeywordTok{class}\NormalTok{ RasterRenderer }\KeywordTok{implements} \BuiltInTok{Renderer} \OperatorTok{\{}
    \KeywordTok{public} \DataTypeTok{void} \FunctionTok{renderCircle}\OperatorTok{(}\DataTypeTok{float}\NormalTok{ r}\OperatorTok{)\{} \BuiltInTok{System}\OperatorTok{.}\FunctionTok{out}\OperatorTok{.}\FunctionTok{println}\OperatorTok{(}\StringTok{"pixels circle r="} \OperatorTok{+}\NormalTok{ r}\OperatorTok{);} \OperatorTok{\}}
\OperatorTok{\}}

\KeywordTok{abstract} \KeywordTok{class} \BuiltInTok{Shape} \OperatorTok{\{}                                                   \CommentTok{// abstraction}
    \KeywordTok{protected} \DataTypeTok{final} \BuiltInTok{Renderer}\NormalTok{ renderer}\OperatorTok{;}          \CommentTok{// the BRIDGE (composition)}
    \BuiltInTok{Shape}\OperatorTok{(}\BuiltInTok{Renderer}\NormalTok{ r}\OperatorTok{)\{} \KeywordTok{this}\OperatorTok{.}\FunctionTok{renderer} \OperatorTok{=}\NormalTok{ r}\OperatorTok{;} \OperatorTok{\}}
    \KeywordTok{abstract} \DataTypeTok{void} \FunctionTok{draw}\OperatorTok{();}
\OperatorTok{\}}
\KeywordTok{class}\NormalTok{ Circle }\KeywordTok{extends} \BuiltInTok{Shape} \OperatorTok{\{}
    \KeywordTok{private} \DataTypeTok{final} \DataTypeTok{float}\NormalTok{ radius}\OperatorTok{;}
    \FunctionTok{Circle}\OperatorTok{(}\BuiltInTok{Renderer}\NormalTok{ r}\OperatorTok{,} \DataTypeTok{float}\NormalTok{ radius}\OperatorTok{)\{} \KeywordTok{super}\OperatorTok{(}\NormalTok{r}\OperatorTok{);} \KeywordTok{this}\OperatorTok{.}\FunctionTok{radius} \OperatorTok{=}\NormalTok{ radius}\OperatorTok{;} \OperatorTok{\}}
    \DataTypeTok{void} \FunctionTok{draw}\OperatorTok{()\{}\NormalTok{ renderer}\OperatorTok{.}\FunctionTok{renderCircle}\OperatorTok{(}\NormalTok{radius}\OperatorTok{);} \OperatorTok{\}}   \CommentTok{// delegate across the bridge}
\OperatorTok{\}}
\CommentTok{// new shape OR new renderer added independently — no n×m explosion:}
\KeywordTok{new} \FunctionTok{Circle}\OperatorTok{(}\KeywordTok{new} \FunctionTok{VectorRenderer}\OperatorTok{(),} \DecValTok{5}\OperatorTok{).}\FunctionTok{draw}\OperatorTok{();}
\KeywordTok{new} \FunctionTok{Circle}\OperatorTok{(}\KeywordTok{new} \FunctionTok{RasterRenderer}\OperatorTok{(),} \DecValTok{5}\OperatorTok{).}\FunctionTok{draw}\OperatorTok{();}
\end{Highlighting}
\end{Shaded}

\textbf{When to use:} Two (or more) orthogonal dimensions that both vary (shape \(\times\) renderer, message \(\times\) transport, UI \(\times\) theme, driver \(\times\) platform); you want to avoid a combinatorial subclass matrix; you want to swap implementation at runtime. \textbf{When NOT to use:} Only one dimension varies (plain Strategy/inheritance suffices); premature when the second axis is hypothetical (YAGNI).

\textbf{Real framework example:} JDBC (the \texttt{Driver} abstraction bridges to vendor implementations), SLF4J (logging API bridged to log4j/logback backends), device-driver architectures, AWT peer architecture (\texttt{Button} \(\leftrightarrow\) platform-specific peer).

\textbf{Trade-offs:} Independent evolution of two hierarchies, no class explosion, runtime impl swap, hides impl from clients. Upfront complexity (two hierarchies + indirection); over-engineering if the second axis never materializes; can be confused with Strategy (it is structurally similar --- intent: Bridge organizes a \emph{structure}, Strategy swaps an \emph{algorithm}).

\subsection{Flyweight}\label{flyweight}

\textbf{Intent:} Use sharing to support a \textbf{large number of fine-grained objects efficiently} --- separate the \textbf{intrinsic} (shared, immutable) state from the \textbf{extrinsic} (context-specific) state, and reuse one shared object for all instances with the same intrinsic state.

\textbf{Problem it solves:} Rendering a document with a million characters, or a game with millions of trees/particles. Creating a distinct object per character (with font, glyph bitmap, kerning tables) blows up memory. Flyweight stores the \emph{intrinsic} state (the glyph \textquotesingle A\textquotesingle, its font) \textbf{once} and shares it; the \emph{extrinsic} state (position on the page) is passed in by the caller per use.

\textbf{Structure:}

\setcodelang{}

\begin{verbatim}
FlyweightFactory: pool { "A@Arial": glyphA, "B@Arial": glyphB }  ← shared, reused
   getGlyph(char, font) → returns cached flyweight (creates once)

Glyph (flyweight): intrinsic = {char, font, bitmap}  (immutable, shared)
   render(x, y)  ← extrinsic (x,y) PASSED IN, not stored

1,000,000 chars on screen → maybe 80 distinct Glyph objects in memory.
\end{verbatim}

\setcodelang{Python}

\begin{Shaded}
\begin{Highlighting}[]
\KeywordTok{class}\NormalTok{ Glyph:                              }\CommentTok{\# FLYWEIGHT — intrinsic state only}
    \KeywordTok{def} \FunctionTok{\_\_init\_\_}\NormalTok{(}\VariableTok{self}\NormalTok{, char, font):}
        \VariableTok{self}\NormalTok{.char, }\VariableTok{self}\NormalTok{.font }\OperatorTok{=}\NormalTok{ char, font   }\CommentTok{\# shared, immutable}
    \KeywordTok{def}\NormalTok{ render(}\VariableTok{self}\NormalTok{, x, y):              }\CommentTok{\# extrinsic (x,y) passed per call}
        \BuiltInTok{print}\NormalTok{(}\SpecialStringTok{f"draw \textquotesingle{}}\SpecialCharTok{\{}\VariableTok{self}\SpecialCharTok{.}\NormalTok{char}\SpecialCharTok{\}}\SpecialStringTok{\textquotesingle{} (}\SpecialCharTok{\{}\VariableTok{self}\SpecialCharTok{.}\NormalTok{font}\SpecialCharTok{\}}\SpecialStringTok{) at (}\SpecialCharTok{\{}\NormalTok{x}\SpecialCharTok{\}}\SpecialStringTok{,}\SpecialCharTok{\{}\NormalTok{y}\SpecialCharTok{\}}\SpecialStringTok{)"}\NormalTok{)}

\KeywordTok{class}\NormalTok{ GlyphFactory:                       }\CommentTok{\# ensures sharing}
\NormalTok{    \_pool }\OperatorTok{=}\NormalTok{ \{\}}
    \AttributeTok{@classmethod}
    \KeywordTok{def}\NormalTok{ get(cls, char, font):}
\NormalTok{        key }\OperatorTok{=}\NormalTok{ (char, font)}
        \ControlFlowTok{if}\NormalTok{ key }\KeywordTok{not} \KeywordTok{in}\NormalTok{ cls.\_pool:          }\CommentTok{\# create once, reuse forever}
\NormalTok{            cls.\_pool[key] }\OperatorTok{=}\NormalTok{ Glyph(char, font)}
        \ControlFlowTok{return}\NormalTok{ cls.\_pool[key]}

\CommentTok{\# Rendering "AAB" — only TWO Glyph objects exist, reused:}
\ControlFlowTok{for}\NormalTok{ i, ch }\KeywordTok{in} \BuiltInTok{enumerate}\NormalTok{(}\StringTok{"AAB"}\NormalTok{):}
\NormalTok{    GlyphFactory.get(ch, }\StringTok{"Arial"}\NormalTok{).render(i }\OperatorTok{*} \DecValTok{10}\NormalTok{, }\DecValTok{0}\NormalTok{)}
\BuiltInTok{print}\NormalTok{(}\BuiltInTok{len}\NormalTok{(GlyphFactory.\_pool))           }\CommentTok{\# 2, not 3}
\end{Highlighting}
\end{Shaded}

\textbf{When to use:} Huge numbers of similar objects where most state is shareable and memory is the bottleneck (text glyphs, game particles/tiles, map markers, syntax-highlight tokens). \textbf{When NOT to use:} Few objects; little shareable state; when the bookkeeping (factory + intrinsic/extrinsic split) costs more than the memory saved (premature optimization).

\textbf{Real framework example:} Java \texttt{Integer.valueOf()} cache (\(-\)128..127 share instances), \texttt{String} interning / the string pool, \texttt{Boolean.TRUE}/\texttt{FALSE}, Python\textquotesingle s small-int caching, glyph caches in text-rendering engines.

\textbf{Trade-offs:} Massive memory savings for many fine-grained objects. Code complexity (intrinsic/extrinsic split is unintuitive), shared state \textbf{must be immutable / thread-safe} or you get cross-talk bugs, trades CPU (re-passing extrinsic state) for RAM.

\subsection{Structural Patterns --- Comparison}\label{structural-patterns--comparison}

\begin{longtable}[]{@{}
  >{\raggedright\arraybackslash}p{(\columnwidth - 6\tabcolsep) * \real{0.0957}}
  >{\raggedright\arraybackslash}p{(\columnwidth - 6\tabcolsep) * \real{0.3404}}
  >{\raggedright\arraybackslash}p{(\columnwidth - 6\tabcolsep) * \real{0.2340}}
  >{\raggedright\arraybackslash}p{(\columnwidth - 6\tabcolsep) * \real{0.3298}}@{}}
\toprule\noalign{}
\rowcolor{acc!16}
\begin{minipage}[b]{\linewidth}\raggedright
Pattern
\end{minipage} & \begin{minipage}[b]{\linewidth}\raggedright
One-line intent
\end{minipage} & \begin{minipage}[b]{\linewidth}\raggedright
Wraps/structures
\end{minipage} & \begin{minipage}[b]{\linewidth}\raggedright
Key distinction
\end{minipage} \\
\midrule\noalign{}
\endhead
\bottomrule\noalign{}
\endlastfoot
\textbf{Adapter} & Make incompatible interfaces fit & One object & Converts an interface \\
\rowcolor{zebra}
\textbf{Decorator} & Add responsibilities at runtime & One object (stackable) & \emph{Adds} behavior, same interface \\
\textbf{Proxy} & Control access to an object & One object & \emph{Controls} access, same interface \\
\rowcolor{zebra}
\textbf{Facade} & Simplify a complex subsystem & Many objects & \emph{New, simpler} interface \\
\textbf{Composite} & Treat tree of objects uniformly & Tree of objects & Part-whole hierarchy \\
\rowcolor{zebra}
\textbf{Bridge} & Split abstraction from impl & Two hierarchies & Avoid n\(\times\)m class explosion \\
\textbf{Flyweight} & Share to save memory & Pool of shared objects & Intrinsic vs extrinsic state \\
\end{longtable}

\section{Behavioral Patterns}\label{behavioral-patterns}

Behavioral patterns are about \textbf{algorithms and the assignment of responsibilities} between objects --- how they communicate and distribute work. The theme: flexible communication that keeps senders and receivers loosely coupled.

\subsection{Strategy}\label{strategy}

\textbf{Intent:} Define a family of algorithms, encapsulate each one, and make them interchangeable. Strategy lets the algorithm vary independently from the clients that use it.

\textbf{Problem it solves:} A route planner must compute routes by car, by walking, by public transit. Hard-coding \texttt{if\ (mode\ ==\ CAR)\ …\ elif\ (mode\ ==\ WALK)\ …} means every new mode edits the planner, and the planner balloons. Strategy extracts each algorithm into its own class behind a common interface; the planner holds a reference and delegates --- adding a mode is a \emph{new class, zero edits} (OCP).

\textbf{Structure:}

\setcodelang{}

\begin{verbatim}
Context ──has-a──> «interface» Strategy { execute() }
   doWork(){ strategy.execute(); }       ▲ realizes
                          ┌──────────────┼──────────────┐
                    ConcreteA        ConcreteB        ConcreteC
                    execute()        execute()        execute()
\end{verbatim}

\setcodelang{Java}

\begin{Shaded}
\begin{Highlighting}[]
\KeywordTok{interface}\NormalTok{ RouteStrategy }\OperatorTok{\{} \DataTypeTok{int} \FunctionTok{estimate}\OperatorTok{(}\BuiltInTok{String}\NormalTok{ from}\OperatorTok{,} \BuiltInTok{String}\NormalTok{ to}\OperatorTok{);} \OperatorTok{\}}      \CommentTok{// algorithm}

\KeywordTok{class}\NormalTok{ CarStrategy    }\KeywordTok{implements}\NormalTok{ RouteStrategy }\OperatorTok{\{} \KeywordTok{public} \DataTypeTok{int} \FunctionTok{estimate}\OperatorTok{(}\BuiltInTok{String}\NormalTok{ a}\OperatorTok{,}\BuiltInTok{String}\NormalTok{ b}\OperatorTok{)\{} \ControlFlowTok{return} \DecValTok{30}\OperatorTok{;} \OperatorTok{\}} \OperatorTok{\}}
\KeywordTok{class}\NormalTok{ WalkStrategy   }\KeywordTok{implements}\NormalTok{ RouteStrategy }\OperatorTok{\{} \KeywordTok{public} \DataTypeTok{int} \FunctionTok{estimate}\OperatorTok{(}\BuiltInTok{String}\NormalTok{ a}\OperatorTok{,}\BuiltInTok{String}\NormalTok{ b}\OperatorTok{)\{} \ControlFlowTok{return} \DecValTok{90}\OperatorTok{;} \OperatorTok{\}} \OperatorTok{\}}
\KeywordTok{class}\NormalTok{ TransitStrategy }\KeywordTok{implements}\NormalTok{ RouteStrategy}\OperatorTok{\{} \KeywordTok{public} \DataTypeTok{int} \FunctionTok{estimate}\OperatorTok{(}\BuiltInTok{String}\NormalTok{ a}\OperatorTok{,}\BuiltInTok{String}\NormalTok{ b}\OperatorTok{)\{} \ControlFlowTok{return} \DecValTok{45}\OperatorTok{;} \OperatorTok{\}} \OperatorTok{\}}

\KeywordTok{class}\NormalTok{ Navigator }\OperatorTok{\{}                                  \CommentTok{// CONTEXT}
    \KeywordTok{private}\NormalTok{ RouteStrategy strategy}\OperatorTok{;}
    \FunctionTok{Navigator}\OperatorTok{(}\NormalTok{RouteStrategy s}\OperatorTok{)\{} \KeywordTok{this}\OperatorTok{.}\FunctionTok{strategy} \OperatorTok{=}\NormalTok{ s}\OperatorTok{;} \OperatorTok{\}}
    \DataTypeTok{void} \FunctionTok{setStrategy}\OperatorTok{(}\NormalTok{RouteStrategy s}\OperatorTok{)\{} \KeywordTok{this}\OperatorTok{.}\FunctionTok{strategy} \OperatorTok{=}\NormalTok{ s}\OperatorTok{;} \OperatorTok{\}}   \CommentTok{// swap at runtime}
    \DataTypeTok{int} \FunctionTok{route}\OperatorTok{(}\BuiltInTok{String}\NormalTok{ from}\OperatorTok{,} \BuiltInTok{String}\NormalTok{ to}\OperatorTok{)\{} \ControlFlowTok{return}\NormalTok{ strategy}\OperatorTok{.}\FunctionTok{estimate}\OperatorTok{(}\NormalTok{from}\OperatorTok{,}\NormalTok{ to}\OperatorTok{);} \OperatorTok{\}}
\OperatorTok{\}}
\NormalTok{Navigator nav }\OperatorTok{=} \KeywordTok{new} \FunctionTok{Navigator}\OperatorTok{(}\KeywordTok{new} \FunctionTok{CarStrategy}\OperatorTok{());}
\NormalTok{nav}\OperatorTok{.}\FunctionTok{route}\OperatorTok{(}\StringTok{"A"}\OperatorTok{,} \StringTok{"B"}\OperatorTok{);}                  \CommentTok{// 30}
\NormalTok{nav}\OperatorTok{.}\FunctionTok{setStrategy}\OperatorTok{(}\KeywordTok{new} \FunctionTok{TransitStrategy}\OperatorTok{());} \CommentTok{// swapped — Navigator code untouched}
\end{Highlighting}
\end{Shaded}

\textbf{When to use:} Multiple interchangeable algorithms (pricing, compression codecs, payment methods, sorting orders, retry policies); you want to swap them at runtime; to eliminate a growing \texttt{if/switch} over algorithm type. \textbf{When NOT to use:} The algorithm never varies; in functional/dynamic languages a plain first-class function often \emph{is} the strategy (don\textquotesingle t ceremony-wrap it in a class).

\textbf{Real framework example:} \texttt{java.util.Comparator} (passed to \texttt{Collections.sort}), \texttt{java.util.concurrent.RejectedExecutionHandler}, Spring \texttt{Resource} loaders, Python\textquotesingle s \texttt{key=} callables, payment-gateway selection.

\textbf{Trade-offs:} OCP (new algorithm = new class), runtime swapping, testable in isolation, replaces conditionals. More classes/objects; clients must know the strategies exist to choose one; overkill when a lambda suffices.

\begin{quote}
\textbf{Strategy is the poster child for composition + OCP + DIP.} If an interviewer asks "name a pattern embodying Open/Closed," answer Strategy.
\end{quote}

\subsection{Observer}\label{observer}

\textbf{Intent:} Define a one-to-many dependency so that when one object (the \textbf{subject}) changes state, all its dependents (\textbf{observers}) are notified and updated automatically.

\textbf{Problem it solves:} A stock price changes; a chart, an alert, and a logger must all react. Hard-wiring the stock object to call each consumer couples them tightly and means adding a consumer edits the stock. Observer lets consumers \emph{subscribe}; the subject broadcasts to whoever is registered, knowing nothing about who they are.

\textbf{Push vs. Pull:}

\begin{longtable}[]{@{}
  >{\raggedright\arraybackslash}p{(\columnwidth - 4\tabcolsep) * \real{0.0600}}
  >{\raggedright\arraybackslash}p{(\columnwidth - 4\tabcolsep) * \real{0.5319}}
  >{\raggedright\arraybackslash}p{(\columnwidth - 4\tabcolsep) * \real{0.4233}}@{}}
\toprule\noalign{}
\rowcolor{acc!16}
\begin{minipage}[b]{\linewidth}\raggedright
Model
\end{minipage} & \begin{minipage}[b]{\linewidth}\raggedright
Subject sends\ldots{}
\end{minipage} & \begin{minipage}[b]{\linewidth}\raggedright
Trade-off
\end{minipage} \\
\midrule\noalign{}
\endhead
\bottomrule\noalign{}
\endlastfoot
\textbf{Push} & The changed data in the notification (\texttt{update(newPrice)}) & Simple for observers; subject must guess what they need \\
\rowcolor{zebra}
\textbf{Pull} & Just "something changed"; observer queries the subject (\texttt{update(subject)} \(\rightarrow\) \texttt{subject.getPrice()}) & Observers fetch exactly what they need; more calls, looser data coupling \\
\end{longtable}

\textbf{Structure:}

\setcodelang{}

\begin{verbatim}
Subject { subscribe(o) / unsubscribe(o) / notify() → for o in observers: o.update() }
   │ holds list of
   ▼
«interface» Observer { update(...) }
   ▲
ConcreteObserverA   ConcreteObserverB   ...
\end{verbatim}

\setcodelang{Python}

\begin{Shaded}
\begin{Highlighting}[]
\ImportTok{from}\NormalTok{ abc }\ImportTok{import}\NormalTok{ ABC, abstractmethod}

\KeywordTok{class}\NormalTok{ Observer(ABC):}
    \AttributeTok{@abstractmethod}
    \KeywordTok{def}\NormalTok{ update(}\VariableTok{self}\NormalTok{, price): ...                 }\CommentTok{\# PUSH model: data sent in}

\KeywordTok{class}\NormalTok{ Stock:                                      }\CommentTok{\# SUBJECT}
    \KeywordTok{def} \FunctionTok{\_\_init\_\_}\NormalTok{(}\VariableTok{self}\NormalTok{):}
        \VariableTok{self}\NormalTok{.\_observers }\OperatorTok{=}\NormalTok{ []}
        \VariableTok{self}\NormalTok{.\_price }\OperatorTok{=} \DecValTok{0}
    \KeywordTok{def}\NormalTok{ subscribe(}\VariableTok{self}\NormalTok{, o): }\VariableTok{self}\NormalTok{.\_observers.append(o)}
    \KeywordTok{def}\NormalTok{ unsubscribe(}\VariableTok{self}\NormalTok{, o): }\VariableTok{self}\NormalTok{.\_observers.remove(o)   }\CommentTok{\# MUST exist (see leak below)}
    \KeywordTok{def}\NormalTok{ set\_price(}\VariableTok{self}\NormalTok{, p):}
        \VariableTok{self}\NormalTok{.\_price }\OperatorTok{=}\NormalTok{ p}
        \ControlFlowTok{for}\NormalTok{ o }\KeywordTok{in} \VariableTok{self}\NormalTok{.\_observers:                 }\CommentTok{\# broadcast}
\NormalTok{            o.update(p)}

\KeywordTok{class}\NormalTok{ PriceAlert(Observer):}
    \KeywordTok{def}\NormalTok{ update(}\VariableTok{self}\NormalTok{, price): }\BuiltInTok{print}\NormalTok{(}\SpecialStringTok{f"ALERT: }\SpecialCharTok{\{}\NormalTok{price}\SpecialCharTok{\}}\SpecialStringTok{"}\NormalTok{)}
\KeywordTok{class}\NormalTok{ Chart(Observer):}
    \KeywordTok{def}\NormalTok{ update(}\VariableTok{self}\NormalTok{, price): }\BuiltInTok{print}\NormalTok{(}\SpecialStringTok{f"chart += }\SpecialCharTok{\{}\NormalTok{price}\SpecialCharTok{\}}\SpecialStringTok{"}\NormalTok{)}

\NormalTok{s }\OperatorTok{=}\NormalTok{ Stock()}
\NormalTok{s.subscribe(PriceAlert())}\OperatorTok{;}\NormalTok{ s.subscribe(Chart())}
\NormalTok{s.set\_price(}\DecValTok{150}\NormalTok{)   }\CommentTok{\# both notified automatically}
\end{Highlighting}
\end{Shaded}

\textbf{The lapsed-listener leak (must mention in interviews):} the subject holds a strong reference to every observer. If an observer is discarded without calling \texttt{unsubscribe()}, the subject keeps it alive forever --- a \textbf{memory leak}, and it keeps receiving (now-useless) notifications. Fixes: explicit unsubscribe (try/finally), weak references (\texttt{WeakHashMap}/\texttt{WeakReference}), or scoping subscriptions to a lifecycle.

\textbf{When to use:} Event systems, MVC (model notifies views), reactive UIs, pub/sub, real-time dashboards, anywhere "many things must react to one thing changing." \textbf{When NOT to use:} When notification order matters (observers fire in undefined order); when update cascades could storm (debounce/batch); when a simple direct call to one consumer is enough.

\textbf{Real framework example:} Java \texttt{java.util.Observer} (deprecated --- for a reason), Swing/AWT \texttt{addActionListener}, \texttt{PropertyChangeListener}, RxJava \texttt{Observable}, React state/\texttt{useState} re-render, Redux store subscriptions, Django signals, Kafka consumers (distributed Observer).

\textbf{Trade-offs:} Loose coupling producer\(\leftrightarrow\)consumer, dynamic subscription, broadcast for free, OCP. Lapsed-listener leaks, undefined notification order, hard-to-trace control flow ("who fired this update?"), update storms / cascade loops, debugging "spooky action at a distance."

\subsection{Command}\label{command}

\textbf{Intent:} Encapsulate a request as an object, thereby letting you parameterize clients with different requests, queue or log requests, and support \textbf{undo/redo}.

\textbf{Problem it solves:} A text editor\textquotesingle s buttons, menu items, and keyboard shortcuts all trigger "actions." Wiring each UI element directly to editor methods couples them and makes undo impossible. Command turns each action into an object with \texttt{execute()} and \texttt{undo()}. Now you can queue them, log them, replay them (macros), and maintain an undo/redo history stack.

\textbf{Structure:}

\setcodelang{}

\begin{verbatim}
Invoker (button) ──holds──> «interface» Command { execute(); undo() }
                                   ▲
                           ConcreteCommand ──calls──> Receiver (the editor)
                           execute(){ receiver.action(); save state for undo }
History: stack<Command>  →  undo() pops & calls cmd.undo()
\end{verbatim}

\setcodelang{Python}

\begin{Shaded}
\begin{Highlighting}[]
\ImportTok{from}\NormalTok{ abc }\ImportTok{import}\NormalTok{ ABC, abstractmethod}

\KeywordTok{class}\NormalTok{ Command(ABC):}
    \AttributeTok{@abstractmethod}
    \KeywordTok{def}\NormalTok{ execute(}\VariableTok{self}\NormalTok{): ...}
    \AttributeTok{@abstractmethod}
    \KeywordTok{def}\NormalTok{ undo(}\VariableTok{self}\NormalTok{): ...}

\KeywordTok{class}\NormalTok{ Document:                                   }\CommentTok{\# RECEIVER}
    \KeywordTok{def} \FunctionTok{\_\_init\_\_}\NormalTok{(}\VariableTok{self}\NormalTok{): }\VariableTok{self}\NormalTok{.text }\OperatorTok{=} \StringTok{""}
    \KeywordTok{def}\NormalTok{ insert(}\VariableTok{self}\NormalTok{, s): }\VariableTok{self}\NormalTok{.text }\OperatorTok{+=}\NormalTok{ s}
    \KeywordTok{def}\NormalTok{ delete(}\VariableTok{self}\NormalTok{, n): }\VariableTok{self}\NormalTok{.text }\OperatorTok{=} \VariableTok{self}\NormalTok{.text[:}\OperatorTok{{-}}\NormalTok{n]}

\KeywordTok{class}\NormalTok{ InsertCommand(Command):                     }\CommentTok{\# CONCRETE COMMAND}
    \KeywordTok{def} \FunctionTok{\_\_init\_\_}\NormalTok{(}\VariableTok{self}\NormalTok{, doc, s): }\VariableTok{self}\NormalTok{.doc, }\VariableTok{self}\NormalTok{.s }\OperatorTok{=}\NormalTok{ doc, s}
    \KeywordTok{def}\NormalTok{ execute(}\VariableTok{self}\NormalTok{): }\VariableTok{self}\NormalTok{.doc.insert(}\VariableTok{self}\NormalTok{.s)}
    \KeywordTok{def}\NormalTok{ undo(}\VariableTok{self}\NormalTok{):    }\VariableTok{self}\NormalTok{.doc.delete(}\BuiltInTok{len}\NormalTok{(}\VariableTok{self}\NormalTok{.s))   }\CommentTok{\# inverse op}

\KeywordTok{class}\NormalTok{ History:                                    }\CommentTok{\# INVOKER + undo/redo stacks}
    \KeywordTok{def} \FunctionTok{\_\_init\_\_}\NormalTok{(}\VariableTok{self}\NormalTok{): }\VariableTok{self}\NormalTok{.\_done, }\VariableTok{self}\NormalTok{.\_undone }\OperatorTok{=}\NormalTok{ [], []}
    \KeywordTok{def}\NormalTok{ run(}\VariableTok{self}\NormalTok{, cmd):}
\NormalTok{        cmd.execute()}\OperatorTok{;} \VariableTok{self}\NormalTok{.\_done.append(cmd)}\OperatorTok{;} \VariableTok{self}\NormalTok{.\_undone.clear()}
    \KeywordTok{def}\NormalTok{ undo(}\VariableTok{self}\NormalTok{):}
        \ControlFlowTok{if} \VariableTok{self}\NormalTok{.\_done:}
\NormalTok{            cmd }\OperatorTok{=} \VariableTok{self}\NormalTok{.\_done.pop()}\OperatorTok{;}\NormalTok{ cmd.undo()}\OperatorTok{;} \VariableTok{self}\NormalTok{.\_undone.append(cmd)}
    \KeywordTok{def}\NormalTok{ redo(}\VariableTok{self}\NormalTok{):}
        \ControlFlowTok{if} \VariableTok{self}\NormalTok{.\_undone:}
\NormalTok{            cmd }\OperatorTok{=} \VariableTok{self}\NormalTok{.\_undone.pop()}\OperatorTok{;}\NormalTok{ cmd.execute()}\OperatorTok{;} \VariableTok{self}\NormalTok{.\_done.append(cmd)}

\NormalTok{doc }\OperatorTok{=}\NormalTok{ Document()}\OperatorTok{;}\NormalTok{ h }\OperatorTok{=}\NormalTok{ History()}
\NormalTok{h.run(InsertCommand(doc, }\StringTok{"hello"}\NormalTok{))   }\CommentTok{\# text="hello"}
\NormalTok{h.run(InsertCommand(doc, }\StringTok{" world"}\NormalTok{))  }\CommentTok{\# text="hello world"}
\NormalTok{h.undo()                             }\CommentTok{\# text="hello"}
\NormalTok{h.redo()                             }\CommentTok{\# text="hello world"}
\end{Highlighting}
\end{Shaded}

\textbf{When to use:} Undo/redo, transactional operations, job/task queues (serialize a command and run it later/elsewhere), macro recording, GUI actions decoupled from handlers, request logging/replay. \textbf{When NOT to use:} Simple one-shot operations with no undo/queue/log requirement --- a direct method call or a lambda is enough.

\textbf{Real framework example:} \texttt{java.lang.Runnable} / \texttt{Callable} (a command for a thread pool), Swing \texttt{Action}, \texttt{javax.swing.undo.UndoableEdit}, job queues (Celery tasks, Sidekiq jobs), the Command pattern underlies CQRS commands and event sourcing.

\textbf{Trade-offs:} Decouples invoker from receiver, enables undo/redo/queue/log/macro, SRP (one command = one action), composability (macro = list of commands). A class per action (proliferation), undo logic can be tricky (must capture inverse state --- Memento often pairs here), memory for long histories.

\subsection{State}\label{state}

\textbf{Intent:} Allow an object to alter its behavior when its internal state changes --- the object appears to change its class.

\textbf{Problem it solves:} An order moves through PENDING \(\rightarrow\) PAID \(\rightarrow\) SHIPPED \(\rightarrow\) DELIVERED, and each method (\texttt{pay()}, \texttt{ship()}, \texttt{cancel()}) behaves differently per state (\texttt{ship()} is illegal before payment). Encoding this as \texttt{if\ (status\ ==\ PENDING)\ …\ elif\ (status\ ==\ PAID)\ …} in \emph{every} method produces parallel sprawling conditionals that all must be edited together. State makes each status a class that implements the behavior \emph{and} decides the next transition.

\textbf{Structure:}

\setcodelang{}

\begin{verbatim}
Context (Order) ──has-a──> «interface» State { pay(ctx); ship(ctx) }
   pay(){ state.pay(this); }                  ▲
                              ┌────────────────┼────────────────┐
                         PendingState       PaidState        ShippedState
                         pay()→PaidState     ship()→Shipped   (terminal-ish)
\end{verbatim}

\setcodelang{Python}

\begin{Shaded}
\begin{Highlighting}[]
\ImportTok{from}\NormalTok{ abc }\ImportTok{import}\NormalTok{ ABC, abstractmethod}

\KeywordTok{class}\NormalTok{ OrderState(ABC):}
    \AttributeTok{@abstractmethod}
    \KeywordTok{def}\NormalTok{ pay(}\VariableTok{self}\NormalTok{, order): ...}
    \AttributeTok{@abstractmethod}
    \KeywordTok{def}\NormalTok{ ship(}\VariableTok{self}\NormalTok{, order): ...}

\KeywordTok{class}\NormalTok{ Pending(OrderState):}
    \KeywordTok{def}\NormalTok{ pay(}\VariableTok{self}\NormalTok{, order):}
        \BuiltInTok{print}\NormalTok{(}\StringTok{"payment ok"}\NormalTok{)}\OperatorTok{;}\NormalTok{ order.state }\OperatorTok{=}\NormalTok{ Paid()     }\CommentTok{\# STATE decides the transition}
    \KeywordTok{def}\NormalTok{ ship(}\VariableTok{self}\NormalTok{, order): }\BuiltInTok{print}\NormalTok{(}\StringTok{"can\textquotesingle{}t ship — unpaid"}\NormalTok{)}

\KeywordTok{class}\NormalTok{ Paid(OrderState):}
    \KeywordTok{def}\NormalTok{ pay(}\VariableTok{self}\NormalTok{, order): }\BuiltInTok{print}\NormalTok{(}\StringTok{"already paid"}\NormalTok{)}
    \KeywordTok{def}\NormalTok{ ship(}\VariableTok{self}\NormalTok{, order):}
        \BuiltInTok{print}\NormalTok{(}\StringTok{"shipping"}\NormalTok{)}\OperatorTok{;}\NormalTok{ order.state }\OperatorTok{=}\NormalTok{ Shipped()}

\KeywordTok{class}\NormalTok{ Shipped(OrderState):}
    \KeywordTok{def}\NormalTok{ pay(}\VariableTok{self}\NormalTok{, order): }\BuiltInTok{print}\NormalTok{(}\StringTok{"already paid"}\NormalTok{)}
    \KeywordTok{def}\NormalTok{ ship(}\VariableTok{self}\NormalTok{, order): }\BuiltInTok{print}\NormalTok{(}\StringTok{"already shipped"}\NormalTok{)}

\KeywordTok{class}\NormalTok{ Order:                                }\CommentTok{\# CONTEXT}
    \KeywordTok{def} \FunctionTok{\_\_init\_\_}\NormalTok{(}\VariableTok{self}\NormalTok{): }\VariableTok{self}\NormalTok{.state }\OperatorTok{=}\NormalTok{ Pending()}
    \KeywordTok{def}\NormalTok{ pay(}\VariableTok{self}\NormalTok{):  }\VariableTok{self}\NormalTok{.state.pay(}\VariableTok{self}\NormalTok{)    }\CommentTok{\# delegate to current state}
    \KeywordTok{def}\NormalTok{ ship(}\VariableTok{self}\NormalTok{): }\VariableTok{self}\NormalTok{.state.ship(}\VariableTok{self}\NormalTok{)}

\NormalTok{o }\OperatorTok{=}\NormalTok{ Order()}
\NormalTok{o.ship()  }\CommentTok{\# "can\textquotesingle{}t ship — unpaid"}
\NormalTok{o.pay()   }\CommentTok{\# "payment ok"   (now Paid)}
\NormalTok{o.ship()  }\CommentTok{\# "shipping"     (now Shipped)}
\end{Highlighting}
\end{Shaded}

\textbf{State vs. Strategy (the classic interview pair):} structurally identical (a context delegates to a swappable helper interface). The \textbf{intent} differs:

\begin{longtable}[]{@{}
  >{\raggedright\arraybackslash}p{(\columnwidth - 4\tabcolsep) * \real{0.2689}}
  >{\raggedright\arraybackslash}p{(\columnwidth - 4\tabcolsep) * \real{0.3060}}
  >{\raggedright\arraybackslash}p{(\columnwidth - 4\tabcolsep) * \real{0.4251}}@{}}
\toprule\noalign{}
\rowcolor{acc!16}
\begin{minipage}[b]{\linewidth}\raggedright
\end{minipage} & \begin{minipage}[b]{\linewidth}\raggedright
Strategy
\end{minipage} & \begin{minipage}[b]{\linewidth}\raggedright
State
\end{minipage} \\
\midrule\noalign{}
\endhead
\bottomrule\noalign{}
\endlastfoot
Who picks the impl? & The \textbf{client} chooses an algorithm & The object \textbf{transitions itself} between states \\
\rowcolor{zebra}
Do the impls know each other? & No --- independent algorithms & Yes --- a state knows its successor states \\
How often does it change? & Usually set once, rarely mid-flow & Changes repeatedly as the object\textquotesingle s lifecycle advances \\
\rowcolor{zebra}
Mental model & "How do I do X?" & "What can I do \emph{now}?" \\
\end{longtable}

\textbf{When to use:} Lifecycle-heavy objects (orders, tickets, connections, documents, vending machines, TCP connections, media players); many state-dependent methods; replacing a multi-method \texttt{switch(status)}. \textbf{When NOT to use:} 2--3 states that rarely change and have trivial behavior --- an enum + if/else is clearer (don\textquotesingle t over-engineer).

\textbf{Real framework example:} TCP connection states, \texttt{java.lang.Thread} states, workflow engines, game character state machines, parser states, Spring State Machine.

\textbf{Trade-offs:} Eliminates sprawling conditionals, each state\textquotesingle s behavior + transitions localized (SRP/OCP), explicit legal transitions. Class per state (proliferation for many states), state-transition logic can be scattered across states (consider a transition table), overkill for trivial state sets.

\subsection{Template Method}\label{template-method}

\textbf{Intent:} Define the \textbf{skeleton} of an algorithm in a base method, deferring some steps to subclasses. Subclasses redefine steps without changing the algorithm\textquotesingle s structure.

\textbf{Problem it solves:} Exporting data to CSV vs. JSON shares a structure --- fetch data, format it, write to file, notify --- but the \emph{format} and \emph{write} steps differ. Duplicating the whole flow in each exporter violates DRY and risks the flows diverging. Template Method puts the invariant flow in a base method and calls overridable hooks for the varying steps (Hollywood Principle --- the base calls down into your steps).

\textbf{Structure:}

\setcodelang{}

\begin{verbatim}
«abstract» AbstractClass
  templateMethod() {           ← FINAL skeleton, defines the order
      step1();                 ← concrete (shared)
      step2();                 ← abstract  (subclass fills in)
      hook();                  ← optional override (default no-op)
  }
       ▲ overrides step2()
  ┌────┴─────┐
ConcreteA   ConcreteB
\end{verbatim}

\setcodelang{Java}

\begin{Shaded}
\begin{Highlighting}[]
\KeywordTok{abstract} \KeywordTok{class}\NormalTok{ DataExporter }\OperatorTok{\{}
    \KeywordTok{public} \DataTypeTok{final} \DataTypeTok{void} \FunctionTok{export}\OperatorTok{(}\BuiltInTok{String}\NormalTok{ file}\OperatorTok{)} \OperatorTok{\{}     \CommentTok{// TEMPLATE METHOD — final = fixed flow}
        \DataTypeTok{var}\NormalTok{ data }\OperatorTok{=} \FunctionTok{fetch}\OperatorTok{();}      \CommentTok{// shared step}
        \DataTypeTok{var}\NormalTok{ out  }\OperatorTok{=} \FunctionTok{format}\OperatorTok{(}\NormalTok{data}\OperatorTok{);} \CommentTok{// varies  → abstract}
        \FunctionTok{write}\OperatorTok{(}\NormalTok{out}\OperatorTok{,}\NormalTok{ file}\OperatorTok{);}        \CommentTok{// varies  → abstract}
        \FunctionTok{onComplete}\OperatorTok{();}            \CommentTok{// hook    → optional override}
    \OperatorTok{\}}
    \KeywordTok{protected} \BuiltInTok{List}\OperatorTok{\textless{}}\BuiltInTok{String}\OperatorTok{\textgreater{}} \FunctionTok{fetch}\OperatorTok{()} \OperatorTok{\{} \ControlFlowTok{return} \BuiltInTok{List}\OperatorTok{.}\FunctionTok{of}\OperatorTok{(}\StringTok{"a"}\OperatorTok{,} \StringTok{"b"}\OperatorTok{);} \OperatorTok{\}}   \CommentTok{// shared}
    \KeywordTok{protected} \KeywordTok{abstract} \BuiltInTok{String} \FunctionTok{format}\OperatorTok{(}\BuiltInTok{List}\OperatorTok{\textless{}}\BuiltInTok{String}\OperatorTok{\textgreater{}}\NormalTok{ data}\OperatorTok{);}          \CommentTok{// subclass step}
    \KeywordTok{protected} \KeywordTok{abstract} \DataTypeTok{void}   \FunctionTok{write}\OperatorTok{(}\BuiltInTok{String}\NormalTok{ content}\OperatorTok{,} \BuiltInTok{String}\NormalTok{ file}\OperatorTok{);} \CommentTok{// subclass step}
    \KeywordTok{protected} \DataTypeTok{void} \FunctionTok{onComplete}\OperatorTok{()} \OperatorTok{\{\}}                                 \CommentTok{// hook (default no{-}op)}
\OperatorTok{\}}
\KeywordTok{class}\NormalTok{ CsvExporter }\KeywordTok{extends}\NormalTok{ DataExporter }\OperatorTok{\{}
    \KeywordTok{protected} \BuiltInTok{String} \FunctionTok{format}\OperatorTok{(}\BuiltInTok{List}\OperatorTok{\textless{}}\BuiltInTok{String}\OperatorTok{\textgreater{}}\NormalTok{ d}\OperatorTok{)\{} \ControlFlowTok{return} \BuiltInTok{String}\OperatorTok{.}\FunctionTok{join}\OperatorTok{(}\StringTok{","}\OperatorTok{,}\NormalTok{ d}\OperatorTok{);} \OperatorTok{\}}
    \KeywordTok{protected} \DataTypeTok{void} \FunctionTok{write}\OperatorTok{(}\BuiltInTok{String}\NormalTok{ c}\OperatorTok{,} \BuiltInTok{String}\NormalTok{ f}\OperatorTok{)\{} \CommentTok{/* write CSV */} \OperatorTok{\}}
\OperatorTok{\}}
\KeywordTok{class}\NormalTok{ JsonExporter }\KeywordTok{extends}\NormalTok{ DataExporter }\OperatorTok{\{}
    \KeywordTok{protected} \BuiltInTok{String} \FunctionTok{format}\OperatorTok{(}\BuiltInTok{List}\OperatorTok{\textless{}}\BuiltInTok{String}\OperatorTok{\textgreater{}}\NormalTok{ d}\OperatorTok{)\{} \ControlFlowTok{return}\NormalTok{ d}\OperatorTok{.}\FunctionTok{toString}\OperatorTok{();} \OperatorTok{\}}
    \KeywordTok{protected} \DataTypeTok{void} \FunctionTok{write}\OperatorTok{(}\BuiltInTok{String}\NormalTok{ c}\OperatorTok{,} \BuiltInTok{String}\NormalTok{ f}\OperatorTok{)\{} \CommentTok{/* write JSON */} \OperatorTok{\}}
    \KeywordTok{protected} \DataTypeTok{void} \FunctionTok{onComplete}\OperatorTok{()\{} \BuiltInTok{System}\OperatorTok{.}\FunctionTok{out}\OperatorTok{.}\FunctionTok{println}\OperatorTok{(}\StringTok{"json done"}\OperatorTok{);} \OperatorTok{\}}  \CommentTok{// overrides hook}
\OperatorTok{\}}
\end{Highlighting}
\end{Shaded}

\textbf{Template Method vs. Strategy:} Template Method varies steps via \textbf{inheritance} (compile-time, subclass overrides) and keeps the algorithm skeleton in the base class; Strategy varies the \textbf{whole algorithm} via \textbf{composition} (runtime swap). Template Method = "fixed flow, pluggable steps"; Strategy = "pluggable whole."

\textbf{When to use:} Several variants share an algorithm skeleton with a few differing steps; you want to enforce the invariant ordering while allowing customization (frameworks do this heavily). \textbf{When NOT to use:} When composition/Strategy gives more flexibility (runtime swap, no inheritance coupling); when the "skeleton" is trivial.

\textbf{Real framework example:} \texttt{java.util.AbstractList}/\texttt{AbstractMap} (you implement \texttt{get}/\texttt{size}, get the rest free), \texttt{javax.servlet.http.HttpServlet.service()} (dispatches to your \texttt{doGet}/\texttt{doPost}), JUnit \texttt{setUp()}/\texttt{test}/\texttt{tearDown()}, Spring\textquotesingle s \texttt{JdbcTemplate.execute()}, Android \texttt{Activity} lifecycle callbacks.

\textbf{Trade-offs:} DRY (shared flow once), enforces invariant structure, hooks for optional customization, framework-friendly. Inheritance coupling (fragile base class), inverted control can be hard to follow, hard to change the skeleton later, only one inheritance slot (Java single inheritance).

\subsection{Iterator}\label{iterator}

\textbf{Intent:} Provide a way to access the elements of an aggregate object \textbf{sequentially} without exposing its underlying representation.

\textbf{Problem it solves:} A collection might be backed by an array, a linked list, a tree, or a remote paginated API. Clients shouldn\textquotesingle t need to know --- they just want to traverse it. Iterator exposes a uniform \texttt{hasNext()/next()} (or \texttt{\_\_iter\_\_/\_\_next\_\_}) cursor, decoupling traversal from storage and allowing multiple independent traversals at once.

\textbf{Structure:}

\setcodelang{Python}

\begin{Shaded}
\begin{Highlighting}[]
\NormalTok{Aggregate \{ createIterator(): Iterator \}}
\NormalTok{Iterator  \{ hasNext(): }\BuiltInTok{bool}\OperatorTok{;}  \BuiltInTok{next}\NormalTok{(): T \}}
\NormalTok{   ConcreteIterator holds a cursor into ConcreteAggregate}\StringTok{\textquotesingle{}s internals.}
\end{Highlighting}
\end{Shaded}

\setcodelang{Python}

\begin{Shaded}
\begin{Highlighting}[]
\KeywordTok{class}\NormalTok{ BinaryTree:}
    \KeywordTok{def} \FunctionTok{\_\_init\_\_}\NormalTok{(}\VariableTok{self}\NormalTok{, value, left}\OperatorTok{=}\VariableTok{None}\NormalTok{, right}\OperatorTok{=}\VariableTok{None}\NormalTok{):}
        \VariableTok{self}\NormalTok{.value, }\VariableTok{self}\NormalTok{.left, }\VariableTok{self}\NormalTok{.right }\OperatorTok{=}\NormalTok{ value, left, right}
    \KeywordTok{def} \FunctionTok{\_\_iter\_\_}\NormalTok{(}\VariableTok{self}\NormalTok{):                    }\CommentTok{\# in{-}order traversal, hides tree structure}
        \ControlFlowTok{if} \VariableTok{self}\NormalTok{.left:  }\ControlFlowTok{yield} \ControlFlowTok{from} \VariableTok{self}\NormalTok{.left}
        \ControlFlowTok{yield} \VariableTok{self}\NormalTok{.value}
        \ControlFlowTok{if} \VariableTok{self}\NormalTok{.right: }\ControlFlowTok{yield} \ControlFlowTok{from} \VariableTok{self}\NormalTok{.right}

\NormalTok{tree }\OperatorTok{=}\NormalTok{ BinaryTree(}\DecValTok{2}\NormalTok{, BinaryTree(}\DecValTok{1}\NormalTok{), BinaryTree(}\DecValTok{3}\NormalTok{))}
\ControlFlowTok{for}\NormalTok{ v }\KeywordTok{in}\NormalTok{ tree:        }\CommentTok{\# client traverses uniformly, knows nothing about nodes}
    \BuiltInTok{print}\NormalTok{(v)          }\CommentTok{\# 1, 2, 3}
\end{Highlighting}
\end{Shaded}

\setcodelang{Java}

\begin{Shaded}
\begin{Highlighting}[]
\CommentTok{// Java: implement Iterable\textless{}T\textgreater{} → works with for{-}each and streams}
\KeywordTok{class}\NormalTok{ Bag}\OperatorTok{\textless{}}\NormalTok{T}\OperatorTok{\textgreater{}} \KeywordTok{implements} \BuiltInTok{Iterable}\OperatorTok{\textless{}}\NormalTok{T}\OperatorTok{\textgreater{}} \OperatorTok{\{}
    \KeywordTok{private} \DataTypeTok{final} \BuiltInTok{List}\OperatorTok{\textless{}}\NormalTok{T}\OperatorTok{\textgreater{}}\NormalTok{ items }\OperatorTok{=} \KeywordTok{new} \BuiltInTok{ArrayList}\OperatorTok{\textless{}\textgreater{}();}
    \KeywordTok{public} \DataTypeTok{void} \FunctionTok{add}\OperatorTok{(}\NormalTok{T t}\OperatorTok{)\{}\NormalTok{ items}\OperatorTok{.}\FunctionTok{add}\OperatorTok{(}\NormalTok{t}\OperatorTok{);} \OperatorTok{\}}
    \KeywordTok{public} \BuiltInTok{Iterator}\OperatorTok{\textless{}}\NormalTok{T}\OperatorTok{\textgreater{}} \FunctionTok{iterator}\OperatorTok{()\{} \ControlFlowTok{return}\NormalTok{ items}\OperatorTok{.}\FunctionTok{iterator}\OperatorTok{();} \OperatorTok{\}}  \CommentTok{// delegate or custom}
\OperatorTok{\}}
\end{Highlighting}
\end{Shaded}

\textbf{When to use:} Any custom collection you want traversable uniformly (for-each, streams); when you want multiple simultaneous traversals; lazy/streaming traversal (generators, paginated APIs). \textbf{When NOT to use:} Built-in collections already provide it; trivial fixed structures where direct access is clearer.

\textbf{Real framework example:} \texttt{java.util.Iterator} / \texttt{Iterable} (for-each loop), Python\textquotesingle s iterator protocol (\texttt{\_\_iter\_\_}/\texttt{\_\_next\_\_}, generators), C\# \texttt{IEnumerable}, database cursors, paginated REST clients exposing an iterator.

\textbf{Trade-offs:} Hides internal structure, uniform traversal, multiple concurrent iterators, supports lazy/infinite sequences. Overkill for simple arrays; \textbf{fail-fast vs. fail-safe} concerns (modifying a collection mid-iteration \(\rightarrow\) \texttt{ConcurrentModificationException} in Java, or stale snapshots); stateful iterators complicate concurrency.

\subsection{Chain of Responsibility}\label{chain-of-responsibility}

\textbf{Intent:} Avoid coupling the sender of a request to its receiver by giving \textbf{more than one object a chance to handle} the request. Chain the receivers and pass the request along until one handles it.

\textbf{Problem it solves:} An incoming HTTP request must pass through authentication, rate-limiting, validation, and logging. Hard-coding \texttt{if\ (!auth)\ …\ if\ (!rateOk)\ …\ if\ (!valid)\ …} in one place is rigid and unordered. Chain of Responsibility makes each concern a handler linked to the next; each decides to handle and/or forward. You can reorder, insert, or remove stages without touching the others.

\textbf{Structure:}

\setcodelang{}

\begin{verbatim}
Request → HandlerA → HandlerB → HandlerC → (end)
            │ each: handle it, and/or pass to next
Handler { setNext(h); handle(req){ ...; if (next) next.handle(req); } }
\end{verbatim}

\setcodelang{Python}

\begin{Shaded}
\begin{Highlighting}[]
\ImportTok{from}\NormalTok{ abc }\ImportTok{import}\NormalTok{ ABC, abstractmethod}

\KeywordTok{class}\NormalTok{ Handler(ABC):}
    \KeywordTok{def} \FunctionTok{\_\_init\_\_}\NormalTok{(}\VariableTok{self}\NormalTok{): }\VariableTok{self}\NormalTok{.\_next }\OperatorTok{=} \VariableTok{None}
    \KeywordTok{def}\NormalTok{ set\_next(}\VariableTok{self}\NormalTok{, h):}
        \VariableTok{self}\NormalTok{.\_next }\OperatorTok{=}\NormalTok{ h}\OperatorTok{;} \ControlFlowTok{return}\NormalTok{ h            }\CommentTok{\# enables fluent chaining}
    \AttributeTok{@abstractmethod}
    \KeywordTok{def}\NormalTok{ handle(}\VariableTok{self}\NormalTok{, req): ...}
    \KeywordTok{def}\NormalTok{ \_forward(}\VariableTok{self}\NormalTok{, req):}
        \ControlFlowTok{if} \VariableTok{self}\NormalTok{.\_next: }\VariableTok{self}\NormalTok{.\_next.handle(req)}

\KeywordTok{class}\NormalTok{ AuthHandler(Handler):}
    \KeywordTok{def}\NormalTok{ handle(}\VariableTok{self}\NormalTok{, req):}
        \ControlFlowTok{if} \KeywordTok{not}\NormalTok{ req.get(}\StringTok{"user"}\NormalTok{): }\ControlFlowTok{return} \BuiltInTok{print}\NormalTok{(}\StringTok{"REJECT: no auth"}\NormalTok{)}
        \BuiltInTok{print}\NormalTok{(}\StringTok{"auth ok"}\NormalTok{)}\OperatorTok{;} \VariableTok{self}\NormalTok{.\_forward(req)}

\KeywordTok{class}\NormalTok{ RateLimitHandler(Handler):}
    \KeywordTok{def}\NormalTok{ handle(}\VariableTok{self}\NormalTok{, req):}
        \ControlFlowTok{if}\NormalTok{ req.get(}\StringTok{"count"}\NormalTok{, }\DecValTok{0}\NormalTok{) }\OperatorTok{\textgreater{}} \DecValTok{100}\NormalTok{: }\ControlFlowTok{return} \BuiltInTok{print}\NormalTok{(}\StringTok{"REJECT: rate limited"}\NormalTok{)}
        \BuiltInTok{print}\NormalTok{(}\StringTok{"rate ok"}\NormalTok{)}\OperatorTok{;} \VariableTok{self}\NormalTok{.\_forward(req)}

\KeywordTok{class}\NormalTok{ ValidateHandler(Handler):}
    \KeywordTok{def}\NormalTok{ handle(}\VariableTok{self}\NormalTok{, req):}
        \ControlFlowTok{if} \KeywordTok{not}\NormalTok{ req.get(}\StringTok{"body"}\NormalTok{): }\ControlFlowTok{return} \BuiltInTok{print}\NormalTok{(}\StringTok{"REJECT: empty body"}\NormalTok{)}
        \BuiltInTok{print}\NormalTok{(}\StringTok{"valid — processing!"}\NormalTok{)}

\NormalTok{auth }\OperatorTok{=}\NormalTok{ AuthHandler()}
\NormalTok{auth.set\_next(RateLimitHandler()).set\_next(ValidateHandler())   }\CommentTok{\# build chain}
\NormalTok{auth.handle(\{}\StringTok{"user"}\NormalTok{: }\StringTok{"alice"}\NormalTok{, }\StringTok{"count"}\NormalTok{: }\DecValTok{5}\NormalTok{, }\StringTok{"body"}\NormalTok{: }\StringTok{"data"}\NormalTok{\})      }\CommentTok{\# passes all stages}
\end{Highlighting}
\end{Shaded}

\textbf{When to use:} Request pipelines / middleware (web frameworks, interceptors), event bubbling, multi-level approval flows (expense \(\rightarrow\) manager \(\rightarrow\) director \(\rightarrow\) VP), exception handling hierarchies. \textbf{When NOT to use:} When exactly one known handler exists (just call it); when a request \emph{must} be handled (a chain can fall off the end unhandled --- guard for that).

\textbf{Real framework example:} Servlet \texttt{Filter} chains, Spring Security filter chain, Java logging handler levels, Netty pipeline, Express/Django middleware, AWS API Gateway request flow.

\textbf{Trade-offs:} Decouples sender from receivers, reorder/add/remove handlers freely (OCP/SRP), each handler is focused. A request can go unhandled (silent failure), harder to debug ("which handler swallowed it?"), no guarantee of handling, runtime configuration can obscure flow.

\subsection{Mediator}\label{mediator}

\textbf{Intent:} Define an object that encapsulates how a set of objects interact. Mediator promotes loose coupling by keeping objects from referring to each other explicitly, centralizing their many-to-many communication.

\textbf{Problem it solves:} In a dialog, toggling a checkbox enables a textbox, which enables a button\ldots{} If each widget references every other widget directly, you get an n\(^2\) tangle where every widget knows every other. A Mediator becomes the single hub: widgets notify the mediator, and the mediator orchestrates everyone else. The n\(^2\) mesh becomes a star.

\textbf{Structure:}

\setcodelang{}

\begin{verbatim}
   Before (n²):           After (Mediator):
   A ↔ B ↔ C              A → Mediator ← B
   ↕ ╳ ↕                       ↑
   D ↔ E                       C, D, E
   (everyone knows         (everyone knows only the Mediator)
    everyone)
\end{verbatim}

\setcodelang{Python}

\begin{Shaded}
\begin{Highlighting}[]
\KeywordTok{class}\NormalTok{ DialogMediator:}
    \KeywordTok{def} \FunctionTok{\_\_init\_\_}\NormalTok{(}\VariableTok{self}\NormalTok{):}
        \VariableTok{self}\NormalTok{.checkbox }\OperatorTok{=} \VariableTok{None}\OperatorTok{;} \VariableTok{self}\NormalTok{.textbox }\OperatorTok{=} \VariableTok{None}\OperatorTok{;} \VariableTok{self}\NormalTok{.button }\OperatorTok{=} \VariableTok{None}
    \KeywordTok{def}\NormalTok{ changed(}\VariableTok{self}\NormalTok{, sender):                 }\CommentTok{\# central orchestration}
        \ControlFlowTok{if}\NormalTok{ sender }\KeywordTok{is} \VariableTok{self}\NormalTok{.checkbox:}
            \VariableTok{self}\NormalTok{.textbox.enabled }\OperatorTok{=} \VariableTok{self}\NormalTok{.checkbox.checked}
            \VariableTok{self}\NormalTok{.button.enabled  }\OperatorTok{=} \VariableTok{self}\NormalTok{.checkbox.checked}

\KeywordTok{class}\NormalTok{ Widget:}
    \KeywordTok{def} \FunctionTok{\_\_init\_\_}\NormalTok{(}\VariableTok{self}\NormalTok{, mediator): }\VariableTok{self}\NormalTok{.mediator }\OperatorTok{=}\NormalTok{ mediator}\OperatorTok{;} \VariableTok{self}\NormalTok{.enabled }\OperatorTok{=} \VariableTok{False}

\KeywordTok{class}\NormalTok{ Checkbox(Widget):}
    \KeywordTok{def} \FunctionTok{\_\_init\_\_}\NormalTok{(}\VariableTok{self}\NormalTok{, m): }\BuiltInTok{super}\NormalTok{().}\FunctionTok{\_\_init\_\_}\NormalTok{(m)}\OperatorTok{;} \VariableTok{self}\NormalTok{.checked }\OperatorTok{=} \VariableTok{False}
    \KeywordTok{def}\NormalTok{ toggle(}\VariableTok{self}\NormalTok{):}
        \VariableTok{self}\NormalTok{.checked }\OperatorTok{=} \KeywordTok{not} \VariableTok{self}\NormalTok{.checked}
        \VariableTok{self}\NormalTok{.mediator.changed(}\VariableTok{self}\NormalTok{)            }\CommentTok{\# tell the hub, not the peers}

\NormalTok{m }\OperatorTok{=}\NormalTok{ DialogMediator()}
\NormalTok{m.checkbox }\OperatorTok{=}\NormalTok{ Checkbox(m)}\OperatorTok{;}\NormalTok{ m.textbox }\OperatorTok{=}\NormalTok{ Widget(m)}\OperatorTok{;}\NormalTok{ m.button }\OperatorTok{=}\NormalTok{ Widget(m)}
\NormalTok{m.checkbox.toggle()      }\CommentTok{\# mediator enables textbox + button}
\end{Highlighting}
\end{Shaded}

\textbf{When to use:} Complex many-to-many interactions among components you want decoupled (UI dialogs, air-traffic control, chat rooms routing messages between users). \textbf{When NOT to use:} Simple interactions (the mediator adds a needless hub); risk the mediator becomes a \textbf{God object} absorbing all logic.

\textbf{Real framework example:} \texttt{java.util.concurrent.ExecutorService} (mediates tasks\(\leftrightarrow\)threads), MVC controllers, message brokers, JMS, chat-server routing, Spring\textquotesingle s event publisher as a lightweight mediator.

\textbf{Trade-offs:} Reduces coupling (n\(^2\) \(\rightarrow\) n), centralizes interaction logic, components reusable in isolation. The mediator can grow into a God object / single point of complexity; centralization can become a bottleneck.

\begin{quote}
\textbf{Mediator vs. Observer:} Observer is one-to-many broadcast (subject doesn\textquotesingle t coordinate observers); Mediator is many-to-many coordination (the hub actively orchestrates who does what). They\textquotesingle re often combined (mediator uses observer to listen to components).
\end{quote}

\subsection{Memento}\label{memento}

\textbf{Intent:} Without violating encapsulation, capture and externalize an object\textquotesingle s internal state so the object can be \textbf{restored} to this state later (undo / snapshots / checkpoints).

\textbf{Problem it solves:} You want undo, but the editor\textquotesingle s internal state is private --- exposing it to the undo-manager would break encapsulation. Memento lets the object produce an opaque snapshot (the memento) that only it can read back, while a caretaker just stores snapshots without peeking inside.

\textbf{Structure:}

\setcodelang{}

\begin{verbatim}
Originator (editor) ── save() ──> Memento (opaque state snapshot)
Originator           ── restore(memento) ──> back to that state
Caretaker (history) stores mementos but CANNOT read their contents.
\end{verbatim}

\setcodelang{Python}

\begin{Shaded}
\begin{Highlighting}[]
\KeywordTok{class}\NormalTok{ EditorMemento:                       }\CommentTok{\# opaque snapshot}
    \KeywordTok{def} \FunctionTok{\_\_init\_\_}\NormalTok{(}\VariableTok{self}\NormalTok{, state): }\VariableTok{self}\NormalTok{.\_state }\OperatorTok{=}\NormalTok{ state   }\CommentTok{\# "private" to originator}
    \KeywordTok{def}\NormalTok{ \_get(}\VariableTok{self}\NormalTok{): }\ControlFlowTok{return} \VariableTok{self}\NormalTok{.\_state}

\KeywordTok{class}\NormalTok{ Editor:                              }\CommentTok{\# ORIGINATOR}
    \KeywordTok{def} \FunctionTok{\_\_init\_\_}\NormalTok{(}\VariableTok{self}\NormalTok{): }\VariableTok{self}\NormalTok{.content }\OperatorTok{=} \StringTok{""}
    \KeywordTok{def} \BuiltInTok{type}\NormalTok{(}\VariableTok{self}\NormalTok{, s): }\VariableTok{self}\NormalTok{.content }\OperatorTok{+=}\NormalTok{ s}
    \KeywordTok{def}\NormalTok{ save(}\VariableTok{self}\NormalTok{): }\ControlFlowTok{return}\NormalTok{ EditorMemento(}\VariableTok{self}\NormalTok{.content)        }\CommentTok{\# create snapshot}
    \KeywordTok{def}\NormalTok{ restore(}\VariableTok{self}\NormalTok{, m): }\VariableTok{self}\NormalTok{.content }\OperatorTok{=}\NormalTok{ m.\_get()             }\CommentTok{\# restore snapshot}

\KeywordTok{class}\NormalTok{ History:                             }\CommentTok{\# CARETAKER — stores, never inspects}
    \KeywordTok{def} \FunctionTok{\_\_init\_\_}\NormalTok{(}\VariableTok{self}\NormalTok{): }\VariableTok{self}\NormalTok{.\_mementos }\OperatorTok{=}\NormalTok{ []}
    \KeywordTok{def}\NormalTok{ push(}\VariableTok{self}\NormalTok{, m): }\VariableTok{self}\NormalTok{.\_mementos.append(m)}
    \KeywordTok{def}\NormalTok{ pop(}\VariableTok{self}\NormalTok{): }\ControlFlowTok{return} \VariableTok{self}\NormalTok{.\_mementos.pop()}

\NormalTok{editor }\OperatorTok{=}\NormalTok{ Editor()}\OperatorTok{;}\NormalTok{ history }\OperatorTok{=}\NormalTok{ History()}
\NormalTok{editor.}\BuiltInTok{type}\NormalTok{(}\StringTok{"hello"}\NormalTok{)}\OperatorTok{;}\NormalTok{ history.push(editor.save())   }\CommentTok{\# checkpoint}
\NormalTok{editor.}\BuiltInTok{type}\NormalTok{(}\StringTok{" world"}\NormalTok{)}
\NormalTok{editor.restore(history.pop())                        }\CommentTok{\# undo → "hello"}
\end{Highlighting}
\end{Shaded}

\textbf{When to use:} Undo/redo, checkpoints, snapshots, transactional rollback, save-game systems --- anywhere you must restore prior state without exposing internals. \textbf{When NOT to use:} When state is huge (snapshots cost memory --- consider command-based undo or incremental deltas instead); when state is trivially public.

\textbf{Real framework example:} \texttt{java.io.Serializable} snapshots, database savepoints/rollback, editor undo stacks, Redux time-travel debugging, VM snapshots, \texttt{git\ stash}.

\textbf{Trade-offs:} Preserves encapsulation, clean undo/snapshot, caretaker stays dumb. Memory cost for many/large snapshots, lifecycle management of mementos, deep-copy concerns. \textbf{Command + Memento} is a common undo combo (command captures a memento before executing).

\subsection{Visitor}\label{visitor}

\textbf{Intent:} Represent an operation to be performed on the elements of an object structure. Visitor lets you define a \textbf{new operation} without changing the classes of the elements on which it operates.

\textbf{Problem it solves:} You have a stable class hierarchy (AST nodes: \texttt{NumberNode}, \texttt{AddNode}, \texttt{MulNode}) and you keep needing \emph{new operations} over it (evaluate, pretty-print, type-check, optimize). Adding each operation as a method on every node class bloats them and edits the whole hierarchy each time. Visitor moves operations \emph{out} into visitor classes; each node just \texttt{accept(visitor)}s and calls back \texttt{visitor.visitX(this)} (\textbf{double dispatch}).

\textbf{Structure:}

\setcodelang{}

\begin{verbatim}
Element { accept(v): v.visitThis(this) }        Visitor { visitNumber(n); visitAdd(a) }
   ▲                                                ▲
NumberNode  AddNode  MulNode                EvalVisitor   PrintVisitor   TypeVisitor
   (accept calls the right visit method — DOUBLE DISPATCH)
\end{verbatim}

\setcodelang{Python}

\begin{Shaded}
\begin{Highlighting}[]
\ImportTok{from}\NormalTok{ abc }\ImportTok{import}\NormalTok{ ABC, abstractmethod}

\KeywordTok{class}\NormalTok{ Node(ABC):}
    \AttributeTok{@abstractmethod}
    \KeywordTok{def}\NormalTok{ accept(}\VariableTok{self}\NormalTok{, visitor): ...}

\KeywordTok{class}\NormalTok{ Number(Node):}
    \KeywordTok{def} \FunctionTok{\_\_init\_\_}\NormalTok{(}\VariableTok{self}\NormalTok{, val): }\VariableTok{self}\NormalTok{.val }\OperatorTok{=}\NormalTok{ val}
    \KeywordTok{def}\NormalTok{ accept(}\VariableTok{self}\NormalTok{, v): }\ControlFlowTok{return}\NormalTok{ v.visit\_number(}\VariableTok{self}\NormalTok{)     }\CommentTok{\# double dispatch}
\KeywordTok{class}\NormalTok{ Add(Node):}
    \KeywordTok{def} \FunctionTok{\_\_init\_\_}\NormalTok{(}\VariableTok{self}\NormalTok{, l, r): }\VariableTok{self}\NormalTok{.l, }\VariableTok{self}\NormalTok{.r }\OperatorTok{=}\NormalTok{ l, r}
    \KeywordTok{def}\NormalTok{ accept(}\VariableTok{self}\NormalTok{, v): }\ControlFlowTok{return}\NormalTok{ v.visit\_add(}\VariableTok{self}\NormalTok{)}

\KeywordTok{class}\NormalTok{ Visitor(ABC):}
    \AttributeTok{@abstractmethod}
    \KeywordTok{def}\NormalTok{ visit\_number(}\VariableTok{self}\NormalTok{, n): ...}
    \AttributeTok{@abstractmethod}
    \KeywordTok{def}\NormalTok{ visit\_add(}\VariableTok{self}\NormalTok{, a): ...}

\KeywordTok{class}\NormalTok{ EvalVisitor(Visitor):                  }\CommentTok{\# NEW operation, no Node edits}
    \KeywordTok{def}\NormalTok{ visit\_number(}\VariableTok{self}\NormalTok{, n): }\ControlFlowTok{return}\NormalTok{ n.val}
    \KeywordTok{def}\NormalTok{ visit\_add(}\VariableTok{self}\NormalTok{, a): }\ControlFlowTok{return}\NormalTok{ a.l.accept(}\VariableTok{self}\NormalTok{) }\OperatorTok{+}\NormalTok{ a.r.accept(}\VariableTok{self}\NormalTok{)}

\KeywordTok{class}\NormalTok{ PrintVisitor(Visitor):                 }\CommentTok{\# ANOTHER new operation}
    \KeywordTok{def}\NormalTok{ visit\_number(}\VariableTok{self}\NormalTok{, n): }\ControlFlowTok{return} \BuiltInTok{str}\NormalTok{(n.val)}
    \KeywordTok{def}\NormalTok{ visit\_add(}\VariableTok{self}\NormalTok{, a): }\ControlFlowTok{return} \SpecialStringTok{f"(}\SpecialCharTok{\{}\NormalTok{a}\SpecialCharTok{.}\NormalTok{l}\SpecialCharTok{.}\NormalTok{accept(}\VariableTok{self}\NormalTok{)}\SpecialCharTok{\}}\SpecialStringTok{ + }\SpecialCharTok{\{}\NormalTok{a}\SpecialCharTok{.}\NormalTok{r}\SpecialCharTok{.}\NormalTok{accept(}\VariableTok{self}\NormalTok{)}\SpecialCharTok{\}}\SpecialStringTok{)"}

\NormalTok{tree }\OperatorTok{=}\NormalTok{ Add(Number(}\DecValTok{3}\NormalTok{), Add(Number(}\DecValTok{4}\NormalTok{), Number(}\DecValTok{5}\NormalTok{)))}
\BuiltInTok{print}\NormalTok{(tree.accept(EvalVisitor()))    }\CommentTok{\# 12}
\BuiltInTok{print}\NormalTok{(tree.accept(PrintVisitor()))   }\CommentTok{\# (3 + (4 + 5))}
\end{Highlighting}
\end{Shaded}

\textbf{When to use:} A \emph{stable} element hierarchy plus \emph{many evolving operations} over it (compilers/AST, document object models, reporting over a fixed domain model). \textbf{When NOT to use:} When the \emph{element hierarchy} changes often --- adding a new element type forces editing \emph{every} visitor (the dual of its strength). This is the famous \textbf{expression problem} trade-off.

\textbf{Real framework example:} Java annotation processing (\texttt{javax.lang.model.element.ElementVisitor}), the JDK \texttt{FileVisitor} (\texttt{Files.walkFileTree}), compiler/AST traversals (ANTLR visitors), \texttt{NodeVisitor} in Python\textquotesingle s \texttt{ast} module, ASM bytecode visitors.

\textbf{Trade-offs:} Add operations without touching elements (OCP for operations), gather related logic in one visitor, accumulate state across a traversal. Adding an \emph{element} type breaks all visitors, double-dispatch boilerplate, visitors may need access to element internals (breaks encapsulation), hard to read for newcomers.

\subsection{Interpreter}\label{interpreter}

\textbf{Intent:} Given a language, define a representation for its grammar along with an interpreter that uses the representation to interpret sentences in the language.

\textbf{Problem it solves:} You have a small, well-defined language to evaluate repeatedly --- a boolean rule engine, a search filter syntax, simple arithmetic, regex-like matching. Interpreter models each grammar rule as a class; an expression tree of these objects evaluates a sentence via recursion.

\textbf{Structure:}

\setcodelang{}

\begin{verbatim}
«interface» Expression { interpret(context): result }
   ▲
TerminalExpr (Number, Variable)    NonTerminalExpr (Add, Or, And — hold sub-Expressions)
   interpret() returns value         interpret() combines children's interpret()
\end{verbatim}

\setcodelang{Python}

\begin{Shaded}
\begin{Highlighting}[]
\ImportTok{from}\NormalTok{ abc }\ImportTok{import}\NormalTok{ ABC, abstractmethod}

\KeywordTok{class}\NormalTok{ Expr(ABC):}
    \AttributeTok{@abstractmethod}
    \KeywordTok{def}\NormalTok{ interpret(}\VariableTok{self}\NormalTok{, ctx: }\BuiltInTok{dict}\NormalTok{): ...}

\KeywordTok{class}\NormalTok{ Var(Expr):                              }\CommentTok{\# terminal}
    \KeywordTok{def} \FunctionTok{\_\_init\_\_}\NormalTok{(}\VariableTok{self}\NormalTok{, name): }\VariableTok{self}\NormalTok{.name }\OperatorTok{=}\NormalTok{ name}
    \KeywordTok{def}\NormalTok{ interpret(}\VariableTok{self}\NormalTok{, ctx): }\ControlFlowTok{return}\NormalTok{ ctx[}\VariableTok{self}\NormalTok{.name]}

\KeywordTok{class}\NormalTok{ And(Expr):                              }\CommentTok{\# non{-}terminal}
    \KeywordTok{def} \FunctionTok{\_\_init\_\_}\NormalTok{(}\VariableTok{self}\NormalTok{, l, r): }\VariableTok{self}\NormalTok{.l, }\VariableTok{self}\NormalTok{.r }\OperatorTok{=}\NormalTok{ l, r}
    \KeywordTok{def}\NormalTok{ interpret(}\VariableTok{self}\NormalTok{, ctx): }\ControlFlowTok{return} \VariableTok{self}\NormalTok{.l.interpret(ctx) }\KeywordTok{and} \VariableTok{self}\NormalTok{.r.interpret(ctx)}

\KeywordTok{class}\NormalTok{ Or(Expr):}
    \KeywordTok{def} \FunctionTok{\_\_init\_\_}\NormalTok{(}\VariableTok{self}\NormalTok{, l, r): }\VariableTok{self}\NormalTok{.l, }\VariableTok{self}\NormalTok{.r }\OperatorTok{=}\NormalTok{ l, r}
    \KeywordTok{def}\NormalTok{ interpret(}\VariableTok{self}\NormalTok{, ctx): }\ControlFlowTok{return} \VariableTok{self}\NormalTok{.l.interpret(ctx) }\KeywordTok{or} \VariableTok{self}\NormalTok{.r.interpret(ctx)}

\CommentTok{\# Rule: "premium AND (verified OR trusted)"}
\NormalTok{rule }\OperatorTok{=}\NormalTok{ And(Var(}\StringTok{"premium"}\NormalTok{), Or(Var(}\StringTok{"verified"}\NormalTok{), Var(}\StringTok{"trusted"}\NormalTok{)))}
\BuiltInTok{print}\NormalTok{(rule.interpret(\{}\StringTok{"premium"}\NormalTok{: }\VariableTok{True}\NormalTok{, }\StringTok{"verified"}\NormalTok{: }\VariableTok{False}\NormalTok{, }\StringTok{"trusted"}\NormalTok{: }\VariableTok{True}\NormalTok{\}))  }\CommentTok{\# True}
\end{Highlighting}
\end{Shaded}

\textbf{When to use:} A \emph{simple, stable} grammar evaluated often (rule engines, filters, small DSLs, arithmetic/boolean expressions). \textbf{When NOT to use:} Complex or evolving grammars --- hand-rolling each rule as a class becomes unmanageable; use a real parser generator (ANTLR, yacc) instead. It\textquotesingle s the least-used GoF pattern in practice.

\textbf{Real framework example:} \texttt{java.util.regex.Pattern} (a compiled regex is an interpreter), \texttt{java.text.Format}, Spring Expression Language (SpEL), rule engines (Drools), database query planners conceptually.

\textbf{Trade-offs:} Easy to extend the grammar (new rule = new class), maps grammar directly to classes. Scales terribly for complex grammars (class per rule), performance overhead vs. a real parser, often the wrong tool --- reach for a parser library.

\subsection{Behavioral Patterns --- Comparison}\label{behavioral-patterns--comparison}

\begin{longtable}[]{@{}
  >{\raggedright\arraybackslash}p{(\columnwidth - 4\tabcolsep) * \real{0.2279}}
  >{\raggedright\arraybackslash}p{(\columnwidth - 4\tabcolsep) * \real{0.3667}}
  >{\raggedright\arraybackslash}p{(\columnwidth - 4\tabcolsep) * \real{0.4054}}@{}}
\toprule\noalign{}
\rowcolor{acc!16}
\begin{minipage}[b]{\linewidth}\raggedright
Pattern
\end{minipage} & \begin{minipage}[b]{\linewidth}\raggedright
Intent
\end{minipage} & \begin{minipage}[b]{\linewidth}\raggedright
Smell it replaces
\end{minipage} \\
\midrule\noalign{}
\endhead
\bottomrule\noalign{}
\endlastfoot
\textbf{Strategy} & Swap interchangeable algorithms & \texttt{if\ type\ ==\ ...} algorithm selection \\
\rowcolor{zebra}
\textbf{Observer} & Notify many dependents on change & Polling / tight producer-consumer coupling \\
\textbf{Command} & Request as an object (undo/queue/log) & No undo, can\textquotesingle t queue/replay actions \\
\rowcolor{zebra}
\textbf{State} & Behavior changes with internal state & Giant \texttt{switch(status)} in every method \\
\textbf{Template Method} & Fixed skeleton, overridable steps & Duplicated algorithm with small diffs \\
\rowcolor{zebra}
\textbf{Iterator} & Sequential access, hide structure & Exposing collection internals \\
\textbf{Chain of Responsibility} & Pipeline of handlers & Nested \texttt{if} deciding who handles \\
\rowcolor{zebra}
\textbf{Mediator} & Centralize many-to-many comms & n\(^2\) object interconnections \\
\textbf{Memento} & Snapshot/restore state & Exposing internals for undo \\
\rowcolor{zebra}
\textbf{Visitor} & New ops without editing elements & Methods bloating a stable hierarchy \\
\textbf{Interpreter} & Evaluate sentences of a grammar & Ad-hoc parsing of a small language \\
\end{longtable}

\section{Modern / Architectural Patterns}\label{modern--architectural-patterns}

The GoF book predates the web, dependency-injection containers, and reactive UIs. These patterns are everyday currency in modern codebases and appear constantly in interviews.

\subsection{Dependency Injection / Inversion of Control (IoC)}\label{dependency-injection--inversion-of-control-ioc}

\textbf{Intent:} Instead of an object creating its own dependencies (\texttt{this.db\ =\ new\ MySqlDb()}), the dependencies are \textbf{supplied} ("injected") from outside --- usually by a container --- so the object depends only on abstractions.

\textbf{Problem it solves:} Hard-coded \texttt{new\ ConcreteClass()} calls weld a class to its collaborators: you can\textquotesingle t substitute a mock in tests, can\textquotesingle t swap implementations, and creation logic spreads everywhere. DI inverts control of construction: a container (or the caller) builds the graph and hands each object what it needs.

\setcodelang{Java}

\begin{Shaded}
\begin{Highlighting}[]
\CommentTok{// WITHOUT DI — OrderService is welded to MySqlRepo; untestable, inflexible}
\KeywordTok{class}\NormalTok{ OrderService }\OperatorTok{\{}
    \KeywordTok{private} \DataTypeTok{final}\NormalTok{ Repo repo }\OperatorTok{=} \KeywordTok{new} \FunctionTok{MySqlRepo}\OperatorTok{();}   \CommentTok{// hard dependency}
\OperatorTok{\}}

\CommentTok{// WITH DI (constructor injection — the preferred form)}
\KeywordTok{class}\NormalTok{ OrderService }\OperatorTok{\{}
    \KeywordTok{private} \DataTypeTok{final}\NormalTok{ Repo repo}\OperatorTok{;}
    \FunctionTok{OrderService}\OperatorTok{(}\NormalTok{Repo repo}\OperatorTok{)} \OperatorTok{\{} \KeywordTok{this}\OperatorTok{.}\FunctionTok{repo} \OperatorTok{=}\NormalTok{ repo}\OperatorTok{;} \OperatorTok{\}} \CommentTok{// depends on the ABSTRACTION}
\OperatorTok{\}}
\CommentTok{// Production: new OrderService(new MySqlRepo());}
\CommentTok{// Test:       new OrderService(new InMemoryRepo());   // trivial to mock}
\end{Highlighting}
\end{Shaded}

\textbf{Three injection styles:} \textbf{constructor} (preferred --- guarantees fully-initialized, immutable, required deps explicit), \textbf{setter} (optional deps, mutable), \textbf{field} (\texttt{@Autowired} on a field --- concise but hides deps and hurts testability). \textbf{When NOT to use:} tiny scripts where a container is overkill; over-injecting (5+ constructor args signals an SRP violation).

\textbf{Real framework example:} Spring (\texttt{@Component}/\texttt{@Autowired}, the \texttt{ApplicationContext} is the IoC container), Google Guice, Dagger (compile-time DI), Angular\textquotesingle s injector, .NET \texttt{IServiceCollection}, Python\textquotesingle s \texttt{dependency-injector}/FastAPI \texttt{Depends}.

\textbf{Trade-offs:} Testability (inject mocks), loose coupling, centralized wiring, swap implementations by config. Indirection/"magic" (where did this object come from?), container learning curve, runtime wiring errors, can be overused. \textbf{DI is how you get Singleton\textquotesingle s "one instance" without Singleton\textquotesingle s global-state downsides} --- the container owns one bean and injects it.

\subsection{MVC / MVP / MVVM}\label{mvc--mvp--mvvm}

\textbf{Intent:} Separate the \textbf{data (Model)}, the \textbf{UI (View)}, and the \textbf{glue logic} so each evolves independently and the model has no UI knowledge.

\begin{longtable}[]{@{}
  >{\raggedright\arraybackslash}p{(\columnwidth - 6\tabcolsep) * \real{0.1346}}
  >{\raggedright\arraybackslash}p{(\columnwidth - 6\tabcolsep) * \real{0.2212}}
  >{\raggedright\arraybackslash}p{(\columnwidth - 6\tabcolsep) * \real{0.3173}}
  >{\raggedright\arraybackslash}p{(\columnwidth - 6\tabcolsep) * \real{0.3269}}@{}}
\toprule\noalign{}
\rowcolor{acc!16}
\begin{minipage}[b]{\linewidth}\raggedright
\end{minipage} & \begin{minipage}[b]{\linewidth}\raggedright
MVC
\end{minipage} & \begin{minipage}[b]{\linewidth}\raggedright
MVP
\end{minipage} & \begin{minipage}[b]{\linewidth}\raggedright
MVVM
\end{minipage} \\
\midrule\noalign{}
\endhead
\bottomrule\noalign{}
\endlastfoot
\textbf{Glue} & Controller & Presenter & ViewModel \\
\rowcolor{zebra}
\textbf{View knows} & Controller + Model & Presenter only & ViewModel (via binding) \\
\textbf{Binding} & Manual & Manual (presenter pushes to view) & \textbf{Data binding} (automatic) \\
\rowcolor{zebra}
\textbf{View passivity} & Active (observes model) & Passive (dumb) & Reactive (binds) \\
\textbf{Typical use} & Web (Rails, Spring MVC) & Android (older), WinForms & Android Jetpack, WPF, Vue/Knockout \\
\end{longtable}

\setcodelang{}

\begin{verbatim}
MVC flow:
  User → View → Controller → updates Model → Model notifies View (Observer) → View re-renders
                              (Controller = Strategy for handling input)
\end{verbatim}

\textbf{Why it matters:} MVC is itself a \emph{composition of patterns} --- the Model\(\leftrightarrow\)View link is \textbf{Observer}, the Controller\textquotesingle s input handling is \textbf{Strategy}, and a View composed of nested widgets is \textbf{Composite}. This is the canonical "patterns combine" example. \textbf{Trade-offs:} separation of concerns, testable model, parallel work. boilerplate, "Massive View Controller" anti-pattern when the controller absorbs everything (SRP violation).

\subsection{Repository}\label{repository}

\textbf{Intent:} Mediate between the domain and data-mapping layers, presenting an \textbf{in-memory collection-like interface} to domain objects (\texttt{save}, \texttt{findById}, \texttt{findByEmail}) while hiding the persistence mechanism.

\textbf{Problem it solves:} Scattering raw SQL / ORM calls through business logic couples your domain to the database, makes testing require a real DB, and means swapping Postgres\(\rightarrow\)Mongo touches everything. A Repository abstracts persistence behind a collection-like interface; business code depends on the interface, not the data store.

\setcodelang{Python}

\begin{Shaded}
\begin{Highlighting}[]
\ImportTok{from}\NormalTok{ abc }\ImportTok{import}\NormalTok{ ABC, abstractmethod}

\KeywordTok{class}\NormalTok{ UserRepository(ABC):                  }\CommentTok{\# the abstraction business code depends on}
    \AttributeTok{@abstractmethod}
    \KeywordTok{def}\NormalTok{ find\_by\_id(}\VariableTok{self}\NormalTok{, }\BuiltInTok{id}\NormalTok{): ...}
    \AttributeTok{@abstractmethod}
    \KeywordTok{def}\NormalTok{ save(}\VariableTok{self}\NormalTok{, user): ...}

\KeywordTok{class}\NormalTok{ SqlUserRepository(UserRepository):     }\CommentTok{\# production impl}
    \KeywordTok{def}\NormalTok{ find\_by\_id(}\VariableTok{self}\NormalTok{, }\BuiltInTok{id}\NormalTok{): ...            }\CommentTok{\# real SQL here}
    \KeywordTok{def}\NormalTok{ save(}\VariableTok{self}\NormalTok{, user): ...}

\KeywordTok{class}\NormalTok{ InMemoryUserRepository(UserRepository):  }\CommentTok{\# test impl — no DB needed}
    \KeywordTok{def} \FunctionTok{\_\_init\_\_}\NormalTok{(}\VariableTok{self}\NormalTok{): }\VariableTok{self}\NormalTok{.\_store }\OperatorTok{=}\NormalTok{ \{\}}
    \KeywordTok{def}\NormalTok{ find\_by\_id(}\VariableTok{self}\NormalTok{, }\BuiltInTok{id}\NormalTok{): }\ControlFlowTok{return} \VariableTok{self}\NormalTok{.\_store.get(}\BuiltInTok{id}\NormalTok{)}
    \KeywordTok{def}\NormalTok{ save(}\VariableTok{self}\NormalTok{, user): }\VariableTok{self}\NormalTok{.\_store[user.}\BuiltInTok{id}\NormalTok{] }\OperatorTok{=}\NormalTok{ user}
\end{Highlighting}
\end{Shaded}

\textbf{When to use:} Domain-driven designs, when you want testable business logic and the freedom to swap data stores. \textbf{When NOT to use:} Trivial CRUD apps where the ORM already gives a clean enough abstraction (a repository over a repository is ceremony). \textbf{Real framework example:} Spring Data \texttt{JpaRepository}, .NET EF repositories, Django\textquotesingle s manager/queryset layer, Ruby\textquotesingle s ActiveRecord (Active Record pattern --- a \emph{different} trade-off where the model knows how to persist itself). \textbf{Trade-offs:} testable, swappable persistence, centralizes queries. extra layer, can leak abstractions (generic repos that just wrap the ORM add little).

\subsection{Null Object}\label{null-object}

\textbf{Intent:} Provide an object with \textbf{neutral ("do nothing") behavior} to represent the absence of a value --- eliminating \texttt{null} checks.

\textbf{Problem it solves:} Code littered with \texttt{if\ (logger\ !=\ null)\ logger.log(...)}. A Null Object is a real object implementing the same interface but doing nothing (or returning a safe default), so callers never special-case absence.

\setcodelang{Java}

\begin{Shaded}
\begin{Highlighting}[]
\KeywordTok{interface} \BuiltInTok{Logger} \OperatorTok{\{} \DataTypeTok{void} \FunctionTok{log}\OperatorTok{(}\BuiltInTok{String}\NormalTok{ msg}\OperatorTok{);} \OperatorTok{\}}
\KeywordTok{class}\NormalTok{ RealLogger }\KeywordTok{implements} \BuiltInTok{Logger} \OperatorTok{\{} \KeywordTok{public} \DataTypeTok{void} \FunctionTok{log}\OperatorTok{(}\BuiltInTok{String}\NormalTok{ m}\OperatorTok{)\{} \BuiltInTok{System}\OperatorTok{.}\FunctionTok{out}\OperatorTok{.}\FunctionTok{println}\OperatorTok{(}\NormalTok{m}\OperatorTok{);} \OperatorTok{\}} \OperatorTok{\}}
\KeywordTok{class}\NormalTok{ NullLogger }\KeywordTok{implements} \BuiltInTok{Logger} \OperatorTok{\{} \KeywordTok{public} \DataTypeTok{void} \FunctionTok{log}\OperatorTok{(}\BuiltInTok{String}\NormalTok{ m}\OperatorTok{)\{} \CommentTok{/* do nothing */} \OperatorTok{\}} \OperatorTok{\}}

\CommentTok{// Caller never checks for null:}
\BuiltInTok{Logger}\NormalTok{ logger }\OperatorTok{=}\NormalTok{ config}\OperatorTok{.}\FunctionTok{isLoggingOn}\OperatorTok{()} \OperatorTok{?} \KeywordTok{new} \FunctionTok{RealLogger}\OperatorTok{()} \OperatorTok{:} \KeywordTok{new} \FunctionTok{NullLogger}\OperatorTok{();}
\NormalTok{logger}\OperatorTok{.}\FunctionTok{log}\OperatorTok{(}\StringTok{"hello"}\OperatorTok{);}   \CommentTok{// safe whether logging is on or off}
\end{Highlighting}
\end{Shaded}

\textbf{When to use:} A sensible "do nothing" default exists and you want to kill null checks (loggers, no-op listeners, empty collections). \textbf{When NOT to use:} When absence is genuinely an error that must surface (silently doing nothing can hide bugs). \textbf{Real framework example:} \texttt{Collections.emptyList()}, \texttt{Optional.empty()} (a richer cousin), SLF4J \texttt{NOPLogger}, Spring\textquotesingle s no-op beans. \textbf{Trade-offs:} removes null checks (LSP-friendly), simpler call sites. can mask errors (a silent no-op when something was expected to happen).

\subsection{Object Pool}\label{object-pool}

\textbf{Intent:} Reuse a fixed set of \textbf{expensive-to-create objects} by borrowing from and returning to a pool, rather than constructing/destroying them repeatedly.

\textbf{Problem it solves:} Opening a DB connection or spinning a thread is expensive (TCP handshake, auth, OS thread setup). Creating one per request kills performance. An Object Pool pre-creates a bounded set, hands them out, and recycles them on return.

\setcodelang{Python}

\begin{Shaded}
\begin{Highlighting}[]
\ImportTok{import}\NormalTok{ queue}

\KeywordTok{class}\NormalTok{ ConnectionPool:}
    \KeywordTok{def} \FunctionTok{\_\_init\_\_}\NormalTok{(}\VariableTok{self}\NormalTok{, size):}
        \VariableTok{self}\NormalTok{.\_pool }\OperatorTok{=}\NormalTok{ queue.Queue(maxsize}\OperatorTok{=}\NormalTok{size)}
        \ControlFlowTok{for}\NormalTok{ \_ }\KeywordTok{in} \BuiltInTok{range}\NormalTok{(size):}
            \VariableTok{self}\NormalTok{.\_pool.put(}\VariableTok{self}\NormalTok{.\_create())      }\CommentTok{\# pre{-}create expensive objects}
    \KeywordTok{def}\NormalTok{ \_create(}\VariableTok{self}\NormalTok{): }\ControlFlowTok{return} \BuiltInTok{object}\NormalTok{()          }\CommentTok{\# imagine: a real DB connection}
    \KeywordTok{def}\NormalTok{ acquire(}\VariableTok{self}\NormalTok{): }\ControlFlowTok{return} \VariableTok{self}\NormalTok{.\_pool.get()  }\CommentTok{\# borrow (blocks if empty)}
    \KeywordTok{def}\NormalTok{ release(}\VariableTok{self}\NormalTok{, conn): }\VariableTok{self}\NormalTok{.\_pool.put(conn)  }\CommentTok{\# return for reuse}

\NormalTok{pool }\OperatorTok{=}\NormalTok{ ConnectionPool(}\DecValTok{5}\NormalTok{)}
\NormalTok{conn }\OperatorTok{=}\NormalTok{ pool.acquire()}
\ControlFlowTok{try}\NormalTok{:    }\ControlFlowTok{pass}                \CommentTok{\# use the connection}
\ControlFlowTok{finally}\NormalTok{: pool.release(conn) }\CommentTok{\# ALWAYS return it (else pool leaks/exhausts)}
\end{Highlighting}
\end{Shaded}

\textbf{When to use:} Expensive, reusable, interchangeable resources under high churn (DB connections, threads, sockets, large buffers, game objects). \textbf{When NOT to use:} Cheap objects (pooling adds overhead and bugs), or when objects hold state that must be reset (forgetting to reset leaks state between borrowers). \textbf{Real framework example:} HikariCP / JDBC connection pools, \texttt{java.util.concurrent.ThreadPoolExecutor}, Apache Commons Pool, Netty buffer pools, game engine object pools. \textbf{Trade-offs:} amortizes creation cost, bounds resource usage, predictable performance. complexity (acquire/release discipline --- leaks if not returned), must reset state between uses, sizing is tricky (too small = contention, too large = waste), thread-safety required.

\section{Which Pattern Do I Use? (Decision Guide)}\label{which-pattern-do-i-use-decision-guide}

Map the \emph{symptom} in the requirement to the pattern. In an interview, say the trigger phrase out loud.

\begin{longtable}[]{@{}
  >{\raggedright\arraybackslash}p{(\columnwidth - 4\tabcolsep) * \real{0.4120}}
  >{\raggedright\arraybackslash}p{(\columnwidth - 4\tabcolsep) * \real{0.3199}}
  >{\raggedright\arraybackslash}p{(\columnwidth - 4\tabcolsep) * \real{0.2681}}@{}}
\toprule\noalign{}
\rowcolor{acc!16}
\begin{minipage}[b]{\linewidth}\raggedright
The requirement says\ldots{}
\end{minipage} & \begin{minipage}[b]{\linewidth}\raggedright
Reach for
\end{minipage} & \begin{minipage}[b]{\linewidth}\raggedright
Why
\end{minipage} \\
\midrule\noalign{}
\endhead
\bottomrule\noalign{}
\endlastfoot
"There must be exactly one X" & \textbf{Singleton} (or DI-managed single bean) & One instance, controlled access \\
\rowcolor{zebra}
"Create an object but the type depends on input/subclass" & \textbf{Factory Method} & Defer instantiation \\
"Create a \emph{family} of matching objects" & \textbf{Abstract Factory} & Guarantee compatible products \\
\rowcolor{zebra}
"Many optional fields / immutable construction" & \textbf{Builder} & Avoid telescoping ctors \\
"Cloning is cheaper than rebuilding" & \textbf{Prototype} & Copy an exemplar \\
\rowcolor{zebra}
"Make this incompatible API fit ours" & \textbf{Adapter} & Interface conversion \\
"Add optional, stackable features at runtime" & \textbf{Decorator} & Compose responsibilities \\
\rowcolor{zebra}
"Hide this messy subsystem behind one call" & \textbf{Facade} & Simplify \\
"Control / defer / secure access to an object" & \textbf{Proxy} & Surrogate with control logic \\
\rowcolor{zebra}
"Treat single items and groups the same (tree)" & \textbf{Composite} & Part-whole hierarchy \\
"Two things vary independently (shape \(\times\) renderer)" & \textbf{Bridge} & Avoid n\(\times\)m explosion \\
\rowcolor{zebra}
"Millions of similar objects, memory-bound" & \textbf{Flyweight} & Share intrinsic state \\
"The algorithm must be swappable at runtime" & \textbf{Strategy} & Interchangeable algorithms \\
\rowcolor{zebra}
"Notify many things when one changes" & \textbf{Observer} & Pub/sub \\
"Need undo/redo, queue, or log of actions" & \textbf{Command} & Request as object \\
\rowcolor{zebra}
"Behavior depends on a status that transitions" & \textbf{State} & State machine \\
"Same flow, a few steps differ" & \textbf{Template Method} & Skeleton + hooks \\
\rowcolor{zebra}
"Traverse a collection without exposing it" & \textbf{Iterator} & Uniform cursor \\
"A request passes through stages until handled" & \textbf{Chain of Responsibility} & Pipeline/middleware \\
\rowcolor{zebra}
"Many components must coordinate (n\(^2\))" & \textbf{Mediator} & Central hub \\
"Snapshot and restore state (undo)" & \textbf{Memento} & Encapsulated snapshot \\
\rowcolor{zebra}
"Add operations to a stable hierarchy" & \textbf{Visitor} & Externalize operations \\
"Evaluate a small custom grammar" & \textbf{Interpreter} & Grammar as classes \\
\rowcolor{zebra}
"Don\textquotesingle t create my own dependencies" & \textbf{Dependency Injection} & Inject abstractions \\
"Separate data, UI, glue" & \textbf{MVC/MVP/MVVM} & Layered UI \\
\rowcolor{zebra}
"Abstract the data store" & \textbf{Repository} & Collection-like persistence \\
"Stop checking for null everywhere" & \textbf{Null Object} & Neutral default \\
\rowcolor{zebra}
"Reuse expensive objects" & \textbf{Object Pool} & Borrow/return \\
\end{longtable}

\begin{quote}
\textbf{Default to the smallest tool.} A plain function, polymorphism, or a first-class function beats a named pattern when the friction isn\textquotesingle t real (KISS/YAGNI). Escalate to a pattern only when you \emph{feel} the specific pain it relieves.
\end{quote}

\section{Patterns in the Wild \& How They Combine}\label{patterns-in-the-wild--how-they-combine}

Real systems are \textbf{layered compositions} of patterns, not single ones. Recognizing the combinations is a senior signal.

\textbf{MVC = Observer + Strategy + Composite.} The View observes the Model (Observer); the Controller plugs in input-handling behavior (Strategy); a View built of nested widgets is a tree (Composite). The classic "patterns collaborate" example.

\textbf{Spring Framework = Factory + Singleton + Proxy + Template Method.} The \texttt{BeanFactory}/\texttt{ApplicationContext} is a Factory producing beans; beans are singleton-scoped by default (DI-managed Singleton); \texttt{@Transactional}/AOP wraps beans in dynamic \textbf{Proxies}; \texttt{JdbcTemplate}/\texttt{RestTemplate} are Template Methods (fixed flow, your callbacks fill steps).

\textbf{Java I/O = Decorator + Adapter + Iterator.} \texttt{BufferedReader(new\ InputStreamReader(new\ FileInputStream(f)))} stacks \textbf{Decorators}; \texttt{InputStreamReader} is an \textbf{Adapter} (bytes\(\rightarrow\)chars); collections traverse via \textbf{Iterator}.

\textbf{React = Observer + Composite + Decorator/Strategy.} Components observe state and re-render (Observer); the component tree is a \textbf{Composite}; Higher-Order Components and hooks wrap behavior (Decorator-ish); render props pass behavior in (Strategy-ish).

\textbf{A plugin system = Factory + Strategy + Observer.} A Factory instantiates plugins by name; each plugin is a Strategy implementing a common interface; the host emits lifecycle events plugins observe.

\textbf{Undo system = Command + Memento.} Each Command captures a Memento (state snapshot) before executing; the history replays/reverts them.

\textbf{Compiler = Visitor + Interpreter + Composite + Factory.} The AST is a Composite tree; passes (type-check, optimize, codegen) are Visitors; a small expression evaluator is an Interpreter; node creation uses Factories.

\textbf{A web request pipeline = Chain of Responsibility + Strategy + Decorator.} Middleware is a Chain; each handler may select behavior via Strategy; route handlers can be wrapped (Decorator) with logging/auth.

\begin{quote}
\textbf{Interview gold:} When asked "how do patterns combine?", give MVC = Observer + Strategy + Composite, or Spring = Factory + Singleton + Proxy. It shows you understand patterns as a \emph{system}, not flashcards.
\end{quote}

\section{Anti-Patterns}\label{anti-patterns}

An anti-pattern is a \emph{common but counterproductive} solution. Recognizing them --- and naming the refactor --- is as valuable as knowing the patterns.

\subsection{God Object (a.k.a. Blob)}\label{god-object-aka-blob}

\textbf{What:} One class that knows/does everything (\texttt{OrderManager} with 4,000 lines handling pricing, payment, email, persistence). \textbf{Why bad:} violates SRP, untestable, every change risks breaking unrelated features, merge-conflict magnet. \textbf{Fix:} Extract responsibilities into focused collaborators (SRP); use Facade if you still want a simple entry point; push behavior into the right domain objects (kill the anemic model).

\subsection{Golden Hammer}\label{golden-hammer}

\textbf{What:} "When all you have is a hammer, everything looks like a nail" --- applying one familiar pattern/tool everywhere (Singleton for every service, inheritance for everything, microservices for a CRUD app). \textbf{Why bad:} forces ill-fitting solutions, adds complexity. \textbf{Fix:} match the tool to the problem; keep a broad toolkit; justify each choice by the \emph{force} it relieves.

\subsection{Spaghetti Code}\label{spaghetti-code}

\textbf{What:} Tangled control flow with no clear structure --- deep nesting, \texttt{goto}-like jumps, hidden global state, no layering. \textbf{Why bad:} unreadable, unmaintainable, fear-driven changes. \textbf{Fix:} introduce layers/modules, extract methods, apply SRP, replace conditionals with polymorphism (Strategy/State), establish clear boundaries.

\subsection{Poltergeist (a.k.a. Gypsy class)}\label{poltergeist-aka-gypsy-class}

\textbf{What:} A short-lived, do-nothing class that just passes control to another (a "manager" that only forwards calls). It haunts the design --- appears, does little, vanishes. \textbf{Why bad:} needless indirection, clutter, obscures real flow. \textbf{Fix:} remove it; call the real worker directly; merge its (trivial) logic into a meaningful class.

\subsection{Premature Abstraction / Speculative Generality}\label{premature-abstraction--speculative-generality}

\textbf{What:} Building interfaces, generics, and extension points for variation that \emph{might} happen but never does (\texttt{AbstractFactoryProviderStrategyFactory} for a feature with one implementation). \textbf{Why bad:} YAGNI violation --- pays complexity now for flexibility never used; harder to read and change. \textbf{Fix:} \textbf{Rule of Three} --- introduce the abstraction on the \emph{third} time you actually need to vary something, not the first. Start concrete, refactor to abstract when a real second case arrives.

\subsection{Singleton Abuse}\label{singleton-abuse}

\textbf{What:} Reaching for Singleton as a glorified global variable for state that isn\textquotesingle t truly single (or just to avoid passing arguments). \textbf{Why bad:} hidden dependencies, global mutable state, untestable (can\textquotesingle t mock), thread-safety hazards, hidden coupling. \textbf{Fix:} use \textbf{Dependency Injection} of a single container-managed instance --- you keep "one instance" but make it explicit, mockable, and swappable.

\subsection{Bonus anti-patterns worth naming}\label{bonus-anti-patterns-worth-naming}

\begin{itemize}
\tightlist
\item
  \textbf{Anemic Domain Model:} classes with only getters/setters, all logic in "Service" classes --- not OOP. \emph{Fix:} put behavior on the object that owns the data.
\item
  \textbf{Lava Flow:} dead code no one dares delete. \emph{Fix:} test coverage + delete.
\item
  \textbf{Boat Anchor:} keeping unused components "just in case." \emph{Fix:} delete; version control remembers.
\item
  \textbf{Yo-Yo Problem:} inheritance so deep you bounce up and down the hierarchy to follow logic. \emph{Fix:} flatten; favor composition.
\end{itemize}

\begin{longtable}[]{@{}
  >{\raggedright\arraybackslash}p{(\columnwidth - 4\tabcolsep) * \real{0.2658}}
  >{\raggedright\arraybackslash}p{(\columnwidth - 4\tabcolsep) * \real{0.3671}}
  >{\raggedright\arraybackslash}p{(\columnwidth - 4\tabcolsep) * \real{0.3671}}@{}}
\toprule\noalign{}
\rowcolor{acc!16}
\begin{minipage}[b]{\linewidth}\raggedright
Anti-pattern
\end{minipage} & \begin{minipage}[b]{\linewidth}\raggedright
Core problem
\end{minipage} & \begin{minipage}[b]{\linewidth}\raggedright
The fix
\end{minipage} \\
\midrule\noalign{}
\endhead
\bottomrule\noalign{}
\endlastfoot
God Object & One class does everything & SRP --- extract collaborators \\
\rowcolor{zebra}
Golden Hammer & One pattern for every problem & Match tool to problem \\
Spaghetti & No structure/layering & Layer, extract, polymorphism \\
\rowcolor{zebra}
Poltergeist & Useless pass-through class & Delete, call worker directly \\
Premature Abstraction & Flexibility never used & Rule of Three; start concrete \\
\rowcolor{zebra}
Singleton Abuse & Global mutable state & Dependency Injection \\
\end{longtable}

\section{Architecture / Diagrams}\label{architecture--diagrams}

\textbf{Strategy (composition + OCP):}

\setcodelang{}

\begin{verbatim}
Context ──has-a──> «interface» Strategy { execute() }
                          ▲ realizes
        ┌─────────────────┼─────────────────┐
   ConcreteA          ConcreteB          ConcreteC
   execute()          execute()          execute()
   (add ConcreteD = new class, ZERO edits to Context — OCP)
\end{verbatim}

\textbf{Decorator stacking (composition over inheritance):}

\setcodelang{}

\begin{verbatim}
[FileSource] ──wrapped by──> [Compression] ──wrapped by──> [Encryption]
   read()=raw                  +gzip layer                  +crypto layer
   write("hi") → encrypt → compress → file    (each adds one responsibility)
\end{verbatim}

\textbf{Observer (one-to-many broadcast):}

\setcodelang{}

\begin{verbatim}
            ┌──> Observer A (chart)
Subject ────┼──> Observer B (alert)     notify(): for o in observers: o.update()
(stock)     └──> Observer C (logger)
   subscribe(o) / unsubscribe(o)   ← unsubscribe to avoid lapsed-listener leak
\end{verbatim}

\textbf{State (object swaps its own behavior):}

\setcodelang{}

\begin{verbatim}
Order ──delegates──> «State»
 Pending ──pay()──> Paid ──ship()──> Shipped ──deliver()──> Delivered
   │                                                          
   └── ship() while Pending → "illegal" (state rejects it)
\end{verbatim}

\textbf{Composite (uniform tree):}

\setcodelang{}

\begin{verbatim}
Directory("root")  ── size() recurses ──┐
 ├── File("a.txt", 100)                  │  client calls size() on ANY node
 └── Directory("docs")                   │  — leaf or composite — identically
      └── File("b.pdf", 2000)            ┘
\end{verbatim}

\textbf{Chain of Responsibility (pipeline):}

\setcodelang{}

\begin{verbatim}
request → [Auth] → [RateLimit] → [Validate] → [Handler]
            │each handler: handle and/or pass to next; may short-circuit (reject)
\end{verbatim}

\textbf{Pattern family map:}

\setcodelang{}

\begin{verbatim}
                     DESIGN PATTERNS
        ┌───────────────┼───────────────┐
   CREATIONAL       STRUCTURAL       BEHAVIORAL
   (create)         (compose)        (communicate)
   Singleton        Adapter          Strategy   State
   Factory Method   Decorator        Observer   Iterator
   Abstract Factory Facade           Command    Chain of Resp.
   Builder          Proxy            Template   Mediator
   Prototype        Composite        Memento    Visitor
                    Bridge           Interpreter
                    Flyweight
\end{verbatim}

\section{Real-World Examples}\label{real-world-examples}

\subsection{Java Standard Library}\label{java-standard-library}

\begin{itemize}
\tightlist
\item
  \textbf{Decorator:} \texttt{BufferedReader}, \texttt{BufferedInputStream}, \texttt{GZIPInputStream}, \texttt{Collections.synchronizedList()}.
\item
  \textbf{Adapter:} \texttt{InputStreamReader} (bytes\(\rightarrow\)chars), \texttt{Arrays.asList()}.
\item
  \textbf{Factory Method:} \texttt{Calendar.getInstance()}, \texttt{NumberFormat.getInstance()}, \texttt{Optional.of()}.
\item
  \textbf{Singleton:} \texttt{Runtime.getRuntime()}, enums.
\item
  \textbf{Strategy:} \texttt{Comparator} passed to \texttt{Collections.sort()}.
\item
  \textbf{Observer:} \texttt{PropertyChangeListener}, deprecated \texttt{java.util.Observer}.
\item
  \textbf{Template Method:} \texttt{AbstractList}, \texttt{HttpServlet.service()}.
\item
  \textbf{Iterator:} the entire \texttt{Iterable}/\texttt{Iterator} framework behind for-each.
\item
  \textbf{Proxy:} \texttt{java.lang.reflect.Proxy}, RMI stubs.
\item
  \textbf{Flyweight:} \texttt{Integer.valueOf()} cache, the \texttt{String} pool.
\end{itemize}

\subsection{Spring Framework}\label{spring-framework}

\texttt{ApplicationContext} (Factory + IoC container), default singleton beans, \texttt{@Transactional}/AOP (Proxy), \texttt{JdbcTemplate}/\texttt{RestTemplate} (Template Method), \texttt{ApplicationEvent} (Observer), \texttt{@Configuration} (Factory + Builder-ish).

\subsection{React / Frontend}\label{react--frontend}

Components observing state and re-rendering (Observer), the component tree (Composite), HOCs and hooks (Decorator), render props / dependency-passing (Strategy), Context (Mediator-ish), Redux store (Observer + Command --- actions are commands, reducers interpret them).

\subsection{Backend / Distributed (the distributed cousins)}\label{backend--distributed-the-distributed-cousins}

\begin{itemize}
\tightlist
\item
  \textbf{Strategy} \(\rightarrow\) pluggable load-balancing / routing / retry policies.
\item
  \textbf{Observer} \(\rightarrow\) Kafka/RabbitMQ pub-sub, webhooks.
\item
  \textbf{Command} \(\rightarrow\) task queues (Celery, Sidekiq), CQRS commands, event sourcing.
\item
  \textbf{Chain of Responsibility} \(\rightarrow\) API gateway / middleware filters (Spring Security, Express).
\item
  \textbf{Proxy} \(\rightarrow\) service mesh sidecars (Envoy), API gateways, caching layers.
\item
  \textbf{Facade} \(\rightarrow\) BFF (backend-for-frontend), SDK clients.
\item
  \textbf{Circuit Breaker / Saga / CQRS} \(\rightarrow\) architectural-scale behavioral patterns (see §System Design).
\end{itemize}

\section{Real-Life Analogies}\label{real-life-analogies}

\emph{One restaurant --- every pattern is a role, a tool, or a workflow in the same kitchen.}

\begin{longtable}[]{@{}
  >{\raggedright\arraybackslash}p{(\columnwidth - 2\tabcolsep) * \real{0.2121}}
  >{\raggedright\arraybackslash}p{(\columnwidth - 2\tabcolsep) * \real{0.7879}}@{}}
\toprule\noalign{}
\rowcolor{acc!16}
\begin{minipage}[b]{\linewidth}\raggedright
Pattern
\end{minipage} & \begin{minipage}[b]{\linewidth}\raggedright
Analogy
\end{minipage} \\
\midrule\noalign{}
\endhead
\bottomrule\noalign{}
\endlastfoot
\textbf{Singleton} & The restaurant\textquotesingle s single head chef --- there is exactly one, everyone routes orders through them \\
\rowcolor{zebra}
\textbf{Factory Method} & "Make me the soup of the day" --- the kitchen decides which soup to cook based on the day \\
\textbf{Abstract Factory} & Ordering a "tasting menu" --- every course (appetizer, main, dessert) comes from one matching themed set \\
\rowcolor{zebra}
\textbf{Builder} & Building a custom burger at the counter --- bun, then patty, then toppings, each step optional, one "done" at the end \\
\textbf{Prototype} & The signature dish recipe card --- clone it and tweak the spice rather than inventing from scratch \\
\rowcolor{zebra}
\textbf{Adapter} & A wine-bottle opener that fits both cork and screw-cap --- bridges two incompatible bottle types \\
\textbf{Decorator} & Building a coffee order --- espresso, then add milk, then add caramel; each layer adds without remaking the cup \\
\rowcolor{zebra}
\textbf{Facade} & The waiter --- you say "I\textquotesingle ll have the special," they coordinate kitchen, bar, and till; you never enter the kitchen \\
\textbf{Proxy} & The host at the door --- screens reservations, checks the list, only seats you if you qualify \\
\rowcolor{zebra}
\textbf{Composite} & A banquet order --- "table of 8" is treated like one order, yet each guest\textquotesingle s plate is also handled individually \\
\textbf{Bridge} & Menu (dishes) \(\times\) cooking method (grill/fry/steam) --- add a dish or a method independently, no combinatorial menu \\
\rowcolor{zebra}
\textbf{Flyweight} & One shared salt shaker per table instead of a personal one per guest --- shared, identical, reused \\
\textbf{Strategy} & Choosing a payment method at the till --- cash, card, or app; same checkout, swappable method \\
\rowcolor{zebra}
\textbf{Observer} & The "order up!" bell --- when a dish is ready, every relevant waiter is notified at once \\
\textbf{Command} & A written order ticket --- queued on the rail, can be re-fired, and "void" undoes it \\
\rowcolor{zebra}
\textbf{State} & A table\textquotesingle s lifecycle --- Empty \(\rightarrow\) Seated \(\rightarrow\) Ordered \(\rightarrow\) Served \(\rightarrow\) Paying; each state allows different actions \\
\textbf{Template Method} & The standard service script --- greet, take order, serve, bill; the steps are fixed, the dishes vary \\
\rowcolor{zebra}
\textbf{Iterator} & The waiter going table to table in order, taking each order without a floor map \\
\textbf{Chain of Responsibility} & A complaint escalating --- waiter \(\rightarrow\) manager \(\rightarrow\) owner, until someone resolves it \\
\rowcolor{zebra}
\textbf{Mediator} & The expediter ("pass") who coordinates between cooks, waiters, and bar so they don\textquotesingle t all shout at each other \\
\textbf{Memento} & Saving a half-eaten meal as a "to-go box" to resume the exact state later \\
\rowcolor{zebra}
\textbf{Visitor} & The health inspector who walks the same kitchen applying a new checklist without changing the kitchen \\
\textbf{Interpreter} & Reading a recipe\textquotesingle s shorthand ("2T butter, fold, rest") --- each token interpreted into an action \\
\rowcolor{zebra}
\textbf{Dependency Injection} & The supplier delivers ingredients to the kitchen --- chefs don\textquotesingle t go forage; they\textquotesingle re handed what they need \\
\textbf{Object Pool} & A rack of clean pans reused across orders rather than buying a new pan per dish \\
\end{longtable}

\section{Memory Tricks / Mnemonics}\label{memory-tricks--mnemonics}

\subsection{The three families}\label{the-three-families}

\textbf{"Create, Structure, Behave"} --- Creational (how objects are \emph{created}), Structural (how they\textquotesingle re \emph{composed}), Behavioral (how they \emph{communicate}).

\subsection{Creational --- "SABFP"}\label{creational--sabfp}

\textbf{S}ingleton \(\cdot\) \textbf{A}bstract Factory \(\cdot\) \textbf{B}uilder \(\cdot\) \textbf{F}actory Method \(\cdot\) \textbf{P}rototype \(\rightarrow\) \emph{"So A Builder Factory Prototypes."}

\subsection{Structural --- "A Big Cat Drinks From Pool" (ABCDFP + Proxy)}\label{structural--a-big-cat-drinks-from-pool-abcdfp--proxy}

\textbf{A}dapter \(\cdot\) \textbf{B}ridge \(\cdot\) \textbf{C}omposite \(\cdot\) \textbf{D}ecorator \(\cdot\) \textbf{F}acade \(\cdot\) fly\textbf{P}roxy/Flyweight.

\subsection{Behavioral --- "Smart Os Cause State To Iterate Calmly, Mostly Visiting Interpreters"}\label{behavioral--smart-os-cause-state-to-iterate-calmly-mostly-visiting-interpreters}

\textbf{S}trategy \(\cdot\) \textbf{O}bserver \(\cdot\) \textbf{C}ommand \(\cdot\) \textbf{S}tate \(\cdot\) \textbf{T}emplate Method \(\cdot\) \textbf{I}terator \(\cdot\) \textbf{C}hain of Resp. \(\cdot\) \textbf{M}ediator \(\cdot\) \textbf{M}emento \(\cdot\) \textbf{V}isitor \(\cdot\) \textbf{I}nterpreter.

\subsection{Killer distinctions (the ones interviewers probe)}\label{killer-distinctions-the-ones-interviewers-probe}

\begin{itemize}
\tightlist
\item
  \textbf{Strategy vs. State:} \emph{Strategy} --- \textbf{you} pick the algorithm (it doesn\textquotesingle t change itself). \emph{State} --- the object \textbf{changes its own} behavior as it transitions.
\item
  \textbf{Decorator vs. Proxy:} same structure (both wrap). \emph{Decorator adds} responsibilities (stackable); \emph{Proxy controls} access (usually one).
\item
  \textbf{Adapter vs. Facade:} \emph{Adapter} makes one interface \emph{match} an expected one. \emph{Facade} makes many interfaces \emph{simpler}.
\item
  \textbf{Factory Method vs. Abstract Factory:} \emph{Method} makes \textbf{one} product (inheritance). \emph{Abstract} makes a \textbf{family} (composition).
\item
  \textbf{Decorator vs. Inheritance:} decorate at \textbf{runtime} (compose); inherit at \textbf{compile time} (fixed).
\item
  \textbf{Template Method vs. Strategy:} \emph{Template} varies \textbf{steps} via inheritance; \emph{Strategy} swaps the \textbf{whole algorithm} via composition.
\item
  \textbf{Mediator vs. Observer:} \emph{Observer} = broadcast (one\(\rightarrow\)many); \emph{Mediator} = coordination hub (many\(\leftrightarrow\)many).
\end{itemize}

\section{Common Interview Questions}\label{common-interview-questions}

\subsection{Q1: What is a design pattern, and when should you NOT use one?}\label{q1-what-is-a-design-pattern-and-when-should-you-not-use-one}

\textbf{Model answer:} A design pattern is a named, reusable \emph{template} for solving a recurring object-oriented design problem --- it captures an intent, a structure (the collaborating classes), and trade-offs. Patterns give a shared vocabulary and proven, SOLID-aligned structures that reduce coupling and ease change. You should NOT use one when it adds indirection without removing real pain --- a premature pattern violates YAGNI/KISS, adds classes, and makes code harder to read. Reach for a pattern only when you \emph{feel} the specific friction it relieves (a growing \texttt{switch} on type \(\rightarrow\) Strategy or State; subclass explosion \(\rightarrow\) Decorator). The senior skill is right-sizing: a plain function or polymorphism first, a named pattern only when justified.

\textbf{Follow-ups:}

\begin{itemize}
\tightlist
\item
  \emph{"Give an example of over-engineering with patterns."} \(\rightarrow\) Wrapping a single, never-varying implementation in a Strategy interface and Factory --- you pay the indirection cost for flexibility you\textquotesingle ll never use.
\item
  \emph{"How do you decide?"} \(\rightarrow\) Rule of Three: introduce the abstraction on the third time you actually need to vary something, not the first.
\end{itemize}

\subsection{Q2: Explain Strategy vs. State --- they look identical.}\label{q2-explain-strategy-vs-state--they-look-identical}

\textbf{Model answer:} Structurally they\textquotesingle re the same: a context object delegates to a swappable helper behind an interface. The \textbf{intent} differs. \textbf{Strategy} lets a \emph{client} choose among interchangeable algorithms that don\textquotesingle t know about each other and rarely change mid-flow --- it answers "\emph{how} do I do this?" (e.g., which sort, which payment method). \textbf{State} encapsulates behavior that depends on the object\textquotesingle s internal status, and crucially the states \emph{transition to one another} (a state often decides the next state) --- it answers "\emph{what can I do now?}" (e.g., an order in PENDING vs. SHIPPED). Strategy is set from outside; State drives itself through a lifecycle.

\textbf{Follow-ups:}

\begin{itemize}
\tightlist
\item
  \emph{"Which replaces a giant switch on a status field across many methods?"} \(\rightarrow\) State.
\item
  \emph{"Do strategies know each other?"} \(\rightarrow\) No. States do (they reference successor states).
\item
  \emph{"Can State be implemented with a transition table instead of classes?"} \(\rightarrow\) Yes --- a \texttt{Map\textless{}State,\ Map\textless{}Event,\ State\textgreater{}\textgreater{}}; cleaner when transitions dominate behavior.
\end{itemize}

\subsection{Q3: Why is Singleton often called an anti-pattern, and what\textquotesingle s the alternative?}\label{q3-why-is-singleton-often-called-an-anti-pattern-and-whats-the-alternative}

\textbf{Model answer:} Singleton introduces hidden \textbf{global mutable state}: a method calling \texttt{Config.get()} doesn\textquotesingle t declare that dependency in its signature, so coupling is invisible. It makes unit testing hard --- you can\textquotesingle t substitute a mock, and state leaks between tests. It adds thread-safety complexity and tends to proliferate ("everything\textquotesingle s a singleton"). The alternative is \textbf{Dependency Injection}: have a container create \emph{one} instance and inject it where needed. You keep the "exactly one instance" guarantee but make the dependency explicit, mockable, and swappable. Spring beans are exactly this --- singleton-scoped but DI-managed.

\textbf{Follow-ups:}

\begin{itemize}
\tightlist
\item
  \emph{"Show a thread-safe Singleton in Java."} \(\rightarrow\) Bill Pugh holder idiom (lazy + thread-safe + no sync cost) or an enum (also serialization/reflection-safe).
\item
  \emph{"Why is \texttt{volatile} required in double-checked locking?"} \(\rightarrow\) Without it, instruction reordering can publish a reference to a partially-constructed object to another thread.
\end{itemize}

\subsection{Q4: Decorator vs. Proxy vs. Inheritance --- distinguish them.}\label{q4-decorator-vs-proxy-vs-inheritance--distinguish-them}

\textbf{Model answer:} \textbf{Decorator} and \textbf{Proxy} both wrap an object behind the same interface. \textbf{Decorator} \emph{adds responsibilities} and you can stack many in any order (Java I/O streams: buffering + compression + encryption). \textbf{Proxy} \emph{controls access} to exactly one object --- lazy creation (virtual), authorization (protection), remote calls (RPC stub), or caching --- without changing what the object does. \textbf{Inheritance} adds behavior at compile time, can\textquotesingle t be combined or removed at runtime, and leads to subclass explosion when features combine; Decorator composes instead. Rule of thumb: adding features \(\rightarrow\) Decorator; gatekeeping access \(\rightarrow\) Proxy; fixed specialization with a true IS-A \(\rightarrow\) inheritance (sparingly).

\textbf{Follow-ups:}

\begin{itemize}
\tightlist
\item
  \emph{"Give the canonical Decorator in the JDK."} \(\rightarrow\) \texttt{new\ BufferedReader(new\ InputStreamReader(new\ FileInputStream(f)))}.
\item
  \emph{"Which does Spring \texttt{@Transactional} use?"} \(\rightarrow\) Proxy (a dynamic proxy opens/commits the transaction around the bean method).
\end{itemize}

\subsection{Q5: Factory Method vs. Abstract Factory?}\label{q5-factory-method-vs-abstract-factory}

\textbf{Model answer:} \textbf{Factory Method} is a single (often abstract) method that creates \emph{one} product, where a \emph{subclass} decides the concrete type --- it uses inheritance. \textbf{Abstract Factory} is an \emph{object} exposing several factory methods to create a \emph{family} of related products that must be used together (a matching \texttt{Button} + \texttt{Checkbox} per OS), guaranteeing compatibility --- it uses composition and is often \emph{implemented with} Factory Methods. Trade-off note: Abstract Factory is flexible against new \emph{families} but rigid against new \emph{product kinds} (adding a \texttt{Slider} edits every factory).

\textbf{Follow-ups:}

\begin{itemize}
\tightlist
\item
  \emph{"When do you prefer a simple static factory over either?"} \(\rightarrow\) Small cases with a finite, stable set of types --- one \texttt{create(type)} method is clearer than a class hierarchy.
\end{itemize}

\subsection{Q6: How does Observer avoid tight coupling, and what\textquotesingle s the classic bug?}\label{q6-how-does-observer-avoid-tight-coupling-and-whats-the-classic-bug}

\textbf{Model answer:} The subject knows only the \texttt{Observer} interface, not concrete observers; observers register/unregister at runtime. The subject broadcasts state changes without knowing who listens or what they do, so producers and consumers evolve independently --- this is the basis of event systems and reactive UIs. Push vs. pull: push sends the changed data in the notification; pull sends "something changed" and lets the observer query the subject. The classic bug is the \textbf{lapsed-listener leak}: the subject holds strong references to observers, so an observer that\textquotesingle s discarded without calling \texttt{unsubscribe()} is kept alive forever (memory leak) and keeps receiving useless notifications. Fixes: explicit unsubscribe, weak references, or lifecycle-scoped subscriptions.

\textbf{Follow-ups:}

\begin{itemize}
\tightlist
\item
  \emph{"What about notification order and storms?"} \(\rightarrow\) Order is undefined --- don\textquotesingle t depend on it; cascading updates can storm, so batch/debounce.
\item
  \emph{"Thread safety?"} \(\rightarrow\) Modifying the observer list while iterating it is a race; use \texttt{CopyOnWriteArrayList} or snapshot before notifying.
\end{itemize}

\subsection{Q7: Give a pattern that embodies the Open/Closed Principle and explain how.}\label{q7-give-a-pattern-that-embodies-the-openclosed-principle-and-explain-how}

\textbf{Model answer:} \textbf{Strategy.} Adding a new algorithm means writing a new class implementing the strategy interface; the \texttt{Context} is \emph{closed for modification} (untouched) yet the system is \emph{open for extension}. It also embodies DIP (the context depends on the abstraction, not concrete strategies) and composition over inheritance. Factory Method and Decorator are equally valid answers --- Factory lets you add product types without editing the creator\textquotesingle s logic; Decorator adds responsibilities without modifying the component.

\textbf{Follow-ups:}

\begin{itemize}
\tightlist
\item
  \emph{"What\textquotesingle s the cost of OCP here?"} \(\rightarrow\) More classes and indirection; only worth it if the variation point will actually be varied (Rule of Three).
\end{itemize}

\subsection{Q8: What pattern fits a request/middleware pipeline?}\label{q8-what-pattern-fits-a-requestmiddleware-pipeline}

\textbf{Model answer:} \textbf{Chain of Responsibility.} Each handler (auth \(\rightarrow\) rate-limit \(\rightarrow\) validation \(\rightarrow\) logging) gets a chance to process the request and decides whether to handle it and/or pass it to the next handler. It decouples the sender from the set of handlers, lets you reorder/insert/remove stages without touching others (OCP/SRP), and short-circuit (reject) early. The risk is a request falling off the end unhandled, so guard for that. This is exactly how Servlet filters, Spring Security, and Express/Django middleware work.

\textbf{Follow-ups:}

\begin{itemize}
\tightlist
\item
  \emph{"How does it differ from Decorator?"} \(\rightarrow\) Decorator always delegates and \emph{augments}; a CoR handler may \emph{stop} the chain (handle-or-forward), and the focus is "who handles this," not "add behavior."
\end{itemize}

\subsection{Q9: How would you implement undo/redo, and which patterns?}\label{q9-how-would-you-implement-undoredo-and-which-patterns}

\textbf{Model answer:} \textbf{Command} + \textbf{Memento}. Each user action becomes a \texttt{Command} object with \texttt{execute()} and \texttt{undo()}. An invoker maintains a history stack of executed commands and a redo stack. For complex state where computing an inverse is hard, the command captures a \textbf{Memento} (an opaque snapshot of the relevant state) before executing and restores it on undo --- this preserves encapsulation since only the originator can read the memento. Redo re-executes the popped command. Watch memory for long histories (snapshots are costly --- prefer command inverses or deltas when possible).

\textbf{Follow-ups:}

\begin{itemize}
\tightlist
\item
  \emph{"Command vs. Memento for undo?"} \(\rightarrow\) Command stores the \emph{operation} (and how to reverse it); Memento stores the \emph{state}. Use Command for invertible ops, Memento when reversal is impractical.
\item
  \emph{"How do macros fit?"} \(\rightarrow\) A composite command (list of commands) --- itself an instance of Composite + Command.
\end{itemize}

\subsection{Q10: When would you choose Composition over Inheritance, and which patterns reflect that?}\label{q10-when-would-you-choose-composition-over-inheritance-and-which-patterns-reflect-that}

\textbf{Model answer:} Prefer composition unless there\textquotesingle s a genuine, stable IS-A relationship and you want the full parent API. Inheritance creates tight coupling (fragile base class --- a parent change breaks children), fixes behavior at compile time, and explodes combinatorially when features combine. Composition delegates to a held object through its interface, is swappable at runtime, and keeps hierarchies flat. Most patterns are composition mechanisms: \textbf{Strategy} (hold an algorithm), \textbf{Decorator} (wrap a component), \textbf{Bridge} (abstraction holds an implementor), \textbf{State} (hold the current state), \textbf{Adapter/Proxy} (wrap an object). Inheritance shows up mainly in \textbf{Template Method} (override steps) and \textbf{Factory Method} (subclass picks the type).

\textbf{Follow-ups:}

\begin{itemize}
\tightlist
\item
  \emph{"When is inheritance the right call?"} \(\rightarrow\) True IS-A, shallow hierarchy (\(\leq\)2 levels), you want polymorphic dispatch and the full base contract (e.g., \texttt{ArrayList\ extends\ AbstractList}).
\end{itemize}

\subsection{Q11: How do real frameworks combine multiple patterns? Give a concrete example.}\label{q11-how-do-real-frameworks-combine-multiple-patterns-give-a-concrete-example}

\textbf{Model answer:} \textbf{MVC = Observer + Strategy + Composite.} The View observes the Model and re-renders on change (Observer); the Controller plugs in input-handling behavior (Strategy); a View built from nested widgets is a tree (Composite). Another: \textbf{Spring = Factory + Singleton + Proxy + Template Method} --- the \texttt{ApplicationContext} is a bean Factory, beans are singleton-scoped (DI-managed Singleton), \texttt{@Transactional} wraps beans in dynamic Proxies, and \texttt{JdbcTemplate} is a Template Method. Recognizing these combinations shows you understand patterns as a collaborating system, not isolated flashcards.

\textbf{Follow-ups:}

\begin{itemize}
\tightlist
\item
  \emph{"How does Java I/O combine patterns?"} \(\rightarrow\) Decorator (stream wrapping) + Adapter (\texttt{InputStreamReader} bytes\(\rightarrow\)chars) + Iterator (collection traversal).
\end{itemize}

\subsection{Q12: What\textquotesingle s an anti-pattern, and how do you fix the God Object?}\label{q12-whats-an-anti-pattern-and-how-do-you-fix-the-god-object}

\textbf{Model answer:} An anti-pattern is a common but counterproductive "solution" that creates more problems than it solves. The \textbf{God Object} is a class that knows and does everything (an \texttt{OrderManager} handling pricing, payment, email, persistence) --- it violates SRP, is untestable, and every change risks breaking unrelated features. The fix: extract each responsibility into a focused collaborator (SRP), push behavior into the right domain objects (kill the anemic model), and if you still want a single easy entry point, put a thin \textbf{Facade} in front of the collaborators --- the Facade simplifies access \emph{without} absorbing the logic. Other anti-patterns to name: Golden Hammer (one pattern everywhere \(\rightarrow\) match tool to problem), Premature Abstraction (Rule of Three), Singleton Abuse (use DI), Spaghetti (layer + polymorphism).

\textbf{Follow-ups:}

\begin{itemize}
\tightlist
\item
  \emph{"How is a Facade different from a God Object?"} \(\rightarrow\) A Facade \emph{delegates} to focused subsystem classes (thin coordinator); a God Object \emph{contains} all the logic itself.
\end{itemize}

\section{Senior-Level Discussion Points}\label{senior-level-discussion-points}

\begin{itemize}
\item
  \textbf{Patterns are trade-offs, not virtues.} Every pattern adds indirection and classes; you must justify the cost against the pain it relieves. The senior move is often \emph{not} using a pattern --- recognizing when a plain function, polymorphism, or a first-class function is enough.
\item
  \textbf{Language shapes patterns.} Many GoF patterns are workarounds for missing language features. First-class functions collapse Strategy/Command/Observer into "pass a function." Python\textquotesingle s duck typing makes Adapter nearly free. Go favors composition + small interfaces and has no inheritance. In functional languages, immutability + higher-order functions replace much of the behavioral catalog. Naming this nuance signals depth.
\item
  \textbf{Encapsulate-what-varies is the meta-principle.} Almost every pattern is the same move: find the axis of change, hide it behind an interface, depend on the interface (OCP + DIP). If you internalize this, you can \emph{derive} the right pattern rather than recall it.
\item
  \textbf{Patterns combine and layer.} Real systems are compositions (MVC, Spring, React, compilers). Discuss how they collaborate, not just individual definitions.
\item
  \textbf{Thread-safety in patterns is a frequent follow-up.} Singleton needs safe lazy init (\texttt{volatile} DCL, holder idiom, enum). Observer needs a thread-safe observer list (\texttt{CopyOnWriteArrayList}) and care iterating while modifying. Flyweight\textquotesingle s shared intrinsic state \emph{must} be immutable. Object Pool needs synchronized borrow/return.
\item
  \textbf{Anti-patterns are half the skill.} Recognizing God objects, golden hammers, premature abstraction, poltergeists, and singleton abuse --- and naming the refactor --- is as valuable as knowing the 23 patterns. Interviewers love "what\textquotesingle s wrong with this design?"
\item
  \textbf{Architectural patterns are the distributed cousins.} CQRS, Saga, Circuit Breaker, Event Sourcing, Sidecar/Ambassador (a distributed Proxy) --- these scale GoF intents to services. Mentioning the bridge from LLD patterns to HLD patterns shows range.
\item
  \textbf{The Rule of Three governs abstraction.} Don\textquotesingle t abstract on the first occurrence (premature), or even the second. On the third real need to vary, introduce the abstraction. Speculative generality is a tax you pay forever for flexibility you may never use.
\end{itemize}

\section{Typical Mistakes Candidates Make}\label{typical-mistakes-candidates-make}

\begin{longtable}[]{@{}
  >{\raggedright\arraybackslash}p{(\columnwidth - 4\tabcolsep) * \real{0.2948}}
  >{\raggedright\arraybackslash}p{(\columnwidth - 4\tabcolsep) * \real{0.3269}}
  >{\raggedright\arraybackslash}p{(\columnwidth - 4\tabcolsep) * \real{0.3783}}@{}}
\toprule\noalign{}
\rowcolor{acc!16}
\begin{minipage}[b]{\linewidth}\raggedright
Mistake
\end{minipage} & \begin{minipage}[b]{\linewidth}\raggedright
Why it\textquotesingle s wrong
\end{minipage} & \begin{minipage}[b]{\linewidth}\raggedright
What to do instead
\end{minipage} \\
\midrule\noalign{}
\endhead
\bottomrule\noalign{}
\endlastfoot
\textbf{Pattern soup} --- forcing patterns everywhere & Over-engineering; indirection with no payoff & Smallest solution first; escalate only on real friction \\
\rowcolor{zebra}
\textbf{Naming the structure, not the intent} & Reciting UML without \emph{why} sounds cargo-cult & Explain the force it relieves and the trade-off \\
\textbf{Singleton everywhere} & Global mutable state, untestable & Dependency Injection of a single managed instance \\
\rowcolor{zebra}
\textbf{Confusing Strategy/State or Decorator/Proxy} & Shows shallow understanding & Memorize the \emph{intent} distinction, not just the shape \\
\textbf{Inheritance-based "decoration"} & Subclass explosion & Compose decorators at runtime \\
\rowcolor{zebra}
\textbf{Ignoring thread-safety} & Broken Singleton/Observer/Flyweight under concurrency & Proactively address it (volatile, immutability, \texttt{CopyOnWriteArrayList}) \\
\textbf{Forgetting the lapsed-listener leak} & Observer memory leaks in long-running systems & Always provide/explain unsubscribe or weak refs \\
\rowcolor{zebra}
\textbf{Premature abstraction} & Interfaces for variation that never comes & Rule of Three --- start concrete \\
\textbf{Shallow vs. deep copy in Prototype} & Cloned objects share mutable state & Choose deep copy deliberately; document it \\
\rowcolor{zebra}
\textbf{Anemic model + God service} & Logic dumped in services, not objects & Put behavior on the object that owns the data \\
\textbf{Adding patterns without trade-off talk} & Misses the senior signal & Always state the cost ("more classes, but OCP") \\
\end{longtable}

\section{How This Connects to Other Topics}\label{how-this-connects-to-other-topics}

\begin{itemize}
\item
  \textbf{OOP (\hyperlink{section.12}{\textcolor{acc}{§12}}) \& LLD (\hyperlink{section.06}{\textcolor{acc}{§06}})} --- Patterns \emph{are} applied OOP; LLD/machine-coding rounds expect you to reach for them appropriately. Patterns are the vocabulary; LLD is where you speak it.
\item
  \textbf{System Design / HLD (\hyperlink{section.05}{\textcolor{acc}{§05}})} --- Architectural patterns (CQRS, Saga, Circuit Breaker, Pub/Sub, Sidecar, API Gateway) are the \emph{distributed} descendants of GoF intents: Pub/Sub is Observer at scale, a Sidecar is a Proxy, a Gateway is a Facade + Chain of Responsibility.
\item
  \textbf{Concurrency (\hyperlink{section.01}{\textcolor{acc}{§01}})} --- Singleton/Flyweight/shared state demand thread-safety; producer-consumer is structurally Command + a queue; Object Pool underlies thread pools and connection pools.
\item
  \textbf{Databases / ORM (\hyperlink{section.08}{\textcolor{acc}{§08}})} --- Repository abstracts the data store; Active Record (model persists itself) is the alternative trade-off; Unit of Work groups operations into a transaction; connection pooling is Object Pool.
\item
  \textbf{Testing} --- Good pattern use enables testing: DI injects mocks, Strategy swaps test algorithms, Repository\textquotesingle s in-memory impl avoids a real DB, Null Object simplifies test setup. Patterns and testability rise together.
\item
  \textbf{Frameworks you use daily} --- Spring, React, the JDK, Django, Hibernet are \emph{built from} these patterns; understanding them lets you read framework source and debug "magic."
\end{itemize}

\section{FAANG Interview Tips}\label{faang-interview-tips}

\begin{enumerate}
\def\labelenumi{\arabic{enumi}.}
\item
  \textbf{Let the requirement reveal the pattern.} Say the trigger out loud: "the algorithm must be swappable" \(\rightarrow\) Strategy; "notify many on change" \(\rightarrow\) Observer; "undo" \(\rightarrow\) Command + Memento; "behavior depends on a status that transitions" \(\rightarrow\) State. Naming it as you go signals pattern literacy.
\item
  \textbf{Prefer the smallest solution.} A function, polymorphism, or a first-class function before a named pattern. Escalate only when the friction is real. Over-patterning is penalized as hard as under-designing.
\item
  \textbf{Always state the trade-off.} "Decorator avoids subclass explosion but produces many small wrapper objects and ordering can matter." The trade-off sentence is the senior signal.
\item
  \textbf{Know the 5 cold:} Strategy, Observer, Factory, Decorator, Singleton --- they cover the majority of LLD prompts. Then State, Command, Adapter, Composite, Builder for the next tier.
\item
  \textbf{Nail the killer distinctions:} Strategy vs. State, Decorator vs. Proxy, Adapter vs. Facade, Factory Method vs. Abstract Factory. Interviewers probe these to separate memorizers from understanders.
\item
  \textbf{Show the extensibility story.} After designing, say: "to add an EV charging spot, I subclass / add a new strategy --- zero edits to existing code (OCP)." This is the \#1 senior moment.
\item
  \textbf{Address concurrency proactively} when the system is multi-threaded --- Singleton (\texttt{volatile} DCL / holder idiom / enum), Observer (\texttt{CopyOnWriteArrayList}, lapsed-listener leak), Flyweight (immutability).
\item
  \textbf{Recognize anti-patterns out loud} if you spot a God object / golden hammer / premature abstraction in a given design --- and name the refactor. "What\textquotesingle s wrong with this design?" is a common prompt.
\item
  \textbf{Demonstrate combination knowledge.} "MVC is Observer + Strategy + Composite; Spring is Factory + Singleton + Proxy." It shows systemic understanding.
\item
  \textbf{Amazon-specific:} map choices to Leadership Principles --- "Dive Deep" (knowing pattern internals like why \texttt{volatile} is needed in DCL), "Invent and Simplify" (choosing the \emph{right} pattern, not the most complex), "Insist on Highest Standards" (clean, SOLID code).
\end{enumerate}

\section{Revision Cheat Sheet}\label{revision-cheat-sheet}

\subsection{Families}\label{families}

\textbf{Creational} (how objects are \emph{created}) \(\cdot\) \textbf{Structural} (how they\textquotesingle re \emph{composed}) \(\cdot\) \textbf{Behavioral} (how they \emph{communicate}).

\subsection{One-liners}\label{one-liners}

\begin{longtable}[]{@{}
  >{\raggedright\arraybackslash}p{(\columnwidth - 4\tabcolsep) * \real{0.2932}}
  >{\raggedright\arraybackslash}p{(\columnwidth - 4\tabcolsep) * \real{0.0880}}
  >{\raggedright\arraybackslash}p{(\columnwidth - 4\tabcolsep) * \real{0.6189}}@{}}
\toprule\noalign{}
\rowcolor{acc!16}
\begin{minipage}[b]{\linewidth}\raggedright
Pattern
\end{minipage} & \begin{minipage}[b]{\linewidth}\raggedright
Family
\end{minipage} & \begin{minipage}[b]{\linewidth}\raggedright
One-liner
\end{minipage} \\
\midrule\noalign{}
\endhead
\bottomrule\noalign{}
\endlastfoot
\textbf{Singleton} & C & One shared instance (use sparingly; prefer DI) \\
\rowcolor{zebra}
\textbf{Factory Method} & C & Subclass decides the concrete type (inheritance) \\
\textbf{Abstract Factory} & C & Create compatible product \emph{families} (composition) \\
\rowcolor{zebra}
\textbf{Builder} & C & Step-by-step / immutable construction; no telescoping ctors \\
\textbf{Prototype} & C & Clone an exemplar instead of \texttt{new} \\
\rowcolor{zebra}
\textbf{Adapter} & S & Convert one interface to another (compatibility) \\
\textbf{Decorator} & S & Stack responsibilities at runtime (Java I/O) \\
\rowcolor{zebra}
\textbf{Facade} & S & Simple front for a complex subsystem \\
\textbf{Proxy} & S & Controlled stand-in (virtual/protection/remote/cache) \\
\rowcolor{zebra}
\textbf{Composite} & S & Treat leaf \& group uniformly (trees) \\
\textbf{Bridge} & S & Split abstraction from impl (avoid n\(\times\)m) \\
\rowcolor{zebra}
\textbf{Flyweight} & S & Share intrinsic state to save memory \\
\textbf{Strategy} & B & Swap interchangeable algorithms (OCP + DIP) \\
\rowcolor{zebra}
\textbf{Observer} & B & Pub/sub; notify many on change (mind the leak) \\
\textbf{Command} & B & Request as object (undo/queue/log) \\
\rowcolor{zebra}
\textbf{State} & B & Behavior changes as state transitions (replaces switch) \\
\textbf{Template Method} & B & Fixed skeleton, overridable steps (Hollywood) \\
\rowcolor{zebra}
\textbf{Iterator} & B & Sequential access, hide structure \\
\textbf{Chain of Responsibility} & B & Pipeline of handlers (middleware) \\
\rowcolor{zebra}
\textbf{Mediator} & B & Central hub for many-to-many comms \\
\textbf{Memento} & B & Snapshot/restore state (undo, encapsulated) \\
\rowcolor{zebra}
\textbf{Visitor} & B & Add operations to a stable hierarchy \\
\textbf{Interpreter} & B & Evaluate sentences of a small grammar \\
\rowcolor{zebra}
\textbf{Dependency Injection} & Modern & Inject abstractions; testable, swappable \\
\textbf{MVC/MVP/MVVM} & Modern & Separate data / UI / glue \\
\rowcolor{zebra}
\textbf{Repository} & Modern & Collection-like abstraction over persistence \\
\textbf{Null Object} & Modern & Neutral do-nothing default (kills null checks) \\
\rowcolor{zebra}
\textbf{Object Pool} & Modern & Borrow/return expensive reusable objects \\
\end{longtable}

\subsection{Killer distinctions (memorize)}\label{killer-distinctions-memorize}

\begin{itemize}
\tightlist
\item
  \textbf{Strategy vs. State:} you pick the algorithm \textbar{} object changes itself.
\item
  \textbf{Decorator vs. Proxy:} adds behavior (stack) \textbar{} controls access (one).
\item
  \textbf{Adapter vs. Facade:} make it \emph{match} \textbar{} make it \emph{simpler}.
\item
  \textbf{Factory Method vs. Abstract Factory:} one product (inherit) \textbar{} family (compose).
\item
  \textbf{Template Method vs. Strategy:} vary steps (inherit) \textbar{} swap whole (compose).
\item
  \textbf{Mediator vs. Observer:} coordinate many\(\leftrightarrow\)many \textbar{} broadcast one\(\rightarrow\)many.
\end{itemize}

\subsection{Singleton thread-safe forms (Java)}\label{singleton-thread-safe-forms-java}

Eager (simple) \(\cdot\) DCL with \texttt{volatile} (lazy) \(\cdot\) \textbf{Bill Pugh holder} (best general) \(\cdot\) \textbf{enum} (Bloch\textquotesingle s choice --- serialization/reflection-safe).

\subsection{\texorpdfstring{Anti-patterns \(\rightarrow\) fixes}{Anti-patterns \textbackslash rightarrow fixes}}\label{anti-patterns-rightarrow-fixes}

God Object \(\rightarrow\) SRP/extract \(\cdot\) Golden Hammer \(\rightarrow\) match tool \(\cdot\) Spaghetti \(\rightarrow\) layer + polymorphism \(\cdot\) Poltergeist \(\rightarrow\) delete \(\cdot\) Premature Abstraction \(\rightarrow\) Rule of Three \(\cdot\) Singleton Abuse \(\rightarrow\) DI.

\subsection{Golden rules}\label{golden-rules}

\begin{enumerate}
\def\labelenumi{\arabic{enumi}.}
\tightlist
\item
  \textbf{Encapsulate what varies, program to interfaces, favor composition.}
\item
  \textbf{Don\textquotesingle t pattern what a function can do} (KISS/YAGNI).
\item
  \textbf{Rule of Three} --- abstract on the third real need, not the first.
\item
  \textbf{Name the pattern \emph{and} the trade-off} --- that\textquotesingle s the senior signal.
\end{enumerate}

\emph{Study tip: For each of the five core patterns (Strategy, Observer, Factory, Decorator, Singleton), be able to (1) state the intent in one line, (2) write the code from memory, (3) name a JDK/Spring/React example, and (4) recite the trade-off. Then drill the killer distinctions until they\textquotesingle re reflexive --- that\textquotesingle s what separates "I memorized the catalog" from "I understand design."}

\end{document}
